% Lesson 1 — Speaking: Exchanges information about interacting with mates, neighbours, adults and school authorities on community celebrations/festivals — Anglais 1ère A — Cours signé OnBuch+
% Programme officiel MINESEC. Compiler avec : tectonic ENA01-speaking-exchanges-information-about-interacting-with-m.tex
\newcommand{\DOCMATIERE}{Anglais}
\newcommand{\DOCNIVEAU}{1ère A}
\newcommand{\DOCMODULE}{Chapter 1 : Family and Social Life}
\newcommand{\DOCLECON}{ENA01}
\newcommand{\DOCTITRE}{Lesson 1 — Speaking: Exchanges information about interacting with mates, neighbours, adults and school authorities on community celebrations/festivals}
% OnBuch+ — préambule « cours signés » ENRICHI (compilé avec tectonic / XeLaTeX)
% Système visuel « Pop » d'OnBuch+ : fond crème (jamais blanc), bordures encre
% épaisses, ombres « dures » décalées, titres Archivo Black en capitales.
% Version MATHÉMATIQUES : graphes (pgfplots/tikz), boîtes pédagogiques
% (définition, propriété, méthode, attention, le savais-tu, exercice résolu,
% exercice, TP, bilan), page de garde et signature OnBuch+ ancrée en bas.
%
% Macros attendues dans main.tex AVANT \input{preamble} :
%   \DOCMATIERE  \DOCNIVEAU  \DOCMODULE  \DOCLECON (numéro)  \DOCTITRE
\documentclass[11pt]{article}

\usepackage{fontspec}
\usepackage[english,french]{babel} % français = langue principale ; l'anglais via \eng{..} / textebox
\usepackage[a4paper,margin=2cm,top=3.4cm,bottom=2.8cm,headheight=1.4cm,footskip=1.3cm]{geometry}
\usepackage{amsmath,amssymb,amsfonts}
\usepackage{mathtools} % \xmapsto, \coloneqq... souvent utilisés par les modèles
\usepackage{cancel} % simplifications barrées (bilans de forces)
% \abs{z} (module, valeur absolue) et \norm{u} (norme), utilisés par les modèles
\providecommand{\abs}{}\let\abs\relax\DeclarePairedDelimiter{\abs}{\lvert}{\rvert}
\providecommand{\norm}{}\let\norm\relax\DeclarePairedDelimiter{\norm}{\lVert}{\rVert}
\usepackage[table]{xcolor} % [table] : fournit aussi \rowcolors (alternance de lignes)
\usepackage{pagecolor}
\usepackage{tikz}
\usetikzlibrary{arrows.meta,positioning,calc,shapes.geometric,shapes.misc,shapes.arrows,decorations.pathmorphing,decorations.pathreplacing,decorations.markings,patterns,shadows,mindmap,trees,fit,backgrounds,matrix,intersections,angles,quotes}
\usepackage{pgfplots}
\usepackage[version=4]{mhchem} % équations nucléaires \ce{^{235}_{92}U}
\usepackage[european,siunitx]{circuitikz} % schémas électriques
\pgfplotsset{compat=1.18}
\usepgfplotslibrary{fillbetween,groupplots} % aires entre courbes (calcul intégral)
\usepackage{enumitem}
\usepackage{fancyhdr}
% [most] charge toutes les bibliothèques tcolorbox (listings, minted, raster,
% documentation...) et gonfle inutilement la mémoire TeX sur un document long
% (~300 boîtes) : on ne charge que ce qui est réellement utilisé.
\usepackage[skins,breakable]{tcolorbox}
\usepackage{graphicx}
\usepackage{colortbl}
\usepackage{booktabs}
\usepackage{tabularx}
\usepackage{diagbox} % case d'en-tête diagonale des tableaux à double entrée
% Colonnes extensibles centrée / gauche / droite (C, L, R), souvent utilisées par les modèles
\newcolumntype{C}{>{\centering\arraybackslash}X}
\newcolumntype{L}{>{\raggedright\arraybackslash}X}
\newcolumntype{R}{>{\raggedleft\arraybackslash}X}
\usepackage{array}
\usepackage{multicol}
\usepackage{multirow}
\usepackage{siunitx}
% parse-numbers=false : accepte aussi \SI{2,5 \times 10^{-3}}{...}
\sisetup{locale=FR,output-decimal-marker={,},group-minimum-digits=4,parse-numbers=false}
\DeclareSIUnit\ohm{\ensuremath{\symup{\Omega}}} % U+2126 absent de la police
\DeclareSIUnit\division{div} % réglages d'oscilloscope (ms/div, V/div)
\DeclareSIUnit\tr{tr} % tours (vitesse de rotation, ex. \SI{1500}{\tr\per\minute})
\DeclareSIUnit\molar{M} % concentration molaire (ex. \SI{3}{\molar}, \SI{1}{\micro\molar})
\DeclareSIUnit\an{an} % année (le modèle confond parfois avec \year, un compteur TeX) — ex. \SI{5}{\kilo\meter\per\an}
\DeclareSIUnit\FCFA{FCFA} % franc CFA (ex. \SI{500}{\FCFA\per\kilo\gram})
\DeclareSIUnit\degreeBrix{\textdegree Brix} % teneur en sucre d'un sirop/jus (réfractométrie)
\DeclareSIUnit\month{mois} % le modèle utilise parfois \month, un compteur TeX, comme unité de durée
\usepackage{qrcode}
% \degree : commande fréquemment utilisée par les modèles pour un angle ou une
% température en dehors de \SI{}{} (non fournie par les paquets chargés).
\providecommand{\degree}{^{\circ}}
% \qed (fin de démonstration), utilisé par les modèles sans amsthm
\providecommand{\qed}{\hfill$\square$}
% \barre : double barre des valeurs interdites dans un tableau de variations (tkz-tab)
\providecommand{\barre}{\ensuremath{\|}}
% environnement proof (amsthm non chargé) utilisé par les modèles
\ifdefined\proof\else\newenvironment{proof}[1][Démonstration]{\par\noindent\textit{#1.}\ }{\qed\par}\fi
\usepackage{lastpage}
\usepackage{hyperref}
\hypersetup{hidelinks}

% ============================================================
% Palette « Pop » OnBuch+ — mêmes hex que la classe P (plus_ui.dart)
% ============================================================
\definecolor{poporange}{HTML}{F59321}
\definecolor{popdark}{HTML}{E07A0C}
\definecolor{popcream}{HTML}{FFF6E8}
\definecolor{popink}{HTML}{1C1714}
\definecolor{poppurple}{HTML}{7A5AE0}
\definecolor{popgreen}{HTML}{2E9E6A}
\definecolor{poppink}{HTML}{E0457B}
\definecolor{popblue}{HTML}{2F7DE1}
\definecolor{popgold}{HTML}{C98A12}
\definecolor{popmuted}{HTML}{8A7F76}
\definecolor{popline}{HTML}{EADFD2}
\definecolor{popgray}{HTML}{8A7F76}
\colorlet{popgrey}{popgray}
% Alias tolérants (le modèle nomme parfois les couleurs différemment)
\colorlet{popwhite}{white}
\colorlet{popblack}{popink}
\colorlet{popred}{poppink}
\colorlet{popyellow}{popgold}
% Teintes claires (fonds de tableaux/encadrés) — le modèle nomme parfois
% les couleurs avec un suffixe L (ex. popblueL) sans les avoir définies.
\colorlet{poporangeL}{poporange!20}
\colorlet{popdarkL}{popdark!20}
\colorlet{poppurpleL}{poppurple!20}
\colorlet{popgreenL}{popgreen!20}
\colorlet{poppinkL}{poppink!20}
\colorlet{popblueL}{popblue!20}
\colorlet{popgoldL}{popgold!20}
\colorlet{popredL}{popred!20}
\colorlet{popgrayL}{popgray!20}
% Teintes claires pour fonds de tableaux / zones
\colorlet{popblueL}{popblue!10!white}
\colorlet{poporangeL}{poporange!14!white}
\colorlet{popgreenL}{popgreen!12!white}
\colorlet{poppurpleL}{poppurple!12!white}
\colorlet{poppinkL}{poppink!12!white}
\colorlet{popgoldL}{popgold!14!white}

\pagecolor{popcream}

% ============================================================
% Typographie
% ============================================================
\defaultfontfeatures{Ligatures=TeX}
\setmainfont{PlusJakartaSans-Medium.ttf}[Path=../../../fonts/,BoldFont=PlusJakartaSans-Bold.ttf]
% Symboles, API et emojis absents de la police principale : repli sur DejaVu Sans (caractères actifs)
\newfontface\symfont{DejaVuSans.ttf}[Path=../../../fonts/]
\makeatletter
\newcommand\symactive[1]{\catcode#1=13 \begingroup\lccode`\~=#1 \lowercase{\endgroup\def~}{{\symfont\char#1}}}
\newcommand\symblank[1]{\catcode#1=13 \begingroup\lccode`\~=#1 \lowercase{\endgroup\def~}{}}
\newcommand\symhyph[1]{\catcode#1=13 \begingroup\lccode`\~=#1 \lowercase{\endgroup\def~}{\mbox{-}}}
\newcount\symc
\newcommand\symrange[2]{\symc=#1 \@whilenum\symc<#2 \do{\expandafter\symactive\expandafter{\the\symc}\advance\symc by 1 }}
\newcommand\symblankrange[2]{\symc=#1 \@whilenum\symc<#2 \do{\expandafter\symblank\expandafter{\the\symc}\advance\symc by 1 }}
\makeatother
\symrange{"0250}{"02FF}\symactive{"2713}\symactive{"2714}\symactive{"2717}\symactive{"2718}\symactive{"2610}\symactive{"2611}\symactive{"2612}
\symhyph{"2011}\symhyph{"2E1A}\symblankrange{"1F300}{"1FAFF}\symblank{"FE0F}
% Maths en sans-serif (Fira Math) avec chiffres et lettres droites de Plus
% Jakarta, pour que les formules se fondent dans le texte.
\usepackage[math-style=ISO,bold-style=ISO]{unicode-math}
\setmathfont{FiraMath-Regular.otf}[Path=../../../fonts/]
\setmathfont{PlusJakartaSans-Medium.ttf}[Path=../../../fonts/,range={up/num,up/{Latin,latin},\mathup}]
\usepackage{icomma} % virgule décimale sans espace en mode math
% Caractères mathématiques Unicode absents des polices de texte (Plus Jakarta,
% Archivo Black) : rendus par leur équivalent mathématique, en texte comme en
% formule — sinon ils s'affichent en carré vide (ex. « Arithmétique dans ℤ »).
\usepackage{newunicodechar}
\newunicodechar{ℝ}{\ensuremath{\mathbb{R}}}
\newunicodechar{ℤ}{\ensuremath{\mathbb{Z}}}
\newunicodechar{ℂ}{\ensuremath{\mathbb{C}}}
\newunicodechar{ℕ}{\ensuremath{\mathbb{N}}}
\newunicodechar{ℚ}{\ensuremath{\mathbb{Q}}}
\newunicodechar{𝔻}{\ensuremath{\mathbb{D}}}
\newunicodechar{α}{\ensuremath{\alpha}}
\newunicodechar{β}{\ensuremath{\beta}}
\newunicodechar{γ}{\ensuremath{\gamma}}
\newunicodechar{δ}{\ensuremath{\delta}}
\newunicodechar{ε}{\ensuremath{\varepsilon}}
\newunicodechar{θ}{\ensuremath{\theta}}
\newunicodechar{λ}{\ensuremath{\lambda}}
\newunicodechar{μ}{\ensuremath{\mu}}
\newunicodechar{σ}{\ensuremath{\sigma}}
\newunicodechar{τ}{\ensuremath{\tau}}
\newunicodechar{φ}{\ensuremath{\varphi}}
\newunicodechar{ψ}{\ensuremath{\psi}}
\newunicodechar{ω}{\ensuremath{\omega}}
\newunicodechar{Σ}{\ensuremath{\Sigma}}
% majuscules produites par \MakeUppercase dans les titres (α-aminés -> Α)
\newunicodechar{Α}{\ensuremath{\alpha}}
\newunicodechar{Β}{\ensuremath{\beta}}
\newunicodechar{Γ}{\ensuremath{\Gamma}}
\newunicodechar{Δ}{\ensuremath{\Delta}}
\newunicodechar{Ε}{\ensuremath{\varepsilon}}
\newunicodechar{Θ}{\ensuremath{\Theta}}
\newunicodechar{Λ}{\ensuremath{\Lambda}}
\newunicodechar{Μ}{\ensuremath{\mu}}
\newunicodechar{Τ}{\ensuremath{\tau}}
\newunicodechar{Φ}{\ensuremath{\Phi}}
\newunicodechar{Ψ}{\ensuremath{\Psi}}
\newunicodechar{Ω}{\ensuremath{\Omega}}
\newunicodechar{↦}{\ensuremath{\mapsto}}
\newunicodechar{∈}{\ensuremath{\in}}
\newunicodechar{∉}{\ensuremath{\notin}}
\newunicodechar{∧}{\ensuremath{\wedge}}
\newunicodechar{⇔}{\ensuremath{\Leftrightarrow}}
\newunicodechar{⟺}{\ensuremath{\Longleftrightarrow}}
\newunicodechar{⇒}{\ensuremath{\Rightarrow}}
\newunicodechar{⟹}{\ensuremath{\Longrightarrow}}
\newunicodechar{∘}{\ensuremath{\circ}}
\newunicodechar{∪}{\ensuremath{\cup}}
\newunicodechar{∩}{\ensuremath{\cap}}
\newunicodechar{≡}{\ensuremath{\equiv}}
\newunicodechar{⊂}{\ensuremath{\subset}}
\newunicodechar{⊥}{\ensuremath{\perp}}
\newunicodechar{∃}{\ensuremath{\exists}}
\newunicodechar{∀}{\ensuremath{\forall}}
\newunicodechar{∛}{\ensuremath{\sqrt[3]{\,}}}
\newunicodechar{∜}{\ensuremath{\sqrt[4]{\,}}}
\newunicodechar{⃗}{\ensuremath{{}^{\to}}}
\newunicodechar{ⁿ}{\ifmmode^{n}\else\textsuperscript{\ensuremath{n}}\fi}
\newunicodechar{ˣ}{\ifmmode^{x}\else\textsuperscript{\ensuremath{x}}\fi}
\newunicodechar{ᵇ}{\ifmmode^{b}\else\textsuperscript{\ensuremath{b}}\fi}
\newunicodechar{ʳ}{\ifmmode^{r}\else\textsuperscript{\ensuremath{r}}\fi}
\newunicodechar{ᵘ}{\ifmmode^{u}\else\textsuperscript{\ensuremath{u}}\fi}
\newunicodechar{ʸ}{\ifmmode^{y}\else\textsuperscript{\ensuremath{y}}\fi}
\newunicodechar{ᵃ}{\ifmmode^{a}\else\textsuperscript{\ensuremath{a}}\fi}
\newunicodechar{ᵉ}{\ifmmode^{e}\else\textsuperscript{\ensuremath{e}}\fi}
\newunicodechar{ᵏ}{\ifmmode^{k}\else\textsuperscript{\ensuremath{k}}\fi}
\newunicodechar{ᵖ}{\ifmmode^{p}\else\textsuperscript{\ensuremath{p}}\fi}
\newunicodechar{ᴺ}{\ifmmode^{N}\else\textsuperscript{\ensuremath{N}}\fi}
\newunicodechar{ᶜ}{\ifmmode^{c}\else\textsuperscript{\ensuremath{c}}\fi}
\newunicodechar{ʰ}{\ifmmode^{h}\else\textsuperscript{\ensuremath{h}}\fi}
\newunicodechar{ᵠ}{\ifmmode^{q}\else\textsuperscript{\ensuremath{q}}\fi}
\newunicodechar{ₙ}{\ifmmode_{n}\else\textsubscript{\ensuremath{n}}\fi}
\newunicodechar{ₐ}{\ifmmode_{a}\else\textsubscript{\ensuremath{a}}\fi}
\newunicodechar{ᵢ}{\ifmmode_{i}\else\textsubscript{\ensuremath{i}}\fi}
\newunicodechar{ᵣ}{\ifmmode_{r}\else\textsubscript{\ensuremath{r}}\fi}
\newunicodechar{ₖ}{\ifmmode_{k}\else\textsubscript{\ensuremath{k}}\fi}
\newunicodechar{ᵦ}{\ifmmode_{\beta}\else\textsubscript{\ensuremath{\beta}}\fi}
\newunicodechar{ₓ}{\ifmmode_{x}\else\textsubscript{\ensuremath{x}}\fi}
\newfontfamily\popdisplay{ArchivoBlack-Regular.ttf}[Path=../../../fonts/]
\newfontfamily\poplogo{SpaceGrotesk-Bold.ttf}[Path=../../../fonts/]
\newfontfamily\popsemi{PlusJakartaSans-SemiBold.ttf}[Path=../../../fonts/]
\newfontfamily\popxbold{PlusJakartaSans-ExtraBold.ttf}[Path=../../../fonts/]
\newfontfamily\popxboldls{PlusJakartaSans-ExtraBold.ttf}[Path=../../../fonts/,LetterSpace=3.0]

\setlist[enumerate]{leftmargin=1.7em,itemsep=5pt,topsep=4pt}
\setlist[itemize]{leftmargin=1.7em,itemsep=4pt,topsep=4pt,label={\color{poporange}\textbullet}}
\setlength{\parindent}{0pt}
\setlength{\parskip}{5pt}
\renewcommand{\arraystretch}{1.25}

% pgfplots : style Pop par défaut
\pgfplotsset{
  every axis/.append style={axis line style={popink,line width=1pt},tick style={popink},
    grid style={popline,line width=0.5pt},label style={font=\small},tick label style={font=\footnotesize}},
  popcurve/.style={poporange,line width=1.8pt,no markers,smooth},
}
% Même style disponible dans un \draw TikZ simple (hors pgfplots)
\tikzset{popcurve/.style={poporange,line width=1.8pt}}
\tikzset{popgrid/.style={popline,very thin}}
\colorlet{popcurve}{poporange} % le nom du style est parfois utilisé comme couleur

% ============================================================
% Pill / squiggle
% ============================================================
\newcommand{\poppill}[3]{%
  \tikz[baseline=-0.55ex]\node[draw=popink,line width=1.2pt,rounded corners=2.6pt,%
    fill=#2,inner xsep=6.5pt,inner ysep=3pt]{%
    \popxboldls\fontsize{7}{7}\selectfont\color{#3}\MakeUppercase{#1}};%
}
\newcommand{\popsquiggle}[1]{%
  \tikz[baseline=-0.4ex]{
    \draw[#1,line width=2.6pt,line cap=round]
      plot [smooth] coordinates {(0,0.09) (0.34,0.01) (0.68,0.21) (1.02,0.01) (1.36,0.10)};
  }%
}
% Mot-clé surligné (marqueur orange sous le texte)
\newcommand{\cle}[1]{{\popxbold\color{popdark}#1}}

% ============================================================
% En-tête / pied
% ============================================================
\pagestyle{fancy}
\fancyhf{}
\renewcommand{\headrulewidth}{0pt}
\renewcommand{\footrulewidth}{0pt}
\lhead{\raisebox{-0.15ex}{\poplogo\fontsize{13}{13}\selectfont\color{popink}OnBuch\color{poporange}+}}
\rhead{\poppill{\DOCMATIERE\ \textbullet\ \DOCNIVEAU\ \textbullet\ Leçon \DOCLECON}{white}{popink}}
\lfoot{\popxbold\fontsize{7.6}{7.6}\selectfont\color{popmuted}\MakeUppercase{Cours signé OnBuch+}\ \textbullet\ {\popsemi\DOCTITRE}}
\rfoot{\tikz[baseline=-0.6ex]\node[rounded corners=8.5pt,fill=popink,minimum height=17pt,minimum width=17pt,inner xsep=5pt,inner sep=0]{\popxbold\fontsize{8}{8}\selectfont\color{popcream}\thepage\,/\,\pageref*{LastPage}};}

% ============================================================
% Titres
% ============================================================
\newcommand{\chaptitre}[2]{%
  \begin{tcolorbox}[enhanced,colback=poporange,colframe=popink,boxrule=2.6pt,arc=11pt,
      drop shadow=popink,left=15pt,right=15pt,top=12pt,bottom=11pt,before skip=3pt,after skip=18pt]
    \poppill{#2}{popink}{popcream}\par\vspace{9pt}
    {\raggedright\hyphenpenalty=10000\exhyphenpenalty=10000\popdisplay\fontsize{22}{24}\selectfont\color{popcream}\MakeUppercase{#1}\par}
  \end{tcolorbox}
}
% Section numérotée : pastille orange avec le numéro + titre Archivo
\newcounter{popsec}
\newcommand{\coursec}[1]{%
  \stepcounter{popsec}\setcounter{popsub}{0}%
  \par\vspace{16pt}\noindent
  \begin{minipage}{\linewidth}
  \tikz[baseline=-0.7ex]\node[draw=popink,line width=1.6pt,fill=poporange,rounded corners=4pt,minimum size=22pt,inner sep=0]{\popdisplay\fontsize{12}{12}\selectfont\color{popcream}\thepopsec};%
  \hspace{8pt}{\popdisplay\fontsize{15}{17}\selectfont\color{popink}\MakeUppercase{#1}}\par
  \vspace{4pt}\noindent{\color{poporange}\rule{36pt}{3.6pt}}
  \end{minipage}\par\nopagebreak\vspace{8pt}
}
\newcounter{popsub}
\newcommand{\courssub}[1]{%
  \stepcounter{popsub}%
  \par\vspace{9pt}\noindent{\popxbold\fontsize{11.5}{13}\selectfont\color{popink}{\color{poporange}\thepopsec.\thepopsub}\hspace{6pt}#1}\par\nopagebreak\vspace{4pt}
}

% ============================================================
% Boîtes pédagogiques — même recette : fond blanc, bordure épaisse colorée,
% ombre dure encre, badge plein. Titre optionnel [#1].
% ============================================================
\tcbset{popbase/.style={breakable,enhanced,colback=white,boxrule=2.3pt,arc=9pt,
  shadow={3pt}{-3pt}{0pt}{popink},left=14pt,right=14pt,top=11pt,bottom=11pt,before skip=11pt,after skip=15pt,
  fontupper=\popsemi\fontsize{10.6}{14.5}\selectfont\color{popink},
  fontlower=\popsemi\fontsize{10.6}{14.5}\selectfont\color{popink}}}
% Coche dessinée (le glyphe ✓ n'existe pas dans les polices du document)
\newcommand{\popcheck}[1]{\tikz[baseline=-0.5ex]\draw[#1,line width=1.6pt,line cap=round,line join=round](-0.13,0.0)--(-0.04,-0.09)--(0.13,0.1);}
\providecommand{\coche}{\popcheck{popgreen}}
\providecommand{\resultat}[1]{\par\smallskip\noindent\fcolorbox{popgreen}{popgreenL}{\parbox{\dimexpr\linewidth-2\fboxsep-2\fboxrule\relax}{\textbf{Résultat.} #1}}\par\smallskip}
\newcommand{\popboxhead}[3]{\poppill{#1}{#2}{white}\ifx\relax#3\relax\else\hspace{6pt}{\popxbold\fontsize{10.6}{12}\selectfont\color{popink}#3}\fi\par\vspace{7pt}}

\newtcolorbox{aretenir}[1][]{popbase,colframe=popgreen,before upper={\popboxhead{À retenir}{popgreen}{#1}}}
% Texte en anglais (document de lecture, dialogue, extrait) : cadre
% d'étude. Un titre optionnel (source, auteur) s'affiche dans l'en-tête.
\newtcolorbox{textebox}[1][]{popbase,colframe=popblue,before upper={\popboxhead{Text}{popblue}{#1}\begin{otherlanguage}{english}},after upper={\end{otherlanguage}}}
% Marques ✓ / ✗ (forme correcte / fautive) souvent utilisées par les modèles
\providecommand{\cross}{\textcolor{poppink}{\ensuremath{\times}}}
\providecommand{\xmark}{\cross}\providecommand{\crossmark}{\cross}\providecommand{\cmark}{\ensuremath{\checkmark}}
% Ligne de réponse / justification à compléter (inventée par les modèles)
\providecommand{\justification}[1]{\par\noindent\textit{Justification~:} \dotfill\par}
% \textipa (package tipa non chargé) : la police ne couvre pas l'API -> simple passage
\providecommand{\textipa}[1]{#1}
% Soulignement (ulem non chargé) : \uline, \ul
\providecommand{\uline}[1]{\underline{#1}}\providecommand{\ul}[1]{\underline{#1}}
% \glossaire{\item ...} inventée par les modèles : liste de glossaire
\providecommand{\glossaire}[1]{\begin{itemize}#1\end{itemize}}
% \alert (beamer) inventée par les modèles : texte mis en évidence
\providecommand{\alert}[1]{\textbf{\textcolor{poppink}{#1}}}
% variantes ASCII d'ordinaux écrites en macro
\providecommand{\iere}{\textsuperscript{re}}\providecommand{\ieme}{\textsuperscript{e}}\providecommand{\ere}{\textsuperscript{re}}\providecommand{\eme}{\textsuperscript{e}}\providecommand{\er}{\textsuperscript{er}}\providecommand{\ie}{\textsuperscript{e}}
% Ordinaux français écrits en macro par les modèles : 1\ière, 3\ième
\newcommand{\ière}{\textsuperscript{re}}\newcommand{\ième}{\textsuperscript{e}}
% \exemple (inventée par les modèles) : étiquette « Exemple : »
\providecommand{\exemple}{\textbf{Exemple~:}\ }
% \square utilisable aussi hors mode math (cases à cocher)
\AtBeginDocument{\let\squareMath\square\renewcommand{\square}{\ensuremath{\squareMath}}}
% Mot ou phrase en anglais à l'intérieur d'un paragraphe français
\newcommand{\eng}[1]{\ifmmode\text{\textit{#1}}\else\begingroup\selectlanguage{english}\itshape #1\endgroup\fi}
\let\esp\eng % alias toléré
% flèches d'intonation inventées par les modèles
\providecommand{\textsearrow}{}\renewcommand{\textsearrow}{\ensuremath{\searrow}}\providecommand{\textnearrow}{}\renewcommand{\textnearrow}{\ensuremath{\nearrow}}
% \fr{..} : mot français (inventé par les modèles)
\providecommand{\fr}[1]{\textit{#1}}
\newtcolorbox{exemplebox}[1][]{popbase,colframe=poporange,before upper={\popboxhead{Exemple}{poporange}{#1}}}
\newtcolorbox{definition}[1][]{popbase,colframe=popblue,before upper={\popboxhead{Définition}{popblue}{#1}}}
\newtcolorbox{propriete}[1][]{popbase,colframe=poppurple,before upper={\popboxhead{Propriété}{poppurple}{#1}}}
\newtcolorbox{methode}[1][]{popbase,colframe=popgold,before upper={\popboxhead{Méthode}{popgold}{#1}}}
\newtcolorbox{attention}[1][]{popbase,colframe=poppink,before upper={\popboxhead{Attention}{poppink}{#1}}}
\newtcolorbox{experience}[1][]{popbase,colframe=popblue,colback=popblueL,before upper={\popboxhead{Expérience}{popblue}{#1}}}
% « Le savais-tu ? » : carte violette pleine (contexte, culture, Cameroun)
\newtcolorbox{savaistu}[1][]{popbase,colback=poppurple,colframe=popink,
  fontupper=\popsemi\fontsize{10.4}{14.2}\selectfont\color{white},
  before upper={\poppill{Le savais-tu\,?}{white}{poppurple}\ifx\relax#1\relax\else\hspace{6pt}{\popxbold\color{white}#1}\fi\par\vspace{7pt}}}
% Exercice résolu : énoncé en haut, \tcblower puis la solution sur fond crème
\newcounter{popexo}\newcounter{popexr}
\newtcolorbox{exoresolu}[1][]{popbase,colframe=poporange,colbacklower=popcream,
  segmentation style={popink,line width=1.2pt,dashed},
  before upper={\stepcounter{popexr}\popboxhead{Exercice résolu \thepopexr}{poporange}{#1}},
  before lower={\poppill{Solution}{popink}{popcream}\par\vspace{6pt}}}
% Exercice (non résolu en place) — la correction est dans la partie Corrigés
\newtcolorbox{exercice}[1][]{popbase,colframe=popink,
  before upper={\stepcounter{popexo}\popboxhead{Exercice \thepopexo}{popink}{#1}}}
% Corrigé
\newtcolorbox{corrige}[1][]{popbase,colframe=popgreen,colback=popgreenL,
  before upper={\popboxhead{Corrigé}{popgreen}{#1}}}
% Objectifs / prérequis / bilan
\newtcolorbox{objectifs}{popbase,colframe=popink,colback=poporangeL,
  before upper={\popboxhead{Objectifs de la leçon}{popink}{}}}
\newtcolorbox{prerequis}{popbase,colframe=popink,colback=popblueL,
  before upper={\popboxhead{Prérequis}{popblue}{}}}
\newtcolorbox{bilan}{popbase,colframe=popink,colback=white,boxrule=2.8pt,
  before upper={\popboxhead{Fiche bilan}{poporange}{L'essentiel à connaître pour l'examen}}}
% Figure encadrée avec légende
\newtcolorbox{popfigure}[1][]{enhanced,breakable=false,colback=white,colframe=popline,boxrule=1.4pt,arc=8pt,
  left=10pt,right=10pt,top=10pt,bottom=8pt,before skip=10pt,after skip=12pt,halign=center,
  lower separated=false,halign lower=center,
  fontlower=\popsemi\fontsize{9}{11.5}\selectfont\color{popmuted},
  #1}
\newcommand{\legende}[1]{\par\vspace{6pt}{\popsemi\fontsize{9}{11.5}\selectfont\color{popmuted}#1}}

% ============================================================
% Signature OnBuch+
% ============================================================
\newcommand{\poplogoinline}[1]{{\poplogo\fontsize{#1}{#1}\selectfont\color{popink}OnBuch\color{poporange}+}}
\newcommand{\signatureonbuch}{%
  \begin{tcolorbox}[enhanced,colback=popink,colframe=popink,boxrule=2.2pt,arc=10pt,
      drop shadow=poporange,left=14pt,right=14pt,top=10pt,bottom=10pt,before skip=0pt,after skip=0pt]
    \tikz[baseline=-0.7ex]\node[circle,fill=popgreen,draw=popcream,line width=1.3pt,minimum size=22pt,inner sep=0]{\popcheck{white}};%
    \hspace{9pt}\begin{minipage}[c]{0.62\linewidth}
      {\popxbold\fontsize{12}{13}\selectfont\color{popcream}Signé\ \textbullet\ {\poplogo OnBuch\color{poporange}+}}\par\vspace{2pt}
      {\raggedright\popsemi\fontsize{8}{10.5}\selectfont\color{popline}\MakeUppercase{Équipe pédagogique OnBuch+}\\\MakeUppercase{Conforme au programme officiel MINESEC}\par}
    \end{minipage}\hfill
    {\poplogo\fontsize{22}{22}\selectfont\color{popcream}OnBuch\color{poporange}+}
  \end{tcolorbox}%
}

% ============================================================
% Page de garde
% #1 titre · #2 matière · #3 niveau · #4 module · #5 n° leçon · #6 durée ·
% #7 description / objectifs (texte ou itemize)
% ============================================================
\newcommand{\pagegarde}[7]{%
  \thispagestyle{empty}
  \newgeometry{a4paper,margin=1.7cm,top=1.6cm,bottom=1.4cm}%
  \begin{tikzpicture}[remember picture,overlay]
    \fill[poporange!18] ([xshift=-3.2cm,yshift=-3.6cm]current page.north east) circle (4.6cm);
    \fill[poppurple!14] ([xshift=2.4cm,yshift=7.6cm]current page.south west) circle (3.8cm);
  \end{tikzpicture}%
  {\poplogo\fontsize{30}{30}\selectfont\color{popink}OnBuch\color{poporange}+}\hfill
  \raisebox{0.6ex}{\poppill{Cours signé \textbullet\ Premium}{popink}{popcream}}\par
  \vspace{1pt}\popsquiggle{poppurple}\par
  \vspace{8pt}
  \begin{tcolorbox}[enhanced,colback=poporange,colframe=popink,boxrule=2.8pt,arc=13pt,
      drop shadow=popink,left=18pt,right=18pt,top=12pt,bottom=13pt,after skip=10pt]
    \poppill{#2 \textbullet\ #3}{popink}{popcream}\hspace{5pt}\poppill{#4}{white}{popink}\par\vspace{8pt}
    {\popdisplay\fontsize{14}{14}\selectfont\color{popink}LEÇON #5}\par\vspace{4pt}
    {\raggedright\hyphenpenalty=10000\exhyphenpenalty=10000\popdisplay\fontsize{25}{27}\selectfont\color{popcream}\MakeUppercase{#1}\par}
    \vspace{8pt}
    {\popxbold\fontsize{9.5}{9.5}\selectfont\color{popink}\MakeUppercase{Durée conseillée : #6}}
  \end{tcolorbox}
  \begin{tcolorbox}[enhanced,colback=white,colframe=popink,boxrule=2.2pt,arc=10pt,
      drop shadow=popink,left=16pt,right=16pt,top=10pt,bottom=10pt,after skip=8pt]
    \poppill{Dans cette leçon}{poppurple}{white}\par\vspace{5pt}
    \popsemi\fontsize{9.3}{12.6}\selectfont\color{popink}#7
  \end{tcolorbox}
  \noindent\begin{minipage}[t]{0.487\linewidth}
    \begin{tcolorbox}[enhanced,colback=popgreen,colframe=popink,boxrule=2.2pt,arc=10pt,
        drop shadow=popink,left=11pt,right=11pt,top=10pt,bottom=10pt,height=3.1cm,valign=center]
      \tcbox[colback=white,colframe=popink,boxrule=1.3pt,arc=5pt,boxsep=0pt,nobeforeafter,
        left=3pt,right=3pt,top=3pt,bottom=3pt,valign=center]{\qrcode[height=1.9cm]{https://play.google.com/store/apps/details?id=com.luvvix.onbuch&hl=fr}}%
      \hspace{8pt}\begin{minipage}[c]{0.5\linewidth}
        {\popdisplay\fontsize{11}{12.5}\selectfont\color{white}\MakeUppercase{Télécharge OnBuch}}\par\vspace{4pt}
        {\raggedright\popsemi\fontsize{8.4}{11}\selectfont\color{white}Gratuit \textbullet\ Google Play\\Tuteur IA, annales, concours, résultats\par}
      \end{minipage}
    \end{tcolorbox}
  \end{minipage}\hfill
  \begin{minipage}[t]{0.487\linewidth}
    \begin{tcolorbox}[enhanced,colback=poppurple,colframe=popink,boxrule=2.2pt,arc=10pt,
        drop shadow=popink,left=11pt,right=11pt,top=10pt,bottom=10pt,height=3.1cm,valign=center]
      \tikz[baseline=-0.7ex]\node[circle,fill=white,draw=popink,line width=1.3pt,minimum size=1.6cm,inner sep=0]{\popdisplay\fontsize{24}{24}\selectfont\color{poppurple}+};%
      \hspace{8pt}\begin{minipage}[c]{0.6\linewidth}
        {\popdisplay\fontsize{11}{12.5}\selectfont\color{white}\MakeUppercase{Passe à OnBuch+}}\par\vspace{4pt}
        {\raggedright\popsemi\fontsize{8.4}{11}\selectfont\color{white}Tous les cours signés, Tuteur IA illimité, fiches PDF — onglet OnBuch+ de l'app\par}
      \end{minipage}
    \end{tcolorbox}
  \end{minipage}
  \vspace{8pt}
  \popcontenu
  \vfill
  \signatureonbuch
  \restoregeometry
  \newpage
}

% Rangée « Ce cours contient » (page de garde)
\newcommand{\poptile}[3]{% #1 couleur #2 gros texte #3 légende
  \begin{tcolorbox}[enhanced,colback=#1,colframe=popink,boxrule=2pt,arc=9pt,shadow={2.5pt}{-2.5pt}{0pt}{popink},
      left=6pt,right=6pt,top=9pt,bottom=9pt,halign=center,valign=center,height=1.75cm,width=\linewidth,nobeforeafter]
    {\popdisplay\fontsize{12.5}{13}\selectfont\color{white}\MakeUppercase{#2}}\par\vspace{3pt}
    {\centering\popsemi\fontsize{7.8}{9.5}\selectfont\color{white}#3\par}
  \end{tcolorbox}}
\newcommand{\popcontenu}{%
  \par\noindent{\popxboldls\fontsize{7.5}{7.5}\selectfont\color{popmuted}CE COURS SIGNÉ CONTIENT}\par\vspace{7pt}
  \noindent\begin{minipage}[t]{0.235\linewidth}\poptile{popblue}{Cours}{complet, illustré, pas à pas}\end{minipage}\hfill
  \begin{minipage}[t]{0.235\linewidth}\poptile{poporange}{Méthodes}{et exercices résolus}\end{minipage}\hfill
  \begin{minipage}[t]{0.235\linewidth}\poptile{poppink}{Exercices}{type examen + corrigés}\end{minipage}\hfill
  \begin{minipage}[t]{0.235\linewidth}\poptile{popgreen}{Fiche bilan}{l'essentiel à retenir}\end{minipage}\par}

% Sommaire
\newtcolorbox{sommaire}{enhanced,colback=white,colframe=popink,boxrule=2.2pt,arc=10pt,
  drop shadow=popink,left=18pt,right=18pt,top=13pt,bottom=13pt,before skip=6pt,after skip=20pt,
  before upper={\poppill{Sommaire}{poppurple}{white}\par\vspace{9pt}\popsemi\fontsize{10.6}{14.5}\selectfont\color{popink}}}

% Signature de fin de cours — poussée tout en bas de la dernière page
\newcommand{\findecours}{%
  \par\vfill\nopagebreak
  \begin{center}\popsquiggle{poporange}\end{center}\vspace{4pt}
  \signatureonbuch
}
\AtBeginDocument{\def\llbracket{\mathopen{[\mkern-3.5mu[}}\def\rrbracket{\mathclose{]\mkern-3.5mu]}}} % intervalles d'entiers (glyphe ⟦ absent de la police)
% Glyphes absents de Fira Math : redéfinis à partir de glyphes disponibles
\AtBeginDocument{%
  \def\setminus{\mathbin{\backslash}}\let\smallsetminus\setminus
  \def\vdots{\mathord{\vbox{\baselineskip3.2pt\lineskiplimit0pt\kern4pt\hbox{.}\hbox{.}\hbox{.}}}}%
  \def\ddots{\mathinner{\mkern1mu\raise6.5pt\hbox{.}\mkern2mu\raise3.6pt\hbox{.}\mkern2mu\raise0.7pt\hbox{.}\mkern1mu}}%
  \def\triangle{\mathord{\scalebox{0.9}{$\symup{\Delta}$}}}%
  \def\nabla{\mathord{\raisebox{\depth}{\scalebox{1}[-1]{$\symup{\Delta}$}}}}%
  \def\checkmark{\popcheck{popdark}}%
  \def\star{\mathbin{\ast}}\def\bigstar{\mathord{\ast}}%
  \def\diamond{\mathbin{\scalebox{0.6}{$\Diamond$}}}%
  \def\bigcup{\mathop{\mathchoice{\raisebox{-0.3em}{\scalebox{1.7}{$\cup$}}}{\scalebox{1.25}{$\cup$}}{\cup}{\cup}}\displaylimits}%
  \def\bigcap{\mathop{\mathchoice{\raisebox{-0.3em}{\scalebox{1.7}{$\cap$}}}{\scalebox{1.25}{$\cap$}}{\cap}{\cap}}\displaylimits}%
  \def\lesseqqgtr{\mathrel{\lessgtr}}\def\bigoplus{\mathop{\oplus}\displaylimits}%
}
\providecommand{\ssi}{si et seulement si} % abréviation fréquente
\AtBeginDocument{\providecommand{\wideparen}{\overparen}} % arc de cercle

%% FIGFIX-BEGIN v5 — correctifs globaux des figures (docs/onbuch-generation/tools/figfix)
\usepackage{adjustbox}\usepackage{etoolbox}\tcbuselibrary{raster}
% toute figure TikZ/pgfplots plus large que la ligne est réduite à la largeur disponible
\BeforeBeginEnvironment{tikzpicture}{\begin{adjustbox}{max width=\linewidth}}
\AfterEndEnvironment{tikzpicture}{\end{adjustbox}}
% légendes pgfplots lisibles par-dessus les courbes
\pgfplotsset{every axis/.append style={legend style={fill=white,fill opacity=0.92,text opacity=1,draw=black!15,inner sep=3pt}}}
% étiquettes d'arêtes (syntaxe edge["..."]) avec fond blanc
\tikzset{every edge quotes/.append style={fill=white,inner sep=1.2pt}}
%% FIGFIX-END
\begin{document}

\pagegarde{Lesson 1 — Speaking: Exchanges information about interacting with mates, neighbours, adults and school authorities on community celebrations/festivals}{Anglais}{1ère A}{Chapter 1 : Family and Social Life}{ENA01}{2 h}{Cette leçon d'expression orale apprend à échanger des informations sur les fêtes et célébrations communautaires (Ngondo, Youth Day, festivals, carnavals) avec des camarades, des voisins, des adultes et les autorités scolaires. À travers des dialogues modèles, des formules fonctionnelles et des jeux de rôle guidés, l'élève s'entraîne à poser des questions, répondre, relancer et conclure une conversation en anglais avec une prononciation soignée.
\vspace{4pt}
\begin{multicols}{2}\small
\begin{itemize}
\item À la fin de cette leçon, je sais saluer et engager une conversation adaptée à mon interlocuteur (camarade, voisin, adulte, autorité scolaire)
\item À la fin de cette leçon, je sais demander des informations sur une fête communautaire (date, lieu, programme, organisateurs)
\item À la fin de cette leçon, je sais donner des informations précises sur une célébration de ma communauté
\item À la fin de cette leçon, je sais relancer une conversation qui ralentit avec des questions ouvertes
\item À la fin de cette leçon, je sais conclure poliment un échange oral
\item À la fin de cette leçon, je sais adapter mon niveau de langue : familier avec un camarade, formel avec le proviseur
\end{itemize}\end{multicols}}

\begin{sommaire}
\begin{multicols}{2}
\begin{enumerate}
\item Engager la conversation selon l'interlocuteur
\item Demander des informations sur une fête
\item Donner des informations et décrire une célébration
\item Relancer et conclure la conversation
\item Dialogues modèles complets
\item Prononciation, accentuation et intonation
\item Jeux de rôle guidés et entraînement
\item Activité d'intégration
\item Méthodes et exercices résolus
\item Exercices
\item Corrigés des exercices
\item Fiche bilan
\end{enumerate}
\end{multicols}
\end{sommaire}

\begin{objectifs}
\begin{itemize}
\item À la fin de cette leçon, je sais saluer et engager une conversation adaptée à mon interlocuteur (camarade, voisin, adulte, autorité scolaire)
\item À la fin de cette leçon, je sais demander des informations sur une fête communautaire (date, lieu, programme, organisateurs)
\item À la fin de cette leçon, je sais donner des informations précises sur une célébration de ma communauté
\item À la fin de cette leçon, je sais relancer une conversation qui ralentit avec des questions ouvertes
\item À la fin de cette leçon, je sais conclure poliment un échange oral
\item À la fin de cette leçon, je sais adapter mon niveau de langue : familier avec un camarade, formel avec le proviseur
\item À la fin de cette leçon, je sais prononcer correctement le lexique des fêtes avec la bonne accentuation et la bonne intonation
\item À la fin de cette leçon, je sais mener un court dialogue complet et naturel en anglais sur un festival
\end{itemize}
\end{objectifs}

\begin{prerequis}
\begin{itemize}
\item Les salutations et formules de politesse de base (Seconde)
\item Les questions en Wh- (what, when, where, who, how) et en yes/no avec auxiliaires
\item Le present simple et le futur avec 'will' et 'going to' pour parler d'événements
\item Le vocabulaire de base de la famille et de la vie sociale (Chapter 1)
\item L'intonation montante et descendante dans les questions (début d'année)
\end{itemize}
\end{prerequis}

\newpage

\coursec{Engager la conversation selon l'interlocuteur}

C'est la période de la fête du \eng{Ngondo} à Douala, et ton quartier prépare aussi le \eng{Youth Day} du 11 février. Ton voisin \eng{Mr Ewusi} veut savoir qui organise la fanfare, la directrice de ton lycée t'interroge sur le programme, et tes camarades te demandent l'heure du défilé. En Grande-Bretagne, les élèves discutent de la même façon du carnaval de \eng{Notting Hill} avec leurs professeurs et leurs voisins. Sauras-tu, en anglais, demander et donner des informations, relancer la conversation et conclure poliment avec un camarade, un adulte ou une autorité scolaire ?

Toute la leçon va répondre à cette question, étape par étape. Mais avant de demander quoi que ce soit, il faut d'abord \cle{engager la conversation} correctement : c'est l'objet de cette première section.

\courssub{Observer : deux façons d'aborder quelqu'un}

Lisons d'abord ce mini-dialogue. Il se passe dans la cour du lycée, entre une élève de Première et le proviseur.

\begin{textebox}[At the principal's office — a mini-dialogue]
\textbf{Student:} Good morning, Sir. May I ask you about the Youth Day programme?

\textbf{Principal:} Good morning. Of course. What would you like to know?

\textbf{Student:} Could you please tell me when the parade starts?

\textbf{Principal:} It starts at eight o'clock, in front of the grandstand.

\textbf{Student:} Really? That's interesting. And who is organising the brass band?

\textbf{Principal:} The music club, with Mr Fon. Is there anything else?

\textbf{Student:} No, Sir. Thank you very much.

\textbf{Principal:} You're welcome. Have a nice day.
\end{textebox}

\textbf{Glossaire :} \eng{parade} (nom, défilé) ; \eng{grandstand} (nom, tribune) ; \eng{brass band} (nom, fanfare) ; \eng{anything else} (locution, autre chose).

\textbf{Questions de compréhension :}
\begin{enumerate}
\item How does the student greet the principal?
\item What two questions does the student ask?
\item How does the student show interest in the answer?
\end{enumerate}

Observons maintenant la même scène entre deux camarades, dans la cour de récréation :

\begin{exemplebox}[Between mates]
\eng{“Hi, Sandra! What's up?”} — « Salut Sandra ! Quoi de neuf ? »

\eng{“Have you heard about the Ngondo festival? It starts on Sunday.”} — « Tu as entendu parler de la fête du Ngondo ? Elle commence dimanche. »

\eng{“Really? Wonderful! What time does the procession begin?”} — « Vraiment ? Génial ! À quelle heure commence la procession ? »
\end{exemplebox}

\textbf{Ce que nous observons :} les mêmes fonctions (saluer, introduire le sujet, poser une question) sont réalisées avec des mots très différents selon l'interlocuteur. \eng{Good morning, Sir} contre \eng{Hi!} ; \eng{Could you please tell me...?} contre \eng{What time does...?}. C'est exactement ce choix que nous allons apprendre à maîtriser.

\courssub{Les quatre types d'interlocuteurs}

Dans la vie scolaire et sociale au Cameroun comme dans les pays anglophones, tu parles chaque jour à quatre grandes catégories de personnes.

\begin{definition}[Les interlocuteurs de la vie quotidienne]
\begin{itemize}
\item \cle{Mates} : camarades de classe, amis, frères et sœurs. Registre \textbf{informel}.
\item \cle{Neighbours} : voisins du quartier, connaissances. Registre \textbf{neutre}.
\item \cle{Adults} : adultes de la famille ou de la communauté (parents, oncles, chefs de quartier). Registre \textbf{formel}.
\item \cle{School authorities} : le proviseur (\eng{the principal}), le censeur (\eng{the vice-principal}), les enseignants (\eng{teachers}). Registre \textbf{formel}.
\end{itemize}
\end{definition}

\begin{definition}[Register — le registre de langue]
Le \cle{register} (registre de langue) est le niveau de langue que l'on choisit selon la personne à qui l'on parle : familier (\eng{informal}), courant (\eng{neutral}) ou soutenu (\eng{formal}). En anglais, le pronom \eng{you} ne change jamais — il n'y a pas de « tu » et de « vous » — mais le \textbf{ton}, les \textbf{salutations} et les \textbf{tournures} changent complètement.
\end{definition}

\begin{attention}[Pas de « tu/vous » en anglais... mais attention au ton !]
Les francophones croient souvent que l'anglais est « plus facile » parce qu'on dit \eng{you} à tout le monde. Erreur ! La politesse se déplace ailleurs : dans les salutations (\eng{Hi} vs \eng{Good morning, Sir}), dans les auxiliaires (\eng{Can you...?} vs \eng{Could you please...?}) et dans les titres d'appel (\eng{Sir}, \eng{Madam}, \eng{Mr Ewusi}). Un élève qui dit \eng{“Hey! Can you give me the programme?”} à son proviseur est aussi impoli qu'un élève qui tutoierait son proviseur en français.
\end{attention}

\courssub{Saluer selon l'interlocuteur}

\begin{propriete}[Salutations informelles (mates)]
Entre camarades, on utilise des formules courtes et décontractées :
\begin{itemize}
\item \eng{Hi!} / \eng{Hello!} — « Salut ! »
\item \eng{How are you doing?} — « Comment ça va ? » (réponse attendue : \eng{Fine, thanks. And you?})
\item \eng{What's up?} — « Quoi de neuf ? » (réponse possible : \eng{Not much!})
\item \eng{Good to see you!} — « Content de te voir ! »
\end{itemize}
\end{propriete}

\begin{propriete}[Salutations formelles (adults, school authorities)]
Avec un adulte ou une autorité, on utilise des formules complètes, souvent accompagnées d'un titre d'appel :
\begin{itemize}
\item \eng{Good morning, Sir.} / \eng{Good afternoon, Madam.} — « Bonjour, Monsieur / Madame. »
\item \eng{Good morning, Mr Ewusi.} / \eng{Good evening, Mrs Ngono.} — avec le nom de famille.
\item \eng{How do you do?} — formule très formelle de première présentation.
\item \eng{Excuse me, please.} — « Excusez-moi, s'il vous plaît », pour aborder poliment quelqu'un.
\end{itemize}
\textbf{Règle d'orthographe :} les titres \eng{Sir}, \eng{Madam}, \eng{Mr}, \eng{Mrs}, \eng{Dr} prennent toujours la majuscule.
\end{propriete}

\begin{attention}[Ne dis jamais “Good morning, teacher!”]
Erreur très fréquente chez les francophones : traduire « Bonjour, professeur ! » par \eng{Good morning, teacher!}. En anglais, \eng{teacher} est un \textbf{métier}, pas un titre d'appel. On dit :
\begin{itemize}
\item \eng{Good morning, Sir.} / \eng{Good morning, Madam.} (à l'école)
\item \eng{Good morning, Mr Fon.} / \eng{Good morning, Mrs Acha.} (si l'on connaît le nom)
\end{itemize}
De même, on ne dit pas \eng{“Good morning, principal!”} mais \eng{“Good morning, Sir.”}
\end{attention}

\begin{center}
\begin{tabularx}{\linewidth}{|X|X|X|}
\hline
\rowcolor{popblueL} Interlocuteur & Salutation anglaise & Équivalent français \\
\hline
Mate (camarade) & Hi! What's up? & Salut ! Quoi de neuf ? \\
\hline
Neighbour (voisin) & Hello, Mr Ewusi. How are you? & Bonjour, M. Ewusi. Comment allez-vous ? \\
\hline
Adult (adulte) & Good afternoon, Madam. & Bonjour, Madame. \\
\hline
School authority & Good morning, Sir. Excuse me, please. & Bonjour, Monsieur. Excusez-moi, s'il vous plaît. \\
\hline
\end{tabularx}
\end{center}

\begin{popfigure}
\begin{center}\tcbox[colback=popblue,colframe=popblue,coltext=white,arc=9pt,boxsep=4pt,left=6pt,right=6pt]{\bfseries\small Starting a conversation}\end{center}
\vspace{-2pt}
\begin{tcbraster}[raster columns=2,raster equal height=rows,raster column skip=6pt,raster row skip=6pt,enhanced,blankest]
\begin{tcolorbox}[enhanced,colback=popblue,colframe=popblue,coltext=white,boxrule=0.9pt,arc=3pt,left=4pt,right=4pt,top=3pt,bottom=3pt,boxsep=1pt,halign=left]\bfseries\scriptsize Mate \emph{“Hi!} \emph{What's up?”}\end{tcolorbox}
\begin{tcolorbox}[enhanced,colback=poporange,colframe=poporange,coltext=white,boxrule=0.9pt,arc=3pt,left=4pt,right=4pt,top=3pt,bottom=3pt,boxsep=1pt,halign=left]\bfseries\scriptsize Neighbour \emph{“Hello,} \emph{Mr Ewusi!”}\end{tcolorbox}
\begin{tcolorbox}[enhanced,colback=popgreen,colframe=popgreen,coltext=white,boxrule=0.9pt,arc=3pt,left=4pt,right=4pt,top=3pt,bottom=3pt,boxsep=1pt,halign=left]\bfseries\scriptsize Adult \emph{“Good afternoon,} \emph{Madam.”}\end{tcolorbox}
\begin{tcolorbox}[enhanced,colback=poppurple,colframe=poppurple,coltext=white,boxrule=0.9pt,arc=3pt,left=4pt,right=4pt,top=3pt,bottom=3pt,boxsep=1pt,halign=left]\bfseries\scriptsize School authority \emph{“Good morning,} \emph{Sir.”}\end{tcolorbox}
\end{tcbraster}

\legende{Carte mentale : une salutation adaptée à chaque interlocuteur.}
\end{popfigure}

\courssub{Introduire le sujet de la conversation}

Après la salutation, on annonce de quoi l'on veut parler. C'est la \cle{transition} vers le sujet : ici, une fête ou une célébration communautaire.

\begin{propriete}[Formules pour introduire le sujet]
\textbf{Registre formel} (adultes, autorités) :
\begin{itemize}
\item \eng{I'd like to ask you about the Youth Day celebration.} — « J'aimerais vous poser des questions sur la fête de la Jeunesse. »
\item \eng{May I ask you about the Ngondo festival?} — « Puis-je vous interroger au sujet de la fête du Ngondo ? »
\item \eng{Excuse me, could I talk to you about the programme?} — « Excusez-moi, pourrais-je vous parler du programme ? »
\end{itemize}
\textbf{Registre informel} (camarades) :
\begin{itemize}
\item \eng{Have you heard about the Ngondo festival?} — « Tu as entendu parler de la fête du Ngondo ? »
\item \eng{Guess what! The parade is on Saturday.} — « Devine quoi ! Le défilé est samedi. »
\item \eng{Do you know anything about the Youth Day match?} — « Tu sais quelque chose sur le match de la fête de la Jeunesse ? »
\end{itemize}
\end{propriete}

\begin{attention}[“I'd like” : la contraction polie]
\eng{I'd like} est la contraction de \eng{I would like}. C'est la forme polie de \eng{I want}. Ne dis jamais \eng{“I want to ask you...”} à une autorité : c'est trop direct, presque agressif. Compare : \eng{I want the programme.} (« Je veux le programme. » — impoli) et \eng{I'd like to see the programme, please.} (« J'aimerais voir le programme, s'il vous plaît. » — poli).
\end{attention}

\courssub{Réagir et montrer son intérêt}

Une conversation n'est pas un interrogatoire. Quand l'autre répond, il faut \cle{réagir} pour montrer qu'on écoute et l'encourager à continuer.

\begin{exemplebox}[Réactions pour montrer son intérêt]
\eng{Really?} — « Vraiment ? »

\eng{That's interesting!} — « C'est intéressant ! »

\eng{I see.} — « Je vois. »

\eng{Wonderful!} / \eng{Great!} — « Merveilleux ! » / « Super ! »

\eng{Oh, I didn't know that!} — « Oh, je ne savais pas ! »
\end{exemplebox}

\begin{exemplebox}[Exemples traduits — registre informel et registre formel]
\eng{What's up, man?} — « Quoi de neuf, mon pote ? » (entre camarades uniquement)

\eng{Could you please tell me when the festival starts?} — « Pourriez-vous me dire quand le festival commence, s'il vous plaît ? » (à un adulte ou une autorité)

\eng{Can you tell me the time of the match?} — « Tu peux me dire l'heure du match ? » (à un camarade)

\eng{How do you do, Mrs Ngono?} — « Enchanté, Madame Ngono. » (première rencontre, très formel)
\end{exemplebox}

\begin{popfigure}
\begin{tikzpicture}[font=\small]
\draw[-{Stealth}, very thick, popmuted] (0,0) -- (0,6.2);
\node[rotate=90, above, popmuted] at (-0.35,3.1) {register of language};
\fill[popgreenL] (0.3,0) rectangle (7.5,1.8);
\node[align=left, anchor=west] at (0.5,0.9) {\textbf{INFORMAL}\quad \emph{“Hi! What's up? Can you tell me the time?”}};
\fill[popgoldL] (0.3,2.1) rectangle (7.5,3.9);
\node[align=left, anchor=west] at (0.5,3.0) {\textbf{NEUTRAL}\quad \emph{“Hello, Mr Ewusi. Do you know who organises the band?”}};
\fill[popblueL] (0.3,4.2) rectangle (7.5,6.0);
\node[align=left, anchor=west] at (0.5,5.1) {\textbf{FORMAL}\quad \emph{“Good morning, Sir. Could you please tell me about the programme?”}};
\node[align=center, popgreen] at (8.4,0.9) {mates};
\node[align=center, popgold] at (8.4,3.0) {neighbours};
\node[align=center, popblue] at (8.4,5.1) {adults,\\authorities};
\end{tikzpicture}
\legende{L'échelle du registre : plus l'interlocuteur est « haut » socialement, plus la langue monte en formalité.}
\end{popfigure}

\courssub{Méthode : engager une conversation en trois étapes}

\begin{methode}[Les trois étapes pour engager la conversation]
\begin{enumerate}
\item \textbf{Saluer selon l'interlocuteur.} Choisis la salutation adaptée : \eng{Hi!} pour un camarade, \eng{Good morning, Sir.} pour une autorité. Ajoute éventuellement \eng{Excuse me, please.} pour aborder quelqu'un d'occupé.
\item \textbf{Introduire le sujet.} Annonce le thème de la fête ou de la célébration : \eng{I'd like to ask you about the Youth Day celebration.} (formel) ou \eng{Have you heard about the Ngondo festival?} (informel).
\item \textbf{Poser une première question ouverte.} Une question ouverte ne se répond pas par « oui » ou « non » : elle commence par \eng{What}, \eng{When}, \eng{Who}, \eng{Where}, \eng{How}. Exemple : \eng{What would you like to know?} / \eng{Who is organising the brass band?}
\end{enumerate}
\end{methode}

\begin{exoresolu}[Choisir la bonne salutation et la bonne formule]
\textbf{Énoncé.} Pour chaque situation, écris en anglais la salutation et la phrase d'introduction du sujet qui conviennent.

1. Tu abordes ta directrice (\eng{Mrs Acha}) pour parler du programme du Youth Day.

2. Tu demandes à ton camarade Tabi s'il connaît l'heure du défilé du Ngondo.

3. Tu t'adresses à ton voisin, \eng{Mr Ewusi}, pour savoir qui organise la fanfare du quartier.

\tcblower
\textbf{Solution détaillée.}

1. Interlocuteur = autorité scolaire, donc registre \textbf{formel} : salutation complète + titre d'appel + \eng{I'd like to...} ou \eng{May I...?}

\eng{Good morning, Madam. I'd like to ask you about the Youth Day programme.}

2. Interlocuteur = camarade, donc registre \textbf{informel} : \eng{Hi} + question directe avec \eng{Do you know...?}

\eng{Hi, Tabi! Do you know what time the Ngondo parade starts?}

3. Interlocuteur = voisin adulte, donc registre \textbf{neutre} : \eng{Hello} + nom + formule polie mais simple.

\eng{Hello, Mr Ewusi. Could you tell me who is organising the brass band, please?}

\textbf{Méthode :} identifier l'interlocuteur → choisir le niveau sur l'échelle du registre → salutation + introduction + question ouverte.
\end{exoresolu}

\begin{savaistu}[“How do you do?” : la question qui n'en est pas une]
Au Royaume-Uni, lorsqu'on est présenté à quelqu'un pour la première fois dans un contexte formel, on dit \eng{“How do you do?”}. Attention : ce n'est \textbf{pas} une vraie question sur la santé ! La réponse traditionnelle est... \eng{“How do you do?”} elle-même, souvent accompagnée d'une poignée de main. Si tu réponds \eng{“I'm fine, thank you”}, on comprendra que tu es étranger — ce qui n'est pas grave, mais tu connais maintenant le code ! Aujourd'hui, les jeunes Britanniques préfèrent souvent \eng{“Nice to meet you.”} (« Enchanté de faire ta connaissance. »), plus naturel.
\end{savaistu}

\begin{exercice}[À toi de jouer]
1. Classe ces formules sur l'échelle du registre (informal / neutral / formal) : \eng{What's up?} — \eng{Good afternoon, Sir.} — \eng{Hello, Mrs Ngono.} — \eng{Could you please help me?} — \eng{Hi, man!}

2. Corrige les erreurs de registre : \eng{“Good morning, teacher! I want to ask you about the festival.”}

3. Traduis en anglais : « Bonjour, Monsieur. J'aimerais vous poser une question sur la fête du Ngondo. »
\end{exercice}

\begin{corrige}[Exercice « À toi de jouer »]
1. \emph{Informal :} \eng{What's up?}, \eng{Hi, man!} — \emph{Neutral :} \eng{Hello, Mrs Ngono.} — \emph{Formal :} \eng{Good afternoon, Sir.}, \eng{Could you please help me?}

2. Deux erreurs : \eng{teacher} n'est pas un titre d'appel, et \eng{I want} est trop direct. Correction : \eng{“Good morning, Sir. I'd like to ask you about the festival.”}

3. \eng{Good morning, Sir. I'd like to ask you a question about the Ngondo festival.}
\end{corrige}

\begin{aretenir}[L'essentiel de la section]
\begin{itemize}
\item En anglais, \eng{you} ne change pas, mais le \textbf{registre} change tout : salutations, titres d'appel, tournures.
\item Quatre interlocuteurs : \eng{mates} (informel), \eng{neighbours} (neutre), \eng{adults} et \eng{school authorities} (formel).
\item Jamais \eng{“Good morning, teacher!”} : on dit \eng{Sir}, \eng{Madam} ou \eng{Mr/Mrs} + nom.
\item Trois étapes : saluer → introduire le sujet → poser une question ouverte.
\item Réagir pour montrer son intérêt : \eng{Really? That's interesting! I see.}
\end{itemize}
\end{aretenir}

\coursec{Demander des informations sur une fête}

Dans la section précédente, nous avons appris à \emph{engager} la conversation en adaptant notre salutation à l'interlocuteur : un camarade de classe (\eng{Hi, how are you?}), un voisin (\eng{Good morning, Mr Ndi.}), un adulte ou une autorité scolaire (\eng{Good morning, Madam. Excuse me...}). Une fois le contact établi, vient l'étape centrale de tout échange : \cle{poser des questions} pour obtenir des informations. C'est le cœur de cette section. Nous allons distinguer deux grands types de questions (\emph{fermées} et \emph{ouvertes}), apprendre à les formuler poliment selon l'interlocuteur, demander des précisions, faire répéter et vérifier que l'on a bien compris.

\courssub{Observer : un dialogue à l'école}

Lisons d'abord ce court dialogue entre une élève et le proviseur (\eng{principal}) de son lycée, la veille de la fête de la Jeunesse (\eng{Youth Day}).

\begin{textebox}[Dialogue : Asking about the Youth Day celebration]
\textbf{Student:} Excuse me, Madam. Could you tell me where the Youth Day march-past will take place?

\textbf{Principal:} At the municipal stadium, at 8 o'clock sharp.

\textbf{Student:} Thank you, Madam. And may I ask who will lead the parade?

\textbf{Principal:} The head boy and the head girl will lead it. All classes must be there by 7.30.

\textbf{Student:} I see. Could you also tell me how long the ceremony will last?

\textbf{Principal:} About three hours. There will be speeches, the march-past and traditional dances.

\textbf{Student:} Sorry, what do you mean by “march-past”?

\textbf{Principal:} It is the moment when students walk in lines in front of the authorities.

\textbf{Student:} So, we must be at the stadium at 7.30, right?

\textbf{Principal:} Exactly. Don't be late!

\textbf{Student:} Thank you very much, Madam. Have a nice day.
\end{textebox}

\textbf{Glossaire :}

\begin{center}
\begin{tabularx}{\linewidth}{|X|X|X|X|}\hline
\rowcolor{popblueL} Mot anglais & Nature & Traduction & Prononciation (approx.) \\ \hline
march-past & nom & défilé & « march-past » (le \eng{t} final se prononce) \\ \hline
sharp & adverbe & précises (heure pile) & « charpe » \\ \hline
head boy / head girl & nom & délégué(e) principal(e) des élèves & « hed boï / hed gueul » \\ \hline
to lead & verbe & diriger, mener & « lide » \\ \hline
to last & verbe & durer & « laste » \\ \hline
speech & nom & discours & « spitch » \\ \hline
Don't be late! & expression & Ne sois pas en retard ! & « donne bi leïte » \\ \hline
\end{tabularx}
\end{center}

\textbf{Questions de compréhension :}
\begin{enumerate}
\item Where will the march-past take place?
\item Who will lead the parade?
\item How long will the ceremony last?
\end{enumerate}

\textbf{Observation.} Remarquez comment l'élève formule ses questions au proviseur : elle ne dit pas \eng{Where will the march-past take place?} mais \eng{Could you tell me where the march-past will take place?}. Avec une autorité, l'anglais préfère la \cle{question indirecte}, plus polie. Observons maintenant chaque type de question en détail.

\courssub{Les questions fermées (yes/no questions)}

Une \cle{question fermée} est une question à laquelle on répond par \eng{yes} ou \eng{no}.

\begin{propriete}[Formation de la question fermée]
La question fermée se construit par \textbf{inversion} : \textbf{auxiliaire + sujet + verbe + complément ?}
\begin{itemize}
\item Avec \eng{be} : \eng{Is the festival next week?} (La fête est-elle la semaine prochaine ?)
\item Avec un verbe ordinaire : l'auxiliaire \eng{do/does/did} porte l'interrogation : \eng{Do people wear traditional clothes?} (Les gens portent-ils des habits traditionnels ?) ; \eng{Does the parade start early?} ; \eng{Did you enjoy the Ngondo festival?}
\item Avec un auxiliaire déjà présent (\eng{will, can, must}) : \eng{Will the mayor attend the ceremony?} (Le maire assistera-t-il à la cérémonie ?)
\end{itemize}
\end{propriete}

\begin{exemplebox}[Questions fermées en contexte]
\eng{Is the Ngondo festival held in Douala?} « La fête du Ngondo a-t-elle lieu à Douala ? »

\eng{Do your neighbours organise a party for Christmas?} « Tes voisins organisent-ils une fête à Noël ? »

\eng{Does the celebration begin at night?} « La célébration commence-t-elle la nuit ? »

\eng{Can students take photos during the ceremony?} « Les élèves peuvent-ils prendre des photos pendant la cérémonie ? »
\end{exemplebox}

\courssub{Les questions ouvertes en Wh-}

Une \cle{question ouverte} commence par un mot interrogatif (\eng{Wh- word}) et appelle une réponse précise, pas un simple \eng{yes} ou \eng{no}.

\begin{propriete}[Formation de la question ouverte]
\textbf{Mot interrogatif + auxiliaire + sujet + verbe + complément ?}

\eng{When does the celebration take place?} ; \eng{Where is the parade?} ; \eng{What time does it start?}

\textbf{Exception importante :} quand le mot interrogatif est le \textbf{sujet} du verbe (\eng{who}, \eng{what} + verbe), il n'y a \textbf{pas d'auxiliaire} : \eng{Who organises the event?} (et non *\eng{Who does organise...?}). Comparez : \eng{Who did you meet?} (\eng{you} est le sujet, donc auxiliaire) / \eng{Who met you?} (\eng{who} est le sujet, pas d'auxiliaire).
\end{propriete}

\begin{center}
\begin{tabularx}{\linewidth}{|X|X|X|}\hline
\rowcolor{popblueL} Mot interrogatif & Question sur une fête & Traduction \\ \hline
When...? & When does the festival begin? & Quand la fête commence-t-elle ? \\ \hline
Where...? & Where is the parade? & Où est le défilé ? \\ \hline
Who...? & Who organises the event? & Qui organise l'événement ? \\ \hline
What...? & What happens during the ceremony? & Que se passe-t-il pendant la cérémonie ? \\ \hline
What time...? & What time does it start? & À quelle heure cela commence-t-il ? \\ \hline
Why...? & Why do people celebrate it? & Pourquoi les gens la célèbrent-ils ? \\ \hline
How...? & How do villagers prepare the feast? & Comment les villageois préparent-ils le festin ? \\ \hline
How long...? & How long does it last? & Combien de temps cela dure-t-il ? \\ \hline
How many...? & How many people attend? & Combien de personnes y assistent ? \\ \hline
\end{tabularx}
\end{center}

\begin{attention}[L'auxiliaire do/does : l'erreur n° 1 des francophones]
En français, on peut poser une question avec la seule intonation : « La fête commence quand ? ». En anglais, c'est \textbf{impossible} : l'auxiliaire est \textbf{obligatoire} dans une question directe avec un verbe ordinaire.

Ne dites jamais : *\eng{When the festival starts?} Dites : \eng{When does the festival start?}

Ne dites jamais : *\eng{Where they dance?} Dites : \eng{Where do they dance?}

Attention aussi à la 3\textsuperscript{e} personne : avec \eng{does}, le verbe principal perd son \textbf{-s} : \eng{When does the festival start?} (et non *\eng{does... starts}).
\end{attention}

\courssub{Les questions polies indirectes : avec les adultes et les autorités}

Avec un camarade, la question directe convient : \eng{When does it start?}. Mais avec un adulte, un voisin âgé, un enseignant ou le proviseur, l'anglais exige une forme plus déférente : la \cle{question indirecte}, introduite par une formule de politesse.

\begin{propriete}[Règle de la question indirecte : l'ordre des mots change]
Après \eng{Could you tell me...}, \eng{I'd like to know...}, \eng{May I ask...}, on reprend l'ordre \textbf{déclaratif} : \textbf{sujet + verbe}, et l'auxiliaire \eng{do/does/did} \textbf{disparaît}.

\eng{Could you tell me when the festival \textbf{starts}?} (et non *\eng{Could you tell me when does the festival start?})

\eng{I'd like to know where the parade \textbf{is}.} (et non *\eng{where is the parade})

Pour une question fermée (yes/no), on utilise \eng{if} ou \eng{whether} : \eng{I'd like to know \textbf{if} the festival takes place in December.} (J'aimerais savoir si la fête a lieu en décembre.)
\end{propriete}

\begin{methode}[Construire une question polie en quatre étapes]
\begin{enumerate}
\item Choisissez la formule d'introduction adaptée à l'interlocuteur : \eng{Could you tell me...} / \eng{I'd like to know...} / \eng{May I ask...}
\item Ajoutez le mot interrogatif (\eng{when, where, who, what time...}) ou \eng{if} pour une question fermée.
\item Remettez les mots dans l'ordre \textbf{sujet + verbe} (comme dans une phrase affirmative).
\item Supprimez l'auxiliaire \eng{do/does} et redonnez au verbe sa marque éventuelle : \eng{does it start} devient \eng{it starts}.
\end{enumerate}
\textbf{Exemple :} \eng{Where does the chief live?} devient \eng{Could you tell me where the chief lives?}
\end{methode}

La figure suivante résume visuellement cette transformation fondamentale.

\begin{popfigure}
\begin{tikzpicture}
\node[draw=popblue, fill=popblueL, rounded corners, text width=5.6cm, align=center, font=\small] (direct) at (0,0) {\textbf{Direct questions}\\[2pt] When does it start?\\ Where is the parade?\\ Do people wear masks?};
\node[draw=poppurple, fill=poppurpleL, rounded corners, text width=6.4cm, align=center, font=\small] (indirect) at (9.5,0) {\textbf{Polite indirect questions}\\[2pt] Could you tell me when it starts?\\ Could you tell me where the parade is?\\ I'd like to know if people wear masks.};
\draw[-{Stealth[length=3mm]}, very thick, poporange] (direct.east) -- (indirect.west) node[midway, above, font=\small\itshape, text=popdark, align=center]{ordre sujet + verbe\\ sans do/does};
\end{tikzpicture}
\legende{De la question directe à la question polie indirecte : l'ordre des mots redevient déclaratif.}
\end{popfigure}

\begin{center}
\begin{tabularx}{\linewidth}{|X|X|}\hline
\rowcolor{popblueL} Question directe (entre camarades) & Question polie indirecte (adulte, autorité) \\ \hline
When does the festival start? & Could you tell me when the festival starts? \\ \hline
Where is the parade? & Could you tell me where the parade is? \\ \hline
Who organises the event? & May I ask who organises the event? \\ \hline
What time does the ceremony begin? & I'd like to know what time the ceremony begins. \\ \hline
Do children receive gifts? & I'd like to know if children receive gifts. \\ \hline
Is the festival open to visitors? & Could you tell me if the festival is open to visitors? \\ \hline
\end{tabularx}
\end{center}

\courssub{Prononciation, accentuation et intonation}

Maîtriser la grammaire ne suffit pas : l'intonation porte le sens de la question et signale la politesse.

\begin{propriete}[Intonation des questions]
\begin{itemize}
\item \textbf{Questions fermées (Yes/No) :} intonation \textbf{montante} à la fin. \eng{Does the festival start at 8?} $\nearrow$
\item \textbf{Questions ouvertes (Wh-) :} intonation \textbf{descendante} à la fin. \eng{When does the festival start?} $\searrow$
\item \textbf{Questions indirectes polies :} l'intonation suit la phrase d'introduction (souvent montante sur \eng{me}, puis descendante sur la proposition). \eng{Could you tell me when it starts?} $\nearrow\searrow$
\item \textbf{Question tags de vérification (\eng{right?}, \eng{isn't it?}) :} intonation montante si on attend une confirmation réelle, descendante si on est quasi certain.
\end{itemize}
\end{propriete}

\begin{exemplebox}[Accentuation de phrase (Sentence stress)]
Dans une question, les \textbf{mots porteurs de sens} (noms, verbes, adjectifs, mots interrogatifs) sont accentués ; les mots grammaticaux (auxiliaires, pronoms, articles) sont faibles.
\eng{\textbf{WHERE} does the \textbf{PARADE} \textbf{TAKE PLACE}?}
\eng{\textbf{COULD} you \textbf{TELL} me \textbf{WHERE} it \textbf{STARTS}?}
\end{exemplebox}

\courssub{Demander des précisions}

Après une première réponse, on a souvent besoin de détails supplémentaires. Voici les outils pour \cle{demander des précisions} :

\begin{exemplebox}[Asking for details]
\eng{What exactly?} « Quoi exactement ? »

\eng{When exactly does the procession leave the village square?} « Quand exactement la procession quitte-t-elle la place du village ? »

\eng{How long does the celebration last?} « Combien de temps dure la célébration ? »

\eng{How many people attend the festival every year?} « Combien de personnes assistent à la fête chaque année ? »

\eng{How often does the community organise this event?} « À quelle fréquence la communauté organise-t-elle cet événement ? »

\eng{Who takes part in the Ngondo festival?} « Qui participe à la fête du Ngondo ? »
\end{exemplebox}

Remarquez la famille des mots en \eng{How + adjectif/adverbe} : \eng{how long} (durée), \eng{how many} (nombre dénombrable), \eng{how much} (quantité indénombrable ou prix), \eng{how often} (fréquence), \eng{how far} (distance), \eng{how old} (âge). C'est un outil très productif pour interroger.

\courssub{Demander de répéter ou d'expliquer}

Dans une conversation réelle, on ne comprend pas toujours du premier coup : bruit du marché, mot inconnu, débit rapide. L'anglais dispose de formules précises, et il est \textbf{poli} de les utiliser plutôt que de rester silencieux.

\begin{exemplebox}[Asking for repetition or explanation]
\eng{Sorry, could you repeat that, please?} « Pardon, pourriez-vous répéter, s'il vous plaît ? »

\eng{Pardon?} « Pardon ? » (avec l'intonation montante)

\eng{Sorry, I didn't catch the time.} « Désolé, je n'ai pas saisi l'heure. »

\eng{What do you mean by “procession”?} « Qu'entendez-vous par “procession” ? »

\eng{Could you speak more slowly, please?} « Pourriez-vous parler plus lentement, s'il vous plaît ? »

\eng{How do you spell “Ngondo”?} « Comment épelle-t-on “Ngondo” ? »
\end{exemplebox}

\begin{attention}[« Pardon ? » : attention à l'intonation]
\eng{Pardon?} se dit avec une intonation \textbf{montante} pour signifier « je n'ai pas compris, pouvez-vous répéter ? ». Avec une intonation descendante, \eng{Pardon.} signifie « je vous demande pardon » (excuse). Évitez aussi *\eng{What?} seul avec un adulte : c'est brutal en anglais. Préférez \eng{Sorry?} ou \eng{Pardon?}.
\end{attention}

\courssub{Vérifier qu'on a compris}

Un bon locuteur \cle{reformule} l'information pour confirmer. C'est une stratégie essentielle de l'interaction orale : elle évite les malentendus (une heure, un lieu, une date).

\begin{exemplebox}[Checking understanding]
\eng{So, the parade starts at 8 a.m., right?} « Donc, le défilé commence à 8 heures, c'est bien cela ? »

\eng{You mean on Saturday morning?} « Vous voulez dire samedi matin ? »

\eng{Let me check: we meet at the stadium at 7.30.} « Je vérifie : on se retrouve au stade à 7 h 30. »

\eng{Did you say 8 o'clock or 9 o'clock?} « Vous avez dit 8 heures ou 9 heures ? »

\eng{So, if I understand correctly, the festival lasts three days.} « Donc, si je comprends bien, la fête dure trois jours. »
\end{exemplebox}

La petite particule finale \eng{right?} (ou \eng{isn't it?}, \eng{doesn't it?}) transforme une affirmation en demande de confirmation : c'est la \emph{question tag}, très fréquente dans la conversation.

\begin{savaistu}[What time...? ou At what time...?]
En anglais de la conversation quotidienne, on dit presque toujours \eng{What time does it start?}. La forme \eng{At what time does it start?} est grammaticalement correcte mais très formelle : on la rencontre surtout à l'écrit (lettres officielles, formulaires). Règle pratique : à l'oral, dites \eng{What time...?} ; à l'écrit formel, \eng{At what time...?} est possible. De même, \eng{Where does it take place?} est plus naturel que \eng{At what place...?}.
\end{savaistu}

\courssub{Dialogues modèles complets : registre familier vs registre formel}

Pour bien fixer les structures, comparons deux échanges complets sur le même thème (la fête du village), l'un entre camarades, l'autre avec une autorité.

\begin{textebox}[Dialogue 1 : Entre camarades (informel)]
\textbf{A:} Hey, \textbf{when's} the village feast this year?
\textbf{B:} \textbf{Next Saturday}. Starts at noon.
\textbf{A:} \textbf{Who's cooking}?
\textbf{B:} The women's association. \textbf{Do you know if} there's a dance after?
\textbf{A:} \textbf{Yeah}, usual band. \textbf{How long} does it go on?
\textbf{B:} Till late. \textbf{Wanna come}?
\textbf{A:} Sure. \textbf{See you there}.
\end{textebox}

\begin{textebox}[Dialogue 2 : Avec le chef de quartier / autorité (formel)]
\textbf{Student:} Good morning, Chief. \textbf{Excuse me, could you tell me when the village feast takes place this year?}
\textbf{Chief:} It will be next Saturday, starting at noon.
\textbf{Student:} \textbf{I'd like to know who organises the cooking.}
\textbf{Chief:} The women's association is in charge.
\textbf{Student:} \textbf{May I ask if there will be traditional dances afterwards?}
\textbf{Chief:} Yes, the cultural troupe will perform until the evening.
\textbf{Student:} \textbf{So, the feast begins at noon and ends in the evening, right?}
\textbf{Chief:} That is correct. You are welcome to attend.
\textbf{Student:} \textbf{Thank you very much, Chief. Have a good day.}
\end{textebox}

\courssub{Entraînons-nous : exercice résolu}

\begin{exoresolu}[Transformer des questions directes en questions polies]
\textbf{Énoncé.} Vous parlez au maire de votre commune. Transformez chaque question directe en question polie indirecte avec \eng{Could you tell me...} ou \eng{I'd like to know...}.

1. When does the cultural festival begin?
2. Where do the dancers perform?
3. Does the festival take place every year?
4. Who pays for the musicians?

\tcblower
\textbf{Solution détaillée.}

1. \eng{Could you tell me when the cultural festival \textbf{begins}?} — On supprime \eng{does} et le verbe reprend son -s de 3\textsuperscript{e} personne (\eng{begins}).

2. \eng{Could you tell me where the dancers \textbf{perform}?} — Ordre sujet + verbe ; \eng{do} disparaît ; \eng{perform} sans -s car \eng{dancers} est pluriel.

3. \eng{I'd like to know \textbf{if} the festival \textbf{takes} place every year.} — Question fermée : on introduit \eng{if} et on reprend l'ordre déclaratif ; \eng{does} disparaît, d'où \eng{takes}.

4. \eng{Could you tell me who \textbf{pays} for the musicians?} — Ici \eng{who} est le sujet du verbe : la question directe n'avait déjà pas d'auxiliaire, la forme ne change donc pas ; on ajoute seulement l'introduction polie.
\end{exoresolu}

\begin{exercice}[À vous de jouer]
Vous interviewez un voisin à Bamenda sur la fête annuelle du village. Formulez :
\begin{enumerate}
\item une question fermée avec \eng{does} sur la date de la fête ;
\item une question ouverte avec \eng{how long} sur la durée de la cérémonie ;
\item une question polie indirecte avec \eng{I'd like to know...} sur le lieu du festin ;
\item une demande de répétition parce que vous n'avez pas entendu l'heure ;
\item une vérification de compréhension avec \eng{So, ... right?}
\end{enumerate}
\end{exercice}

\begin{corrige}[Exercice : À vous de jouer]
Propositions de réponses (d'autres formulations correctes sont possibles) :

1. \eng{Does the festival take place in December?}

2. \eng{How long does the ceremony last?}

3. \eng{I'd like to know where the feast takes place.} (ordre sujet + verbe, sans \eng{does})

4. \eng{Sorry, could you repeat the time, please?}

5. \eng{So, the feast starts at noon, right?}
\end{corrige}

\begin{aretenir}[L'essentiel de la section]
\begin{itemize}
\item Question fermée : \textbf{auxiliaire + sujet + verbe} : \eng{Do people wear traditional clothes?}
\item Question ouverte : \textbf{Wh- + auxiliaire + sujet + verbe} : \eng{When does the celebration take place?} — sauf si le mot interrogatif est le sujet : \eng{Who organises the event?}
\item Question polie indirecte (adultes, autorités) : \eng{Could you tell me / I'd like to know} + mot interrogatif + \textbf{sujet + verbe}, sans \eng{do/does} : \eng{Could you tell me when the festival starts?}
\item Précisions : \eng{How long...? How many...? What exactly...?}
\item Répétition : \eng{Sorry, could you repeat that, please?} ; \eng{What do you mean by...?}
\item Vérification : \eng{So, the parade starts at 8 a.m., right?} ; \eng{You mean on Saturday morning?}
\item Intonation : montante (fermée), descendante (ouverte), montante-descendante (indirecte polie).
\end{itemize}
\end{aretenir}

Maintenant que nous savons interroger efficacement et poliment, il reste à apprendre l'autre moitié de l'échange : \emph{répondre} et \emph{décrire} une célébration. C'est l'objet de la section suivante.

\coursec{Donner des informations et décrire une célébration}

Dans la section précédente, nous avons appris à \cle{demander} des informations sur une fête : \eng{When does it take place? Where is it held? What happens during the festival?} Il est maintenant temps de passer de l'autre côté du dialogue : c'est à vous de \cle{répondre}, de \cle{donner} des informations et de \cle{décrire} une célébration de votre communauté. C'est le cœur de cette leçon : savoir présenter une fête camerounaise — le Ngondo, la fête de la jeunesse, un festival de village — à un camarade, à un voisin ou à un visiteur anglophone.

\courssub{Observer d'abord : un paragraphe modèle}

Avant toute règle, lisons attentivement ce paragraphe. Un élève de Première décrit le festival Ngondo à un correspondant britannique. Repérez comment il donne la date, le lieu, le programme et son opinion.

\begin{textebox}[The Ngondo festival — a model paragraph]
The Ngondo festival takes place every December in Douala. It is held on the banks of the Wouri river, and it usually starts at 9 a.m. sharp. First, the Sawa people gather on the river in their beautiful traditional outfits. Then, there are exciting canoe races and traditional dances. After that, the chiefs give speeches and bless the community. Finally, everybody shares a big meal near the water. Thousands of people come every year: families, neighbours, tourists and even school groups. In my opinion, it is the most exciting event of the year in Douala. I really enjoy the drumming and the songs. Why don't you come with us next December? You should see it! You are welcome to join us.
\end{textebox}

\textbf{Glossaire}

\begin{center}
\begin{tabularx}{\linewidth}{|X|X|X|}\hline
\rowcolor{popblueL} English & Nature & Français \\ \hline
to take place & verbe & avoir lieu \\ \hline
the banks of the river & nom & les berges du fleuve \\ \hline
sharp & adverbe & précises (heure pile) \\ \hline
to gather & verbe & se rassembler \\ \hline
a canoe race & nom & une course de pirogues \\ \hline
to bless & verbe & bénir \\ \hline
to share a meal & expression & partager un repas \\ \hline
drumming & nom & les tambours, le jeu de tambour \\ \hline
\end{tabularx}
\end{center}

\textbf{Questions d'observation} (répondez mentalement) :
\begin{enumerate}
\item When does the festival take place?
\item What happens \emph{first}? What happens \emph{finally}?
\item Which words show the order of events?
\item What is the speaker's opinion?
\end{enumerate}

Vous avez remarqué ? Le paragraphe suit un plan très simple : \cle{date et lieu} $\rightarrow$ \cle{programme en étapes} $\rightarrow$ \cle{participants} $\rightarrow$ \cle{opinion} $\rightarrow$ \cle{invitation}. C'est exactement le plan que vous allez reproduire à l'oral.

\courssub{Donner la date, le lieu et l'heure}

\begin{propriete}[Date, lieu, heure : les trois informations de base]
Pour situer une fête, on utilise trois structures fixes :
\begin{itemize}
\item \textbf{Date} : \eng{It takes place on 11th February.} — préposition \cle{on} + date, \textbf{sans article}.
\item \textbf{Lieu} : \eng{It is held at the community hall.} / \eng{It is held in Douala.} — \cle{at} + lieu précis, \cle{in} + ville ou quartier.
\item \textbf{Heure} : \eng{It starts at 9 a.m. sharp.} / \eng{It begins at 8 o'clock.} — préposition \cle{at} + heure.
\end{itemize}
Le verbe \eng{to be held} (« se tenir ») est un passif très fréquent pour les événements : \eng{The festival is held every year in Buea.} Attention, \eng{held} est le participe passé irrégulier de \eng{hold} (\eng{hold – held – held}).
\end{propriete}

\begin{exemplebox}[Exemples traduits]
\eng{The Youth Day parade takes place on 11th February.} « Le défilé de la fête de la jeunesse a lieu le 11 février. »

\eng{It is held at the community hall.} « Il se tient à la salle des fêtes communale. »

\eng{The ceremony starts at 9 a.m. sharp.} « La cérémonie commence à 9 heures précises. »

\eng{The match begins at 4 p.m. at the stadium.} « Le match commence à 16 heures au stade. »
\end{exemplebox}

\begin{attention}[Erreur fréquente des francophones]
Ne dites jamais \eng{*It takes place the 11th February.} En français on dit « le 11 février », mais en anglais il faut la préposition \cle{on} et \textbf{pas d'article} : \eng{It takes place on 11th February.} Retenez la logique des prépositions de temps : \cle{on} + jour ou date (\eng{on Monday}, \eng{on 20th May}), \cle{in} + mois ou année (\eng{in December}, \eng{in 2025}), \cle{at} + heure (\eng{at 9 a.m.}). Autre piège : \eng{sharp} se place \emph{après} l'heure et signifie « précises, pile » : \eng{at 9 a.m. sharp} « à 9 heures pile ».
\end{attention}

\courssub{Décrire le programme : les connecteurs de séquence}

Pour raconter le déroulement d'une fête, on ordonne les événements avec des \cle{connecteurs de séquence}.

\begin{definition}[Connecteurs de séquence]
Les connecteurs de séquence sont des mots qui ordonnent le récit d'un programme : \cle{first} (d'abord), \cle{then} (puis, ensuite), \cle{after that} (après cela), \cle{next} (ensuite), \cle{finally} (finalement, pour finir). On les place en général en début de phrase, suivis d'une virgule : \eng{First, there is a parade. Then, traditional groups perform dances.}
\end{definition}

\begin{popfigure}
\begin{center}
\begin{tikzpicture}[
box/.style={draw=popblue, fill=popblueL, rounded corners=3pt, align=center, minimum width=2.6cm, minimum height=1.3cm, font=\small},
lbl/.style={font=\small\bfseries, text=popdark},
arr/.style={-{Stealth[length=3mm]}, thick, poporange}]
\node[box, align=center] (a) at (0,0){First\\ parade};
\node[box, align=center] (b) at (3.6,0){Then\\ speeches};
\node[box, align=center] (c) at (7.2,0){After that\\ dances};
\node[box, align=center] (d) at (10.8,0){Finally\\ meal};
\draw[arr] (a) -- (b);
\draw[arr] (b) -- (c);
\draw[arr] (c) -- (d);
\node[lbl] at (1.8,-0.75) {then};
\node[lbl] at (5.4,-0.75) {after that};
\node[lbl] at (9,-0.75) {finally};
\end{tikzpicture}
\end{center}
\legende{Frise chronologique d'un programme de fête : les connecteurs de séquence ordonnent le récit.}
\end{popfigure}

\begin{exemplebox}[Décrire le programme — exemples traduits]
\eng{First, there is a parade through the streets of Bamenda.} « D'abord, il y a un défilé dans les rues de Bamenda. »

\eng{Then, traditional groups perform dances.} « Ensuite, des groupes traditionnels exécutent des danses. »

\eng{After that, the mayor gives a speech.} « Après cela, le maire prononce un discours. »

\eng{Finally, we share a meal with our neighbours.} « Pour finir, nous partageons un repas avec nos voisins. »
\end{exemplebox}

\courssub{Décrire les participants et les activités}

Pour décrire qui fait quoi, on utilise le \cle{présent simple} : c'est le temps des habitudes et des événements qui se répètent chaque année.

\begin{exemplebox}[Participants et activités — exemples traduits]
\eng{People wear traditional outfits.} « Les gens portent des tenues traditionnelles. »

\eng{The chief gives a speech and blesses the community.} « Le chef prononce un discours et bénit la communauté. »

\eng{Children sing and dance in the school yard.} « Les enfants chantent et dansent dans la cour de l'école. »

\eng{Women prepare traditional dishes like eru and ndolé.} « Les femmes préparent des plats traditionnels comme l'eru et le ndolé. »

\eng{The school band plays music all afternoon.} « La fanfare de l'école joue de la musique tout l'après-midi. »
\end{exemplebox}

\begin{attention}[La troisième personne du singulier]
Au présent simple, n'oubliez pas le \cle{-s} à la troisième personne du singulier : \eng{The chief give\textbf{s} a speech.} / \eng{The festival start\textbf{s} at 9 a.m.} Les francophones l'oublient souvent à l'oral parce que le français ne marque pas cette différence à l'oreille. Attention aussi aux verbes en \emph{-ss, -sh, -ch, -o} qui prennent \cle{-es} : \eng{she watch\textbf{es}}, \eng{he go\textbf{es}}, \eng{it bless\textbf{es}}. Entraînez-vous : \eng{he wear\textbf{s}}, \eng{she dance\textbf{s}}, \eng{it happen\textbf{s}}.
\end{attention}

\courssub{Exprimer la fréquence, la tradition et l'opinion}

\begin{propriete}[Fréquence et tradition]
Pour dire qu'un événement est régulier ou ancien :
\begin{itemize}
\item \eng{It happens every year.} « Cela a lieu chaque année. »
\item \eng{It takes place every December.} « Cela a lieu chaque mois de décembre. »
\item \eng{It is an old tradition in our village.} « C'est une vieille tradition dans notre village. »
\item \eng{We have celebrated it for generations.} « Nous le célébrons depuis des générations. »
\end{itemize}
Pour donner son opinion, on utilise : \eng{In my opinion, ...} (« à mon avis »), \eng{I think (that) ...} (« je pense que »), \eng{I really enjoy ...} + nom ou verbe en \cle{-ing} : \eng{I really enjoy the drumming.}
\end{propriete}

\begin{exemplebox}[Opinion — exemples traduits]
\eng{In my opinion, it's the most exciting event of the year.} « À mon avis, c'est l'événement le plus passionnant de l'année. »

\eng{I think the dances are wonderful.} « Je trouve que les danses sont magnifiques. »

\eng{I really enjoy the drumming.} « J'adore vraiment les tambours. »
\end{exemplebox}

\begin{attention}[\eng{enjoy} + verbe en -ing]
Après \eng{enjoy}, le verbe qui suit est toujours en \cle{-ing} : \eng{I enjoy danc\textbf{ing}.} et jamais \eng{*I enjoy to dance.} Piège de traduction : « j'aime danser » se dit \eng{I like dancing} ou \eng{I like to dance}, mais avec \eng{enjoy}, seule la forme en -ing est correcte. Autre faux ami : \eng{exciting} signifie « passionnant » (la fête est passionnante), tandis que \eng{excited} décrit la personne (« je suis impatient, enthousiaste ») : \eng{The festival is exciting. I am excited about it.}
\end{attention}

\courssub{Inviter l'interlocuteur}

Terminer la description par une invitation rend le dialogue naturel et chaleureux.

\begin{center}
\begin{tabularx}{\linewidth}{|X|X|X|}\hline
\rowcolor{popblueL} English & Français & Remarque \\ \hline
Why don't you come with us? & Pourquoi ne viens-tu pas avec nous ? & suggestion amicale \\ \hline
You should see it! & Tu devrais voir ça ! & conseil enthousiaste \\ \hline
You're welcome to join us. & Tu es le bienvenu pour te joindre à nous. & formule polie, tout public \\ \hline
Would you like to come? & Voudrais-tu venir ? & invitation directe et polie \\ \hline
\end{tabularx}
\end{center}

Avec un adulte ou une autorité scolaire, préférez les formules plus polies : \eng{You're welcome to join us, Sir.} / \eng{Would you like to attend the ceremony, Madam?} Le verbe \eng{to attend} (« assister à ») est plus soutenu que \eng{to come}.

\courssub{Méthode : décrire une fête en quatre points}

\begin{methode}[Comment décrire une célébration à l'oral]
\begin{enumerate}
\item \textbf{Nom et date} : annoncez la fête et quand elle a lieu. \eng{The Ngondo festival takes place every December.}
\item \textbf{Lieu et participants} : dites où elle se tient et qui y participe. \eng{It is held on the Wouri river. Thousands of people come.}
\item \textbf{Déroulement} : racontez le programme avec les connecteurs \eng{first, then, after that, finally}. \eng{First, people gather. Then, there are canoe races. Finally, we share a meal.}
\item \textbf{Opinion et invitation} : terminez par votre avis et invitez l'interlocuteur. \eng{In my opinion, it's wonderful. Why don't you come with us?}
\end{enumerate}
\end{methode}

\begin{exoresolu}[Décrire la fête de la jeunesse]
Énoncé : Un correspondant américain vous demande : \eng{What happens in your town on 11th February?} Décrivez la fête de la jeunesse en quatre ou cinq phrases, en suivant la méthode des quatre points.
\tcblower
Solution détaillée :
\eng{Youth Day takes place on 11th February in Yaoundé. It is held at the Ahmadou Ahidjo stadium, and it starts at 8 a.m. First, pupils and students parade in their school uniforms. Then, the authorities give speeches. After that, there are songs and traditional dances. Finally, we go home and share a meal with our families. In my opinion, it is a great day because the whole community is together. You should see it!}

Traduction : « La fête de la jeunesse a lieu le 11 février à Yaoundé. Elle se tient au stade Ahmadou Ahidjo et commence à 8 heures. D'abord, les élèves et les étudiants défilent en uniforme. Ensuite, les autorités prononcent des discours. Après cela, il y a des chants et des danses traditionnelles. Pour finir, nous rentrons et partageons un repas en famille. À mon avis, c'est une belle journée parce que toute la communauté est réunie. Tu devrais voir ça ! »

Points de méthode : date avec \eng{on} sans article ; lieu avec \eng{at} ; connecteurs \eng{first, then, after that, finally} ; \eng{-s} à la troisième personne (\eng{takes}, \eng{starts}, \eng{give}) ; opinion avec \eng{in my opinion}.
\end{exoresolu}

\begin{savaistu}[Le carnaval de Notting Hill]
À Londres, le carnaval de Notting Hill, créé dans les années 1960 par la communauté caribéenne, attire plus d'un million de visiteurs chaque mois d'août. Défilés costumés, orchestres de steel drums et plats antillais envahissent les rues pendant deux jours. Comme le Ngondo à Douala, c'est une fête communautaire devenue un événement national : une belle preuve que les célébrations locales peuvent rayonner dans le monde entier.
\end{savaistu}

\begin{exercice}[À vous de parler]
Décrivez à voix haute une fête de votre village ou de votre quartier (fête traditionnelle, tournoi de football, journée culturelle de l'école) en suivant la méthode des quatre points : nom et date, lieu et participants, déroulement avec \eng{first, then, after that, finally}, opinion et invitation. Visez cinq phrases minimum.
\end{exercice}

\coursec{Relancer et conclure la conversation}

Dans la section précédente, nous avons appris à \cle{donner des informations} et à \cle{décrire une célébration} : parler des danses, des plats, des courses de pirogues ou des cérémonies traditionnelles. Mais une conversation n'est pas un monologue ! Un bon locuteur sait aussi \cle{relancer} son interlocuteur pour que l'échange continue, et \cle{conclure} poliment lorsqu'il est temps de partir. C'est précisément l'objet de cette section : apprendre à ne jamais laisser mourir un dialogue, puis à le terminer avec élégance, en adaptant son langage à la personne à qui l'on parle — camarade, voisin, adulte ou autorité scolaire.

\courssub{Observer : la fin d'un dialogue modèle}

Avant toute règle, observons comment deux élèves terminent naturellement leur conversation à propos d'une fête communautaire.

\begin{textebox}[End of a model dialogue — At Government High School, Buea]
A: So the Ngondo festival takes place every December on the banks of the Wouri river?

B: Exactly. The water spirits are honoured, and there are wonderful canoe races.

A: Really? Tell me more about the canoe races!

B: Well, the best paddlers from the Sawa community compete. The atmosphere is incredible — drums, songs, colourful pirogues...

A: That sounds amazing! And you, have you ever attended the festival?

B: Yes, twice. Last year my uncle won the race!

A: Wow! And then what happened?

B: The whole family celebrated until midnight. But tell me — what about your village? Do you have similar celebrations?

A: We have a harvest festival in September, with traditional dances.

B: Interesting! I would love to see that one day... Well, I must go now. Class is about to start. Thank you very much for the information.

A: You're welcome. It was nice talking to you. See you at the celebration!

B: See you later, and take care!
\end{textebox}

\textbf{Glossaire}

\begin{center}
\begin{tabularx}{\linewidth}{|l|l|X|}
\hline
\rowcolor{popblueL} \textbf{Mot anglais} & \textbf{Nature} & \textbf{Traduction} \\
\hline
banks & nom (pluriel) & les rives, les bords (d'une rivière) \\
\hline
canoe race & nom & course de pirogues \\
\hline
paddler & nom & pagayeur \\
\hline
to attend & verbe & assister à, être présent à \\
\hline
harvest & nom & la récolte, la moisson \\
\hline
You're welcome & expression & je vous en prie, de rien \\
\hline
\end{tabularx}
\end{center}

\textbf{Questions d'observation}

\begin{enumerate}
\item Which expressions does speaker A use to show interest? (\emph{Really? Tell me more! That sounds amazing!})
\item Which question does speaker B ask to give the floor back to A? (\emph{What about your village?})
\item How does speaker B announce that he must leave? (\emph{Well, I must go now. Class is about to start.})
\item Which sentences close the conversation politely? (\emph{Thank you very much for the information. / It was nice talking to you. / See you later!})
\end{enumerate}

\courssub{Relancer la conversation : les questions ouvertes}

\begin{definition}[Follow-up question]
Une \cle{follow-up question} (question de relance) est une question qui \textbf{prolonge la conversation à partir de la réponse reçue}. Elle montre à l'interlocuteur qu'on l'a écouté et qu'on veut en savoir davantage. Elle s'oppose à la question fermée, qui n'appelle qu'un « yes » ou un « no » et fait mourir le dialogue.
\end{definition}

\begin{propriete}[Les trois techniques de relance]
Pour relancer un dialogue en anglais, on dispose de trois techniques complémentaires :
\begin{enumerate}
\item \textbf{La question ouverte générale} : elle redonne la parole à l'interlocuteur sur sa propre expérience. \\
\eng{And you, what do you think of the festival?} « Et toi, que penses-tu de la fête ? » \\
\eng{What about your village?} « Et ton village ? » \\
\eng{Do you have similar celebrations?} « Avez-vous des célébrations semblables ? »
\item \textbf{La réaction d'intérêt} : une courte exclamation suivie d'une invitation à continuer. \\
\eng{Really? Tell me more!} « Vraiment ? Raconte-m'en plus ! » \\
\eng{And then what happened?} « Et ensuite, que s'est-il passé ? » \\
\eng{That sounds amazing!} « Ça a l'air extraordinaire ! »
\item \textbf{Le rebond sur une information précise} : on reprend un mot ou une idée de la réponse pour poser une question ciblée. \\
\eng{You mentioned canoe races — when exactly do they start?} « Tu as parlé de courses de pirogues — quand commencent-elles exactement ? »
\end{enumerate}
\end{propriete}

\begin{exemplebox}[Rebondir sur une information]
Le rebond consiste à reprendre un élément de la réponse avec \eng{You said...}, \eng{You mentioned...} ou \eng{You told me...} :

\eng{You said the dances start at sunset — where exactly do they take place?} « Tu as dit que les danses commencent au coucher du soleil — où ont-elles lieu exactement ? »

\eng{You mentioned the chief's speech — what does he usually say?} « Tu as mentionné le discours du chef — que dit-il d'habitude ? »

\eng{You told me about the food — which dish is the most popular?} « Tu m'as parlé de la nourriture — quel plat est le plus populaire ? »
\end{exemplebox}

\begin{methode}[Ne jamais laisser mourir un dialogue : la méthode R.R.R.]
Pour entretenir n'importe quelle conversation en anglais, applique le cycle en trois étapes :
\begin{enumerate}
\item \textbf{Réagir} : montre que tu écoutes avec une exclamation courte — \eng{Really?}, \eng{Wow!}, \eng{That's interesting!}, \eng{No way!}
\item \textbf{Relancer} : pose une question ouverte — \eng{What about you?}, \eng{Tell me more!}, \eng{And then what happened?}
\item \textbf{Rebondir} : reprends un mot de la réponse pour creuser — \eng{You said that the race starts early — at what time exactly?}
\end{enumerate}
Puis recommence le cycle : réagir, relancer, rebondir. Avec cette méthode, ton dialogue ne s'arrêtera jamais de façon gênante.
\end{methode}

\begin{attention}[L'intonation des réactions d'intérêt]
Les petites phrases de relance ne fonctionnent que si l'\textbf{intonation} est juste. \eng{Really?} se dit avec une intonation \textbf{montante} (la voix monte à la fin), comme un « Vraiment ? » enthousiaste en français : « \emph{ri-li} ? ». Si la voix descend, la phrase sonne comme un doute ou une moquerie ! De même, \eng{Tell me more!} se dit avec énergie, l'accent principal sur \emph{more} : « tèl mi \textbf{mo} ! ». Entraîne-toi à voix haute : l'intonation montante est la marque de l'intérêt en anglais.
\end{attention}

\courssub{Signaler qu'on doit terminer}

Une conversation ne s'interrompt pas brutalement. En anglais, on \textbf{annonce d'abord} qu'on doit partir, avec une formule d'excuse ou d'explication, avant de conclure.

\begin{propriete}[Formules pour signaler la fin de la conversation]
\begin{itemize}
\item \eng{Well, I have to go now.} « Bon, je dois y aller maintenant. » (neutre, tous interlocuteurs)
\item \eng{I'm afraid I must leave.} « Je crains de devoir partir. » (plus formel, avec un adulte ou une autorité)
\item \eng{Class is about to start.} « Le cours va bientôt commencer. » (justification pratique, au lycée)
\item \eng{Sorry, I have to run — my mother is waiting for me.} « Désolé, je dois filer — ma mère m'attend. » (familier, entre camarades)
\end{itemize}
Remarque : le mot \eng{Well} placé en tête de phrase sert de « signal » : il prépare l'interlocuteur à un changement. C'est l'équivalent du français « Bon, ... » ou « Écoute, ... ».
\end{propriete}

\courssub{Conclure poliment et prendre congé}

\begin{propriete}[Les deux étapes d'une conclusion réussie]
Une conclusion polie en anglais comporte presque toujours \textbf{deux étapes} :
\begin{enumerate}
\item \textbf{Remercier ou apprécier} : \eng{Thank you very much for the information.} « Merci beaucoup pour ces informations. » / \eng{It was nice talking to you.} « C'était un plaisir de discuter avec vous. »
\item \textbf{Prendre congé} : formule d'au revoir adaptée à l'interlocuteur.
\end{enumerate}
\end{propriete}

\begin{propriete}[Prendre congé selon l'interlocuteur]
Le choix de la formule d'au revoir dépend du \textbf{statut social} de la personne à qui l'on parle :
\begin{itemize}
\item Avec un \textbf{camarade} ou un ami : \eng{See you later!} « À plus tard ! » / \eng{Bye!} / \eng{Take care!} « Prends soin de toi ! » / \eng{See you at the celebration!} « À la fête ! »
\item Avec un \textbf{voisin} ou un adulte : \eng{Goodbye, Mr Atangana. Have a nice day.} « Au revoir, M. Atangana. Bonne journée. »
\item Avec une \textbf{autorité scolaire} (proviseur, censeur, enseignant) : \eng{Goodbye, Madam. Thank you very much, Sir.} « Au revoir, Madame. Merci beaucoup, Monsieur. »
\end{itemize}
\end{propriete}

\begin{center}
\begin{tabularx}{\linewidth}{|X|X|X|}
\hline
\rowcolor{popblueL} \textbf{Fonction} & \textbf{English} & \textbf{Français} \\
\hline
Relancer (question ouverte) & And you, what do you think of the festival? & Et toi, que penses-tu de la fête ? \\
\hline
Relancer (redonner la parole) & What about your village? & Et ton village ? \\
\hline
Montrer de l'intérêt & Really? Tell me more! & Vraiment ? Raconte-m'en plus ! \\
\hline
Montrer de l'intérêt & That sounds amazing! & Ça a l'air extraordinaire ! \\
\hline
Faire continuer le récit & And then what happened? & Et ensuite, que s'est-il passé ? \\
\hline
Rebondir & You mentioned canoe races — when do they start? & Tu as parlé des courses de pirogues — quand commencent-elles ? \\
\hline
Signaler le départ & Well, I have to go now. & Bon, je dois y aller. \\
\hline
Signaler le départ (formel) & I'm afraid I must leave. & Je crains de devoir partir. \\
\hline
Remercier & Thank you very much for the information. & Merci beaucoup pour ces informations. \\
\hline
Apprécier & It was nice talking to you. & C'était un plaisir de discuter avec vous. \\
\hline
Congé (camarade) & See you later! Take care! & À plus tard ! Prends soin de toi ! \\
\hline
Congé (autorité) & Goodbye, Madam. Have a nice day. & Au revoir, Madame. Bonne journée. \\
\hline
\end{tabularx}
\end{center}

\begin{attention}[Ne traduis pas « au revoir » systématiquement par « Goodbye » !]
Erreur très fréquente chez les francophones : utiliser \eng{Goodbye} dans toutes les situations. Or, entre camarades, \eng{Goodbye} sonne froid, distant, voire définitif (comme « adieu »). On préfère :
\begin{itemize}
\item \eng{See you!} ou \eng{See you later!} — « À plus ! »
\item \eng{Bye!} ou \eng{Bye-bye!} — « Salut ! »
\item \eng{Take care!} — « Prends soin de toi ! »
\end{itemize}
Réserve \eng{Goodbye} aux situations formelles (adultes, autorités). Autre piège : à la fin d'une conversation, la forme la plus naturelle est \eng{It was nice talking to you} (gérondif en \emph{-ing} après \eng{nice}) ; la tournure \eng{It was nice to talk to you} existe, mais elle est moins idiomatique dans ce contexte — retiens le gérondif, tu ne te tromperas jamais.
\end{attention}

\begin{savaistu}[Pourquoi les anglophones « traînent » avant de dire au revoir]
Dans les cultures anglophones, conclure une conversation comporte presque toujours une \textbf{appréciation} avant le congé : \eng{It was lovely talking to you.}, \eng{I really enjoyed our chat.} Sauter cette étape et dire directement \eng{Goodbye} peut sembler brusque, voire impoli — l'interlocuteur peut croire qu'il vous a ennuyé ou offensé. Au Cameroun anglophone comme en Grande-Bretagne, on ajoute souvent une promesse de se revoir : \eng{See you at the celebration!}, \eng{Let's talk again soon!} C'est ce qu'on appelle le « pre-closing » : la conversation se ferme en douceur, par étapes.
\end{savaistu}

\courssub{Vue d'ensemble : le déroulement complet d'un dialogue}

La relance et la conclusion s'inscrivent dans la structure générale d'un dialogue. Voici l'organigramme complet des sept étapes d'un échange réussi.

\begin{popfigure}
\begin{center}
\begin{tikzpicture}[
box/.style={rectangle, rounded corners=3pt, draw=popblue, fill=popblueL, text width=3.1cm, align=center, font=\small, minimum height=0.85cm},
fleche/.style={-{Stealth[length=2.5mm]}, thick, popdark}]
\node[box, align=center] (g) at (0,0){\textbf{1. Greeting}\\ \eng{Hello! Good morning, Sir.}};
\node[box, align=center] (i) at (4.3,0){\textbf{2. Introduce topic}\\ \eng{I'd like to know about the festival.}};
\node[box, align=center] (a1) at (8.6,0){\textbf{3. Ask}\\ \eng{When does it take place?}};
\node[box, align=center] (a2) at (8.6,-2.2){\textbf{4. Answer}\\ \eng{It takes place in December.}};
\node[box, align=center] (f) at (4.3,-2.2){\textbf{5. Follow up}\\ \eng{Really? Tell me more!}};
\node[box, align=center] (c) at (0,-2.2){\textbf{6. Conclude}\\ \eng{Thank you. It was nice talking to you.}};
\node[box, draw=poporange, fill=poporangeL, align=center] (l) at (0,-4.4){\textbf{7. Take leave}\\ \eng{See you later! / Goodbye, Madam.}};
\draw[fleche] (g) -- (i);
\draw[fleche] (i) -- (a1);
\draw[fleche] (a1) -- (a2);
\draw[fleche] (a2) -- (f);
\draw[fleche] (f) -- (c);
\draw[fleche] (c) -- (l);
\draw[fleche, dashed, popmuted] (f.north) to[bend right=25] (a1.west);
\node[font=\footnotesize\itshape, text=popmuted] at (6.4,-1.05) {on peut boucler : ask / answer / follow up};
\end{tikzpicture}
\end{center}
\legende{Les sept étapes d'un dialogue complet : la boucle Ask -- Answer -- Follow up peut se répéter plusieurs fois avant la conclusion.}
\end{popfigure}

Remarque la flèche en pointillés : les étapes 3, 4 et 5 forment une \textbf{boucle}. Tant que la conversation continue, on enchaîne question, réponse et relance. Ce n'est que lorsqu'on doit partir qu'on passe aux étapes 6 et 7.

\courssub{Entraînement guidé}

\begin{exoresolu}[Compléter la fin d'un dialogue]
\textbf{Énoncé.} Complete the end of this dialogue between Nadia and her neighbour Mr Ewusi, using expressions from the lesson: \emph{Really? Tell me more! / You mentioned / Well, I have to go now / It was nice talking to you / Have a nice day}.

Nadia: The whole village dances around the fire until dawn.

Mr Ewusi: (1) \dots\dots\dots\dots

Nadia: Well, the oldest dancer leads the group, and the drums never stop.

Mr Ewusi: (2) \dots\dots\dots\dots the drums — who plays them?

Nadia: Young men from the quarter. But (3) \dots\dots\dots\dots — my bus is leaving.

Mr Ewusi: Of course. (4) \dots\dots\dots\dots, Nadia.

Nadia: Thank you, Mr Ewusi. (5) \dots\dots\dots\dots!

\tcblower
\textbf{Solution détaillée.}
\begin{enumerate}
\item \eng{Really? Tell me more!} — Mr Ewusi réagit à l'information et invite Nadia à continuer (technique de la réaction d'intérêt).
\item \eng{You mentioned} — il rebondit sur un mot précis de la réponse (\emph{the drums}) pour poser une question ciblée.
\item \eng{Well, I have to go now} — Nadia signale poliment qu'elle doit partir, en justifiant (\emph{my bus is leaving}).
\item \eng{It was nice talking to you} — l'appréciation avant le congé, étape indispensable de la conclusion polie.
\item \eng{Have a nice day} — formule de congé adaptée : Nadia s'adresse à un adulte, son voisin ; \eng{See you!} serait trop familier ici.
\end{enumerate}
\end{exoresolu}

\begin{exercice}[À toi de jouer]
\begin{enumerate}
\item Put these sentences in the correct order to build the end of a dialogue: \\
a) \eng{See you at the festival!} \\
b) \eng{Well, I'm afraid I must leave now.} \\
c) \eng{Thank you very much for the information, Sir.} \\
d) \eng{You're welcome. Goodbye!} \\
e) \eng{It was nice talking to you.}
\item Write two follow-up questions to react to this sentence: \eng{In my village, we celebrate the yam festival every August.}
\item You are talking to your headmistress about the school open day. Write the last three lines of the dialogue (signal your departure, thank her, take leave appropriately).
\end{enumerate}
\end{exercice}

\begin{corrige}[Exercice « À toi de jouer »]
\begin{enumerate}
\item Ordre correct : b → c → e → d → a. On signale d'abord le départ (b), puis on remercie (c), on exprime son appréciation (e), l'interlocuteur répond au remerciement (d), et on prend congé (a).
\item Réponses possibles : \eng{Really? Tell me more about the yam festival!} / \eng{You mentioned August — on which day exactly does it start?} / \eng{That sounds interesting! What do people eat during the celebration?}
\item Modèle de réponse : \eng{Well, I'm afraid I must leave now — class is about to start. Thank you very much for the information, Madam. It was nice talking to you. Goodbye, Madam, and have a nice day.}
\end{enumerate}
\end{corrige}

\begin{aretenir}[L'essentiel de la section]
\begin{itemize}
\item Une \cle{follow-up question} prolonge la conversation à partir de la réponse reçue : \eng{And you? What about your village? Tell me more!}
\item Méthode R.R.R. : \textbf{Réagir} (\eng{Really?}) → \textbf{Relancer} (\eng{What about...?}) → \textbf{Rebondir} (\eng{You said that...?}).
\item Les réactions d'intérêt se disent avec une intonation \textbf{montante} : \eng{Really?} « \emph{ri-li} ? ».
\item On annonce son départ avant de conclure : \eng{Well, I have to go now.} / \eng{I'm afraid I must leave.}
\item Une conclusion polie = remerciement + appréciation + congé : \eng{Thank you for the information. It was nice talking to you.}
\item Le congé s'adapte à l'interlocuteur : \eng{See you later!} entre camarades ; \eng{Goodbye, Madam. Have a nice day.} avec une autorité.
\item Piège : \eng{Goodbye} n'est pas la formule universelle ; entre amis, préfère \eng{See you!}, \eng{Bye!}, \eng{Take care!}
\end{itemize}
\end{aretenir}

\coursec{Dialogues modèles complets}

Dans la section précédente, nous avons appris à \cle{relancer} une conversation (\eng{What about you?}, \eng{Really? Tell me more!}) et à la \cle{conclure} poliment (\eng{See you there!}, \eng{Thank you for your help.}). Il est temps d'assembler toutes ces pièces : nous allons maintenant étudier trois \cle{dialogues modèles complets}, du début à la fin, dans trois situations différentes. Vous constaterez que la même information peut être demandée de trois façons très différentes selon l'interlocuteur : un camarade, un voisin ou le proviseur.

\courssub{Dialogue 1 : entre camarades (registre informel)}

\begin{textebox}[Dialogue 1 — Two mates talking about Youth Day]
A: “Hi Bih! Are you going to the Youth Day parade?”\\
B: “Sure! It starts at 8 at the stadium. Why don't we go together?”\\
A: “Great idea! See you there!”\\
B: “OK, don't forget your uniform!”
\end{textebox}

\textbf{Observation.} Lisez ce dialogue à voix haute. Quelles formules reconnaissez-vous ?

\begin{itemize}
\item \cle{Ouverture} : \eng{Hi Bih!} — salutation familière, avec le prénom.
\item \cle{Demande d'information} : \eng{Are you going to the Youth Day parade?} — question directe, courte.
\item \cle{Information donnée} : \eng{It starts at 8 at the stadium.}
\item \cle{Invitation} : \eng{Why don't we go together?} — littéralement « Pourquoi n'irions-nous pas ensemble ? » ; c'est la formule typique de proposition entre amis.
\item \cle{Clôture légère} : \eng{Great idea! See you there!} — « Super idée ! On se retrouve là-bas ! »
\end{itemize}

\begin{propriete}[Le registre informel entre camarades]
Entre amis (\eng{mates}), l'anglais utilise :
\begin{itemize}
\item des salutations courtes : \eng{Hi!}, \eng{Hello!}, \eng{Hey!} ;
\item des questions \textbf{directes} : \eng{Are you going...?}, \eng{When does it start?} ;
\item des contractions : \eng{don't}, \eng{it's}, \eng{let's} ;
\item des exclamations : \eng{Sure!}, \eng{Great!}, \eng{Cool!} ;
\item des invitations avec \eng{Why don't we...?} ou \eng{Let's...}.
\end{itemize}
\end{propriete}

\begin{exemplebox}[Exemples traduits — registre informel]
\eng{Why don't we go together?} « Pourquoi n'irions-nous pas ensemble ? »\\
\eng{Let's meet at the stadium.} « Retrouvons-nous au stade. »\\
\eng{See you there!} « On se voit là-bas ! »\\
\eng{Don't forget your uniform!} « N'oublie pas ta tenue ! »
\end{exemplebox}

\courssub{Dialogue 2 : élève et proviseur (registre formel)}

\begin{textebox}[Dialogue 2 — A student asking the Principal about the 11th February programme]
Student: “Good morning, Sir. Could you tell me what time the march-past begins?”\\
Principal: “Good morning. At 8 o'clock sharp. Don't be late.”\\
Student: “Thank you very much, Sir. And where will the prize-giving ceremony take place?”\\
Principal: “In the school hall, after the march-past.”\\
Student: “I see. Thank you for your help, Sir. Goodbye.”\\
Principal: “Goodbye, and see you on the 11th of February.”
\end{textebox}

\textbf{Observation.} Comparez ce dialogue avec le premier. L'information demandée est presque la même (« à quelle heure ça commence ? »), mais la formulation change totalement.

\begin{itemize}
\item \cle{Ouverture formelle} : \eng{Good morning, Sir.} — jamais \eng{Hi!} avec une autorité scolaire.
\item \cle{Question indirecte} : \eng{Could you tell me what time the march-past begins?} — « Pourriez-vous me dire à quelle heure commence le défilé ? »
\item \cle{Réaction} : \eng{I see.} — « Je vois. »
\item \cle{Remerciements} : \eng{Thank you very much, Sir.} / \eng{Thank you for your help.}
\item \cle{Clôture formelle} : \eng{Goodbye.}
\end{itemize}

\begin{propriete}[La question indirecte — forme et emploi]
Pour s'adresser à un adulte ou à une autorité, on enrobe la question directe dans une formule polie :
\begin{center}
\textbf{Could you tell me / Can you tell me / I'd like to know} + sujet + verbe
\end{center}
Attention à l'\cle{ordre des mots} : après \eng{Could you tell me}, on reprend l'ordre de la phrase affirmative (sujet + verbe), \textbf{sans inversion} :
\begin{itemize}
\item Question directe : \eng{What time does the march-past begin?}
\item Question indirecte : \eng{Could you tell me what time the march-past \textbf{begins}?} (et non \emph{does begin})
\end{itemize}
Autres exemples : \eng{Could you tell me where the ceremony takes place?} « Pourriez-vous me dire où a lieu la cérémonie ? » ; \eng{I'd like to know if the festival is on Saturday.} « J'aimerais savoir si la fête a lieu samedi. » (Pour les questions fermées, on utilise \eng{if} ou \eng{whether}).
\end{propriete}

\begin{attention}[Erreurs fréquentes des francophones]
\begin{itemize}
\item Ne dites jamais \eng{Could you tell me what time does it begin?} — double marque interrogative (inversion + \eng{does}) : faute classique ! Après \eng{Could you tell me}, ordre affirmatif obligatoire.
\item \eng{Don't be late.} signifie « Ne sois pas en retard » (et non « ne sois pas tard »).
\item En anglais, on ne dit pas \eng{Mister the Principal} : on dit simplement \eng{Sir} ou \eng{Mr Atanga}.
\item \eng{Sharp} après une heure signifie « précises, pile » : \eng{at 8 o'clock sharp} = « à 8 heures précises ».
\end{itemize}
\end{attention}

\courssub{Dialogue 3 : entre voisins (registre neutre)}

\begin{textebox}[Dialogue 3 — Two neighbours talking about the neighbourhood festival]
Mrs Ngono: “Good evening, Mr Ekane. I heard there is a festival in our quarter this weekend.”\\
Mr Ekane: “Yes, that's right. The Mvog-Ada festival takes place on Saturday.”\\
Mrs Ngono: “Oh, interesting! What is the programme?”\\
Mr Ekane: “Well, there is traditional dancing in the morning, a football match in the afternoon, and a big meal in the evening.”\\
Mrs Ngono: “Really? That sounds wonderful! What time does it start?”\\
Mr Ekane: “At 9 o'clock, at the community field. You are welcome to come with your family.”\\
Mrs Ngono: “Thank you very much. We will be there!”\\
Mr Ekane: “Fine. See you on Saturday, then.”
\end{textebox}

\textbf{Glossaire.} \emph{quarter} (n.) : quartier ; \emph{to take place} (v.) : avoir lieu ; \emph{traditional dancing} (n.) : danses traditionnelles ; \emph{community field} (n.) : terrain communautaire ; \emph{you are welcome to come} : vous êtes invité(e) à venir ; \emph{then} (adv.) : alors, dans ce cas.

\textbf{Observation.} Ce registre est \cle{neutre} : poli mais détendu. On y trouve :
\begin{itemize}
\item \cle{Présentation du sujet} : \eng{I heard there is a festival in our quarter this weekend.} — « J'ai entendu dire qu'il y a une fête dans notre quartier ce week-end. »
\item \cle{Demande polie mais directe} : \eng{What is the programme?}, \eng{What time does it start?}
\item \cle{Description du programme} : \eng{there is traditional dancing..., a football match..., and a big meal...}
\item \cle{Relance} : \eng{Really? That sounds wonderful!} — « Vraiment ? Ça a l'air merveilleux ! »
\item \cle{Invitation} : \eng{You are welcome to come with your family.}
\item \cle{Clôture} : \eng{See you on Saturday, then.}
\end{itemize}

\begin{definition}[Vocabulaire clé : fêtes et célébrations]
\begin{itemize}
\item \eng{festival} (n.) /« fes-ti-val/ : fête, festival
\item \eng{celebration} (n.) /« se-le-brei-chn/ : célébration
\item \eng{parade} (n.) /« pa-reid/ : défilé
\item \eng{march-past} (n.) /« march-pa-st/ : défilé (militaire/scolaire)
\item \eng{ceremony} (n.) /« se-re-mo-ni/ : cérémonie
\item \eng{prize-giving} (n.) : remise des prix
\item \eng{uniform} (n.) /« u-ni-form/ : uniforme, tenue
\item \eng{quarter} (n.) /« kor-teur/ : quartier (usage camerounais)
\item \eng{community field} (n.) : terrain communautaire
\item \eng{sharp} (adv.) /« charp/ : pile, précises (après l'heure)
\end{itemize}
\end{definition}

\courssub{Comparer les registres : la même question, trois formulations}

\begin{popfigure}
\begin{tikzpicture}
\node[draw=popblue, fill=popblueL, rounded corners, align=center, minimum width=4.2cm, minimum height=2.6cm] (mate) at (0,0) {\textbf{Mate (informel)}\\[2pt] \eng{When does it start?}\\ \eng{Are you going?}};
\node[draw=popgreen, fill=popgreenL, rounded corners, align=center, minimum width=4.2cm, minimum height=2.6cm] (neigh) at (5.2,0) {\textbf{Neighbour (neutre)}\\[2pt] \eng{What time does it start?}\\ \eng{That sounds wonderful!}};
\node[draw=poporange, fill=poporangeL, rounded corners, align=center, minimum width=4.2cm, minimum height=2.6cm] (princ) at (10.4,0) {\textbf{Principal (formel)}\\[2pt] \eng{Could you tell me what time}\\ \eng{the march-past begins?}};
\node[draw=poppurple, fill=poppurpleL, rounded corners, align=center, minimum width=5.5cm] (q) at (5.2,2.6) {\textbf{Même question : « Quand ça commence ? »}};
\draw[-{Stealth}, thick, poppurple] (q) -- (mate);
\draw[-{Stealth}, thick, poppurple] (q) -- (neigh);
\draw[-{Stealth}, thick, poppurple] (q) -- (princ);
\end{tikzpicture}
\legende{Une même information demandée à trois niveaux de registre}
\end{popfigure}

\begin{center}
\begin{tabularx}{\linewidth}{|X|X|X|}
\hline
\rowcolor{popblueL} \textbf{Mate (informel)} & \textbf{Neighbour (neutre)} & \textbf{Principal (formel)} \\
\hline
\eng{Hi Bih!} & \eng{Good evening, Mr Ekane.} & \eng{Good morning, Sir.} \\
\hline
\eng{Are you going to the parade?} & \eng{I heard there is a festival this weekend.} & \eng{Could you tell me what time the march-past begins?} \\
\hline
\eng{Why don't we go together?} & \eng{You are welcome to come.} & \eng{I'd like to know the programme, please.} \\
\hline
\eng{Great! See you there!} & \eng{See you on Saturday, then.} & \eng{Thank you very much, Sir. Goodbye.} \\
\hline
\end{tabularx}
\end{center}

\begin{aretenir}[L'essentiel]
\begin{itemize}
\item Le \cle{registre} dépend de l'interlocuteur : informel avec un camarade, neutre avec un voisin, formel avec une autorité.
\item Formel = salutation complète + \eng{Sir/Madam} + question indirecte (\eng{Could you tell me...?}) + remerciements.
\item Informel = salutation courte + question directe + contractions + exclamations.
\item Dans tous les cas, un dialogue complet suit le schéma : ouverture $\rightarrow$ demande $\rightarrow$ information $\rightarrow$ relance $\rightarrow$ clôture.
\end{itemize}
\end{aretenir}

\courssub{Prononciation, accentuation et intonation}

\begin{propriete}[Points de prononciation à maîtriser]
\begin{itemize}
\item \textbf{Accent de mot (word stress)} : \eng{fes-TI-val}, \eng{cel-e-BRA-tion}, \eng{com-MU-ni-ty}, \eng{au-THO-ri-ty}, \eng{pa-RADE}, \eng{uni-FORM}, \eng{pro-GRAMME}. L'accent ne tombe jamais sur le préfixe \eng{re-}, \eng{un-}, \eng{com-} (sauf noms composés).
\item \textbf{Accent de phrase (sentence stress)} : Mots de contenu (noms, verbes, adjectifs, adverbes) accentués ; mots grammaticaux (articles, auxiliaires, prépositions, pronoms) faibles.
\item \textbf{Intonation des questions} :
  \begin{itemize}
  \item Questions fermées (Yes/No) : intonation \textbf{montante} à la fin. Ex. \eng{Are you going?} (voix qui monte sur \eng{going?}).
  \item Questions ouvertes (Wh-) : intonation \textbf{descendante} à la fin. Ex. \eng{What time does it start?} (voix qui descend sur \eng{start?}).
  \end{itemize}
\item \textbf{Enchaînements (linking)} : \eng{What time} $\rightarrow$ « wot-taïm » ; \eng{Don't forget} $\rightarrow$ « doun-for-guet » ; \eng{See you} $\rightarrow$ « si-you » ; \eng{at 8} $\rightarrow$ « a-tei ».
\item \textbf{Contractions obligatoires à l'oral} : \eng{don't} /« dount/ (pas « do not »), \eng{it's} /« its/ (pas « it is »), \eng{let's} /« lets/ (pas « let us »), \eng{I'd} /« aid/.
\item \textbf{\eng{Sharp}} : se prononce /« charp/ (ch comme dans \emph{chat}), collé à l'heure : \eng{at 8 o'clock sharp} /« at eit o-klok charp/.
\end{itemize}
\end{propriete}

\begin{methode}[Travailler la prononciation d'un dialogue en 3 temps]
\begin{enumerate}
\item \textbf{Écoute mimétique} : écoutez le modèle (professeur ou audio) sans lire le texte. Repérez la musique de la phrase (montées/descentes).
\item \textbf{Lecture chorale} : lisez le dialogue à voix haute en même temps que le modèle, en imitant exactement l'intonation et les liaisons.
\item \textbf{Enregistrement personnel} : enregistrez-vous sur téléphone, réécoutez et comparez au modèle. Corrigez l'accent de mot et l'intonation des questions.
\end{enumerate}
\end{methode}

\courssub{Méthode et entraînement}

\begin{methode}[Exploiter un dialogue modèle en 4 temps]
\begin{enumerate}
\item \textbf{Écouter / lire} : lisez le dialogue silencieusement, puis écoutez-le (ou lisez-le à voix haute) pour la prononciation et l'intonation.
\item \textbf{Repérer les formules} : soulignez en couleur l'ouverture, la demande, la relance et la clôture ; identifiez le registre.
\item \textbf{Jouer tel quel} : en binôme, jouez le dialogue mot pour mot, avec contact visuel et gestuelle.
\item \textbf{Adapter} : remplacez les informations par les vôtres (une autre fête, une autre date, un autre lieu) et rejouez le dialogue.
\end{enumerate}
\end{methode}

\begin{exoresolu}[Repérer les formules dans le Dialogue 3]
Énoncé : dans le Dialogue 3 (voisins), relevez : (a) la formule d'ouverture, (b) une question, (c) une relance, (d) la formule de clôture.
\tcblower
Solution détaillée : (a) Ouverture : \eng{Good evening, Mr Ekane.} — salutation neutre avec le titre \eng{Mr}. (b) Question : \eng{What is the programme?} ou \eng{What time does it start?} — questions directes mais polies, typiques du registre neutre. (c) Relance : \eng{Really? That sounds wonderful!} — réaction qui montre l'intérêt et encourage le partenaire à continuer. (d) Clôture : \eng{See you on Saturday, then.} — prise de congé avec rendez-vous. Remarquez que le dialogue suit bien le schéma ouverture $\rightarrow$ demande $\rightarrow$ description $\rightarrow$ relance $\rightarrow$ invitation $\rightarrow$ clôture.
\end{exoresolu}

\begin{attention}[Ne récitez pas : dialoguez !]
L'erreur la plus fréquente des élèves francophones est de \textbf{réciter} le dialogue les yeux fixés sur la feuille, sans regarder le partenaire ni réagir à ce qu'il dit. Un vrai dialogue oral exige :
\begin{itemize}
\item le \cle{contact visuel} : regardez votre partenaire ;
\item des \cle{réactions} naturelles : \eng{Really?}, \eng{I see!}, \eng{Oh, interesting!}, \eng{That sounds great!} ;
\item une \cle{intonation} vivante : la voix monte dans les questions fermées (\eng{Are you going?}) et descend dans les questions en \eng{wh-} (\eng{What time does it start?}).
\end{itemize}
Un dialogue récité sans vie est une mauvaise note ; un dialogue imparfait mais vivant est une bonne note.
\end{attention}

\begin{experience}[Entraînement en binôme : substitution]
\begin{enumerate}
\item En binôme, lisez le Dialogue 1 à voix haute, chacun un rôle, puis inversez les rôles.
\item Remplacez maintenant les informations : choisissez une autre fête (\eng{the National Day}, \eng{the school cultural week}, \eng{the village festival}), une autre heure et un autre lieu (\eng{the school playground}, \eng{the town hall}, \eng{the church hall}).
\item Rejouez le dialogue avec ces nouvelles informations, sans regarder le texte.
\item Recommencez avec le Dialogue 2 (un élève demande le programme de la fête au proviseur) puis le Dialogue 3 (deux voisins parlent de la fête du quartier de votre choix : Bastos, Nkoulouloun, Molyko...).
\end{enumerate}
\end{experience}

\begin{savaistu}[Le dialogue et l'épreuve orale]
Aux examens camerounais (Probatoire et Baccalauréat), l'épreuve orale d'anglais valorise l'\cle{aisance} et l'\cle{interaction réelle} bien plus que la récitation par cœur. L'examinateur écoute : votre prononciation, votre intonation, votre capacité à réagir (\eng{Really? I see!}) et à poser des questions spontanées. S'entraîner régulièrement avec des dialogues modèles — puis les adapter — est la meilleure préparation : c'est exactement ce que vous faites dans cette leçon.
\end{savaistu}

\coursec{Prononciation, accentuation et intonation}

Dans les dialogues modèles que nous venons d'étudier, les mots et les formules fonctionnelles sont maintenant bien connus. Mais une question correcte sur le papier peut devenir incompréhensible à l'oral si l'accent tombe sur la mauvaise syllabe ou si la voix monte au lieu de descendre. Cette dernière étape est donc consacrée à la \cle{prononciation}, à l'\cle{accentuation} (\eng{word stress}) et à l'\cle{intonation}, afin que vos échanges sur les fêtes communautaires sonnent naturels et clairs.

\courssub{Observer avant de prononcer : l'accentuation des mots-clés}

Commençons par une observation. Écoutez (ou imaginez) ces deux versions de la même phrase :

\eng{We are going to the FES-ti-val.} (correct)

\eng{We are going to the fes-ti-VAL.} (incorrect, accent « à la française »)

En français, l'accent tombe presque toujours sur la \textbf{dernière syllabe} du groupe de mots : « nous allons à la fes-ti-\textbf{VAL} ». En anglais, chaque mot possède son propre accent, fixé une fois pour toutes, et il tombe \textbf{souvent vers le début} du mot. C'est la différence de rythme la plus importante entre les deux langues.

\begin{propriete}[Règle d'accentuation des mots]
En anglais, chaque mot de deux syllabes ou plus possède une \cle{syllabe accentuée} (\eng{stressed syllable}), prononcée plus fort, plus longtemps et sur un ton légèrement plus haut. Les autres syllabes sont faibles et rapides. Contrairement au français, l'accent ne se déplace pas : il fait partie du mot, comme son orthographe. Il faut donc \textbf{apprendre chaque mot avec son accent}.
\end{propriete}

Voici les mots-clés de notre chapitre, avec leur accentuation (la syllabe forte est en MAJUSCULES) et une prononciation guidée en lettres françaises (approximation pour francophones) :

\begin{center}
\begin{tabularx}{\linewidth}{|l|l|X|X|}
\hline
\rowcolor{popblueL} Mot & Accentuation & Prononciation guidée & Traduction \\
\hline
festival & \cle{FES}-ti-val & « fèss-ti-veul » & la fête, le festival \\
\hline
celebration & ce-le-\cle{BRA}-tion & « sè-lè-BREI-cheune » & la célébration \\
\hline
tradition & tra-\cle{DI}-tion & « treu-DI-cheune » & la tradition \\
\hline
parade & pa-\cle{RADE} & « peu-REID » & le défilé \\
\hline
community & com-\cle{MU}-ni-ty & « keu-MIOU-ni-ti » & la communauté \\
\hline
invitation & in-vi-\cle{TA}-tion & « ine-vi-TEI-cheune » & l'invitation \\
\hline
\end{tabularx}
\end{center}

\begin{exemplebox}[Exemples en contexte, avec traduction]
\eng{The community celebration is wonderful.} — « La célébration de la communauté est merveilleuse. »

Prononciation guidée : « ze keu-\textbf{MIOU}-ni-ti sè-lè-\textbf{BREI}-cheune iz \textbf{OUEN}-deur-feul ».

\eng{Did you receive an invitation to the traditional parade?} — « As-tu reçu une invitation au défilé traditionnel ? »

Prononciation guidée : « did you ri-\textbf{SIVD} èn ine-vi-\textbf{TEI}-cheune tou ze treu-\textbf{DI}-cheu-neul peu-\textbf{REID} ? »
\end{exemplebox}

\begin{attention}[L'erreur n° 1 du francophone : accentuer la fin du mot]
Un élève francophone dit spontanément *\eng{festi-VAL}, *\eng{celebra-TION}, *\eng{communi-TY}, par habitude du français. Résultat : l'interlocuteur anglophone peine à reconnaître le mot. Corrige-toi systématiquement : \textbf{FES}-ti-val, ce-le-\textbf{BRA}-tion, com-\textbf{MU}-ni-ty, pa-\textbf{RADE} (accent sur la deuxième syllabe en anglais britannique standard). Astuce : frappe doucement la table sur la syllabe forte quand tu répètes.
\end{attention}

\begin{savaistu}[Quand l'accent change le sens]
En anglais, mal placer l'accent peut changer le sens du mot ou le rendre incompréhensible. Ainsi \eng{PREsent} (accent sur la première syllabe) est un nom qui signifie « un cadeau » ou « le présent », tandis que \eng{preSENT} (accent sur la deuxième syllabe) est un verbe qui signifie « présenter ». De même : \eng{REcord} (un disque, un record) contre \eng{reCORD} (enregistrer). L'accentuation n'est pas un détail : c'est une partie du mot !
\end{savaistu}

\courssub{L'intonation : la voix monte ou descend ?}

L'\cle{intonation} est la mélodie de la phrase : la voix monte-t-elle ou descend-elle à la fin ? En anglais, cette mélodie suit une règle simple et très fiable.

\begin{propriete}[Règle d'intonation]
\begin{itemize}
\item \textbf{Question fermée} (réponse \eng{yes} ou \eng{no}, commençant par un auxiliaire : \eng{are you...?}, \eng{do you...?}, \eng{is the...?}) : la voix \textbf{monte} à la fin ↗.
\item \textbf{Question en Wh-} (\eng{when, where, what, who, how}) et \textbf{affirmation} : la voix \textbf{descend} à la fin ↘.
\end{itemize}
Exception utile : les questions fermées très longues retombent parfois à la toute fin, mais au lycée, applique la règle de base sans hésiter.
\end{propriete}

\begin{textebox}[Phrases à répéter avec intonation]
Is the parade on Saturday? ↗

Are you coming to the festival with your neighbours? ↗

Where does the festival take place? ↘

When does the celebration start? ↘

It starts at nine o'clock in the community hall. ↘

The traditional dances are wonderful. ↘
\end{textebox}

Comparons avec le français : en français, la voix monte aussi pour les questions fermées (« Tu viens ? ↗ »), donc ce point est facile. En revanche, pour les questions en \eng{wh-}, le français monte souvent (« Où est-ce que tu \textbf{vas} ? ↗ »), alors que l'anglais descend fermement : \eng{WHERE are you GOing? ↘}. C'est cette chute finale qui donne à l'anglais son ton assuré.

\begin{popfigure}
\begin{tikzpicture}[scale=1, every node/.style={font=\small}]
% Question fermée : montante
\draw[->, thick, popmuted] (0,0) -- (6.5,0);
\node[popdark, anchor=north] at (0.5,-0.3) {Are};
\node[popdark, anchor=north] at (1.5,-0.3) {you};
\node[popblue, anchor=north, font=\small\bfseries] at (2.7,-0.3) {COM-ing};
\node[popdark, anchor=north] at (4.2,-0.3) {to the};
\node[popblue, anchor=north, font=\small\bfseries] at (5.5,-0.3) {FES-tival?};
\draw[very thick, popblue, -{Stealth[length=3mm]}] plot[smooth, tension=0.8] coordinates {(0,0.2) (1.5,0.3) (2.7,0.5) (4.2,0.7) (5.5,1.3)};
\node[popblue, font=\small\bfseries, align=center, anchor=south] at (3.2,1.7) {Question fermée : la voix monte ↗};

% Question Wh- : descendante
\draw[->, thick, popmuted] (8,0) -- (14.5,0);
\node[poporange, anchor=north, font=\small\bfseries] at (8.8,-0.3) {WHERE};
\node[popdark, anchor=north] at (10,-0.3) {is the};
\node[poporange, anchor=north, font=\small\bfseries] at (11.5,-0.3) {pa-RADE?};
\draw[very thick, poporange, -{Stealth[length=3mm]}] plot[smooth, tension=0.8] coordinates {(8,1.3) (9.5,1.0) (11,0.6) (12.5,0.2)};
\node[poporange, font=\small\bfseries, align=center, anchor=south] at (11.5,1.7) {Question en Wh- : la voix descend ↘};
\end{tikzpicture}
\legende{Deux courbes mélodiques : intonation montante sur une question fermée, descendante sur une question en Wh-.}
\end{popfigure}

\courssub{Les sons difficiles pour les francophones}

Trois sons reviennent constamment dans nos dialogues sur les fêtes et méritent un entraînement ciblé.

\textbf{1. Le son \eng{th} (voisé /ð/) de \eng{the, this, that, there}.} Ce son n'existe pas en français. Il se prononce « z » doux et soufflé, la langue entre les dents — jamais « z » sec ni « d ».

\begin{methode}[Bien prononcer le son \eng{th} /ð/]
\begin{enumerate}
\item Place le \textbf{bout de la langue entre les dents}, légèrement visible (ne la mord pas).
\item Souffle doucement en faisant vibrer les cordes vocales : le son ressemble à un « z » bourdonné et soufflé.
\item Vérifie dans un miroir : si ta langue n'est pas entre les dents, tu produis « z » ou « d », pas \eng{th}.
\item Entraîne-toi sur la phrase du chapitre : \eng{the traditional celebration this Thursday} (« ze treu-DI-cheu-neul sè-lè-BREI-cheune ziss THEURZ-deï »), d'abord lentement, puis à vitesse normale.
\end{enumerate}
\end{methode}

\textbf{2. Le \eng{h} aspiré de \eng{hall, happy, home}.} En français, le « h » est toujours muet. En anglais, il faut \textbf{souffler} un petit « h » chaud, comme quand on souffle sur des lunettes pour les nettoyer : \eng{the community hall} (« ze keu-MIOU-ni-ti \textbf{h}ôl »), \eng{happy} (« \textbf{h}æ-pi »).

\textbf{3. Le \eng{r} anglais de \eng{parade, traditional, neighbours}.} Le \eng{r} anglais ne se roule pas comme le « r » français : la pointe de la langue recule vers le palais mou \textbf{sans le toucher}. Prononce « peu-REID » avec un « r » doux, presque un « w » léger : \eng{traditional} (« treu-DI-cheu-neul »).

\begin{attention}[Faux amis de prononciation]
Ne prononce jamais \eng{the} « dé », \eng{this} « diss » sec, ni \eng{hall} « al » (sans souffle). Et attention à \eng{celebration} : le « c » initial se dit « s » (« sè-lè-BREI-cheune »), pas « k ».
\end{attention}

\courssub{Le rythme anglais : frapper fort, avaler le reste}

Dernière clé du naturel : le \cle{rythme}. Le français donne à peu près la même force à chaque syllabe (rythme syllabique) ; l'anglais, lui, \textbf{accentue les mots importants} (noms, verbes pleins, adjectifs, mots interrogatifs, adverbes) et \textbf{avale les petits mots grammaticaux} (\eng{the, a, of, to, and, are, do}), qui deviennent rapides, faibles et souvent réduits au \eng{schwa} /ə/.

\begin{exemplebox}[Observer le rythme (mots forts en gras, flèche d'intonation finale)]
\eng{The \textbf{FES}-ti-val \textbf{STARTS} at \textbf{NINE} o'\textbf{CLOCK}.} ↘ — « La fête commence à neuf heures. »

Mots forts (lexicaux) : \eng{festival, starts, nine, clock}. Mots faibles (grammaticaux), presque « mangés » : \eng{the} (« ze »), \eng{at} (« èt »).

\eng{\textbf{WHERE} does the \textbf{pa-RADE} take \textbf{PLACE}?} ↘ — « Où le défilé a-t-il lieu ? »

Forts : \eng{where, parade, place}. Faibles : \eng{does the} (« déz ze »), \eng{take} (verbe léger ici, souvent décentré).
\end{exemplebox}

\begin{center}
\begin{tabularx}{\linewidth}{|X|X|X|}
\hline
\rowcolor{popblueL} Aspect & Français & Anglais \\
\hline
Accent du mot & Sur la dernière syllabe du groupe & Fixe, souvent au début du mot \\
\hline
Question fermée & Voix montante ↗ & Voix montante ↗ \\
\hline
Question en « où, quand » & Souvent montante ↗ & Descendante ↘ \\
\hline
Rythme & Syllabes régulières (rythme syllabique) & Mots importants forts/lents, petits mots faibles/rapides (rythme accentuel) \\
\hline
Lettre « h » & Toujours muette & Soufflée : \eng{hall, happy} \\
\hline
Lettre « r » & Roulée (uvulaire) & Douce, langue reculée (alvéolaire) \\
\hline
\end{tabularx}
\end{center}

\courssub{Entraînement guidé}

\begin{exoresolu}[Placer l'accent et l'intonation]
Énoncé : dans les phrases suivantes, souligne la syllabe accentuée des mots principaux et indique par une flèche (↗ ou ↘) l'intonation finale.
\begin{enumerate}
\item \eng{Is the celebration in the community hall?}
\item \eng{When does the traditional parade begin?}
\item \eng{Our neighbours received an invitation.}
\end{enumerate}
\tcblower
Solution détaillée :
\begin{enumerate}
\item \eng{Is the ce-le-\textbf{BRA}-tion in the com-\textbf{MU}-ni-ty \textbf{HALL}?} ↗ — Question fermée (\eng{Is...?}), donc voix montante à la fin.
\item \eng{\textbf{WHEN} does the tra-\textbf{DI}-tional pa-\textbf{RADE} be-\textbf{GIN}?} ↘ — Question en \eng{wh-}, donc voix descendante.
\item \eng{Our \textbf{NEIGH}-bours re-\textbf{CEIVED} an in-vi-\textbf{TA}-tion.} ↘ — Affirmation, donc voix descendante. Attention : \eng{neighbours} s'accentue sur la première syllabe et se prononce « \textbf{NÉ}-beurz » ; \eng{received} sur la deuxième « re-\textbf{SIVD} ».
\end{enumerate}
\end{exoresolu}

\begin{exercice}[À ton tour]
\begin{enumerate}
\item Répète en chœur après le professeur les six mots-clés du tableau en frappant la syllabe forte.
\item Lis à voix haute les phrases de la boîte « Phrases à répéter avec intonation » en respectant les flèches d'intonation, puis enregistre-toi avec ton téléphone.
\item Auto-évaluation : réécoute ton enregistrement et coche :
      \begin{itemize}
      \item Mon \eng{th} est-il entre les dents (bourdonné) ?
      \item Mon \eng{h} de \eng{hall} est-il soufflé ?
      \item Ma voix monte-t-elle sur les questions fermées et descend-elle sur les questions en \eng{wh-} et les affirmations ?
      \item Les petits mots (\eng{the, a, to, of, does}) sont-ils rapides et faibles ?
      \end{itemize}
      Refais les phrases qui posent problème.
\item Avec un camarade, rejoue le dialogue modèle de la section précédente en exagérant volontairement les accents et l'intonation, comme des acteurs.
\end{enumerate}
\end{exercice}

\begin{corrige}[Exercice « À ton tour » : critères de réussite]
Il n'y a pas de réponse écrite unique : la réussite se juge à l'oral. Tu as réussi si :
\begin{itemize}
\item[(a)] Les syllabes fortes tombent au bon endroit (\textbf{FES}-ti-val, ce-le-\textbf{BRA}-tion, com-\textbf{MU}-ni-ty, in-vi-\textbf{TA}-tion, pa-\textbf{RADE}, tra-\textbf{DI}-tional).
\item[(b)] Les questions fermées montent ↗ et les questions en \eng{wh-} ainsi que les affirmations descendent ↘.
\item[(c)] Le \eng{th} est soufflé langue entre les dents (bourdonné), le \eng{h} de \eng{hall} est expiré, et le \eng{r} de \eng{parade} est doux (non roulé).
\item[(d)] Les mots grammaticaux (\eng{the, a, to, of, does, the}) sont rapides, faibles et contiennent souvent un schwa /ə/.
\end{itemize}
Compare ton enregistrement avec la lecture du professeur (ou un modèle audio fiable) et note tes progrès.
\end{corrige}

\begin{aretenir}[L'essentiel de la prononciation]
\begin{itemize}
\item Chaque mot anglais a un accent fixe, souvent au début : \textbf{FES}-ti-val, com-\textbf{MU}-ni-ty — jamais à la française sur la dernière syllabe.
\item Question fermée (oui/non) = voix qui \textbf{monte} ↗ ; question en \eng{wh-} et affirmation = voix qui \textbf{descend} ↘.
\item \eng{th} (/ð/) = langue entre les dents, son bourdonné ; \eng{h} = soufflé ; \eng{r} = doux, langue reculée sans toucher le palais.
\item Rythme anglais : mots lexicaux forts et lents, mots grammaticaux faibles et rapides (schwa).
\end{itemize}
\end{aretenir}

Tu es maintenant prêt à passer aux jeux de rôle guidés : avec un vocabulaire sûr, des formules fonctionnelles maîtrisées et une prononciation claire, ton dialogue sur les fêtes communautaires sera à la fois correct et naturel.

\coursec{Jeux de rôle guidés et entraînement}

Nous savons maintenant saluer selon l'interlocuteur, demander et donner des informations, relancer et conclure une conversation, et prononcer avec une intonation correcte. Il est temps de \textbf{tout réunir dans la pratique} : c'est le rôle des \eng{role plays} (jeux de rôle). Un jeu de rôle est une simulation : chaque élève reçoit un rôle et doit parler \emph{sans lire}, comme dans la vraie vie. C'est exactement l'exercice qui vous attendra à l'oral du Baccalauréat.

\begin{methode}[Réussir un jeu de rôle en 4 étapes]
\begin{enumerate}
\item \textbf{Read your card.} Lisez attentivement votre carte de rôle : qui êtes-vous ? que devez-vous demander ou dire ? Soulignez les informations à échanger.
\item \textbf{Prepare.} Préparez mentalement (ou sur un brouillon) \textbf{3 questions} et \textbf{3 formules} (salutation, relance, conclusion). Ne rédigez pas le dialogue entier !
\item \textbf{Perform without reading.} Jouez la scène \emph{sans lire} : regardez votre partenaire, parlez fort, lentement, en articulant.
\item \textbf{Listen to the feedback.} Écoutez les remarques de l'observateur et du professeur, puis rejouez si possible en corrigeant vos erreurs.
\end{enumerate}
\end{methode}

\begin{popfigure}
\begin{tikzpicture}[node distance=1.6cm and 2.2cm]
\node[draw=popblue, fill=popblueL, rounded corners, thick, minimum width=2.6cm, minimum height=1.1cm, align=center, font=\bfseries] (prep) {1. Prepare};
\node[draw=poporange, fill=poporangeL, rounded corners, thick, minimum width=2.6cm, minimum height=1.1cm, align=center, font=\bfseries, right=of prep] (perf) {2. Perform};
\node[draw=popgreen, fill=popgreenL, rounded corners, thick, minimum width=2.6cm, minimum height=1.1cm, align=center, font=\bfseries, below=of perf] (obs) {3. Observe};
\node[draw=poppurple, fill=poppurpleL, rounded corners, thick, minimum width=2.6cm, minimum height=1.1cm, align=center, font=\bfseries, left=of obs] (feed) {4. Feedback};
\draw[-{Stealth[length=3mm]}, very thick, popblue] (prep) -- (perf);
\draw[-{Stealth[length=3mm]}, very thick, poporange] (perf) -- (obs);
\draw[-{Stealth[length=3mm]}, very thick, popgreen] (obs) -- (feed);
\draw[-{Stealth[length=3mm]}, very thick, poppurple] (feed) -- (prep);
\end{tikzpicture}
\legende{Le cercle des 4 étapes du jeu de rôle : Prepare, Perform, Observe, Feedback.}
\end{popfigure}

\courssub{Jeu de rôle 1 : organiser la sortie au festival du quartier (entre camarades)}

\textbf{Mise en situation.} Deux camarades de classe veulent aller ensemble au festival du quartier. Mais attention : \textbf{chacun possède des informations différentes} (c'est ce qu'on appelle un \eng{information gap}, un « écart d'information »). Ils doivent donc se poser des questions pour compléter leurs informations et décider ensemble du programme de la journée.

\begin{textebox}[Role card A — Student 1]
You are a new neighbour in the district. You heard about the community festival but you know nothing about it. Ask your friend about: the date, the place, the activities, and whether you can join. Use polite questions and show interest.
\end{textebox}

\begin{textebox}[Role card B — Student 2]
You have lived in the neighbourhood for ten years. You know the community festival very well: it takes place on Saturday 15th June at the community hall, from 10 a.m. to 8 p.m. There is traditional dancing, a food fair and a football match. Describe the festival and invite your neighbour to join your family.
\end{textebox}

\textbf{Consigne.} L'élève A commence la conversation. L'élève B répond avec ses informations. À la fin, vous devez avoir décidé ensemble : l'heure de départ, le moyen de transport (taxi, bendskin, marche) et ce que chacun apporte.

\begin{exemplebox}[Formules utiles pour ce jeu de rôle]
\begin{itemize}
\item \eng{Have you heard about the community festival?} « As-tu entendu parler de la fête du quartier ? »
\item \eng{When exactly does it take place?} « Quand exactement a-t-elle lieu ? »
\item \eng{May I join the celebration?} « Puis-je me joindre à la célébration ? »
\item \eng{You're most welcome!} « Vous êtes le bienvenu ! »
\item \eng{Shall we meet at the bus stop at nine?} « On se retrouve à l'arrêt de bus à neuf heures ? »
\end{itemize}
\end{exemplebox}

\courssub{Jeu de rôle 2 : interviewer le proviseur sur le Youth Day (élève / autorité)}

\textbf{Mise en situation.} Vous êtes le délégué de votre classe. Le proviseur (\eng{the principal}) accepte de répondre à vos questions sur le programme de la fête de la jeunesse (\eng{Youth Day}, le 11 février au Cameroun). Vous êtes face à une \textbf{autorité scolaire} : le registre doit donc être \emph{formel}.

\begin{propriete}[Questions indirectes : obligatoires avec une autorité]
Avec une personne en position d'autorité, on évite la question directe trop brusque (\eng{When does the parade start?}). On utilise la \cle{question indirecte}, plus polie :
\begin{center}
\textbf{Could you tell me / May I ask you / I would like to know + sujet + verbe ?}
\end{center}
Attention à l'\textbf{ordre des mots} : après l'introducteur, on revient à l'ordre de la phrase affirmative (sujet + verbe), sans inversion et sans \eng{do/does/did}.
\begin{itemize}
\item Direct : \eng{When does the parade start?}
\item Indirect : \eng{Could you tell me when the parade starts?} « Pourriez-vous me dire quand le défilé commence ? »
\item Direct : \eng{Will the Minister attend the ceremony?}
\item Indirect : \eng{May I ask you whether the Minister will attend the ceremony?} « Puis-je vous demander si le Ministre assistera à la cérémonie ? »
\end{itemize}
\end{propriete}

\begin{attention}[L'erreur classique de la question indirecte]
Les francophones gardent souvent l'inversion dans la question indirecte : \emph{Could you tell me when does the parade start?} est \textbf{faux}. Dites : \eng{Could you tell me when the parade starts?} Règle simple : après \eng{Could you tell me}, construisez une phrase affirmative normale.
\end{attention}

\textbf{Consigne.} L'élève-journaliste pose au moins 4 questions indirectes (date, programme, invités, tenue recommandée). L'élève-proviseur répond poliment et complètement, puis conclut : \eng{Thank you for your interest. See you on February 11th!}

\courssub{Jeu de rôle 3 : le nouveau voisin s'informe sur les traditions (entre voisins)}

\textbf{Mise en situation.} Une famille vient d'emménager dans votre quartier. Le nouveau voisin frappe à votre porte : il veut comprendre les traditions locales — le \eng{Ngondo} (festival traditionnel du peuple Sawa, à Douala), la fête annuelle du village, ou un mariage traditionnel auquel il est invité. Le registre est \textbf{semi-formel} : poli mais chaleureux.

\begin{center}
\begin{tabularx}{\linewidth}{|X|X|X|}
\hline
\rowcolor{popblueL} \textbf{Français} & \textbf{English} & \textbf{Remarque} \\
\hline
Je viens d'emménager. & I have just moved in. & \eng{move in} = emménager \\
\hline
Pourriez-vous me parler des traditions du quartier ? & Could you tell me about the local traditions? & question indirecte polie \\
\hline
Qu'est-ce que le Ngondo exactement ? & What exactly is the Ngondo? & \eng{exactly} renforce la question \\
\hline
Les visiteurs sont-ils les bienvenus ? & Are visitors welcome? & adjectif \eng{welcome} \\
\hline
Que dois-je porter pour le mariage ? & What should I wear for the wedding? & \eng{should} = conseil \\
\hline
Merci pour ces précieuses informations. & Thank you for this valuable information. & \eng{information} est indénombrable \\
\hline
\end{tabularx}
\end{center}

\begin{exoresolu}[Transformer des questions directes en questions polies]
Énoncé : le nouveau voisin veut être poli. Transformez ces questions directes en questions indirectes avec \eng{Could you tell me...}.
\begin{enumerate}
\item \eng{When does the Ngondo festival take place?}
\item \eng{What time does the traditional wedding begin?}
\item \eng{Can visitors take photos?}
\end{enumerate}
\tcblower
Solution détaillée :
\begin{enumerate}
\item \eng{Could you tell me when the Ngondo festival takes place?} — on supprime l'auxiliaire \eng{does} et on remet le verbe à la forme affirmative (\eng{takes}, car le sujet est à la 3\textsuperscript{e} personne).
\item \eng{Could you tell me what time the traditional wedding begins?} — même transformation : \eng{does ... begin} devient \eng{begins}.
\item \eng{Could you tell me whether visitors can take photos?} — pour une question fermée (oui/non) sans mot interrogatif, on introduit \eng{whether} (= « si ») et on garde l'ordre sujet + verbe.
\end{enumerate}
Piège : ne jamais écrire \emph{Could you tell me when does the festival take place?}
\end{exoresolu}

\courssub{Observer et s'auto-évaluer}

\textbf{Le rôle d'observateur.} Quand deux élèves jouent une scène, un troisième élève est \eng{observer} : il ne parle pas, il écoute et \textbf{coche sur une fiche d'écoute} les formules réellement entendues. Cela entraîne la compréhension orale et responsabilise tout le groupe.

\begin{center}
\begin{tabularx}{\linewidth}{|X|c|}
\hline
\rowcolor{popgreenL} \textbf{Listening checklist — tick what you hear} & \textbf{Heard} \\
\hline
A greeting adapted to the person (\eng{Hello Sir} / \eng{Hi}) & $\square$ \\
\hline
At least three questions & $\square$ \\
\hline
An indirect question (\eng{Could you tell me...}) & $\square$ \\
\hline
A follow-up (\eng{Really?} / \eng{What else?} / \eng{I see.}) & $\square$ \\
\hline
An invitation or a suggestion & $\square$ \\
\hline
A polite conclusion (\eng{Thank you very much. Goodbye!}) & $\square$ \\
\hline
\end{tabularx}
\end{center}

\textbf{Grille d'auto-évaluation.} Après votre prestation, posez-vous honnêtement ces questions :

\begin{center}
\begin{tabularx}{\linewidth}{|X|c|c|}
\hline
\rowcolor{poporangeL} \textbf{Self-assessment} & \textbf{Yes} & \textbf{Not yet} \\
\hline
Did I greet my partner correctly? & $\square$ & $\square$ \\
\hline
Did I ask at least three questions? & $\square$ & $\square$ \\
\hline
Did I react and follow up (\eng{Oh really? That's interesting!})? & $\square$ & $\square$ \\
\hline
Did I speak loudly and clearly? & $\square$ & $\square$ \\
\hline
Did I conclude politely? & $\square$ & $\square$ \\
\hline
\end{tabularx}
\end{center}

\begin{attention}[Le piège de la timidité : parler trop bas et trop vite]
Beaucoup d'élèves francophones, par nervosité, parlent \emph{trop bas} et \emph{trop vite} : l'examinateur n'entend rien et note mal. À l'oral, la règle d'or est : \textbf{speak slowly, loudly and clearly} — parlez lentement, fort et en articulant. Une phrase simple bien prononcée vaut mieux qu'une phrase compliquée marmonnée. N'oubliez pas : l'intonation montante dans les questions fermées (\eng{May I join you?} $\nearrow$) montre que vous maîtrisez la langue.
\end{attention}

\begin{savaistu}[Ce que les examinateurs du Bac adorent]
À l'oral du Baccalauréat, les examinateurs apprécient particulièrement les candidats qui \textbf{relancent eux-mêmes la conversation} au lieu d'attendre passivement les questions. Un simple \eng{And what about you, Sir?}, \eng{What else should I know?} ou \eng{Oh really? That's very interesting!} transforme un interrogatoire en véritable échange — et fait monter la note. Montrez que vous êtes un \emph{communicateur actif} !
\end{savaistu}

\courssub{Entraînement libre et variante Bac+}

\begin{experience}[Improvisation sans support — variante Bac+]
Quand les trois jeux de rôle guidés sont réussis, passez au niveau supérieur : \textbf{improviser un dialogue sans aucun support écrit}, à partir d'une simple consigne tirée au sort, comme à l'oral du Bac. Exemples de consignes :
\begin{enumerate}
\item \eng{You meet your neighbour on the day before the village festival. Talk about it.}
\item \eng{You are a journalist. Interview the mayor about the Youth Day celebrations.}
\item \eng{Your friend has never seen a traditional wedding. Describe one and invite him or her.}
\end{enumerate}
Règles : 30 secondes de préparation mentale, puis 2 minutes de dialogue. L'observateur utilise la fiche d'écoute. Ce complément, signalé comme variante Bac+, prépare directement à l'épreuve orale du baccalauréat.
\end{experience}

\begin{exercice}[À vous de jouer]
\begin{enumerate}
\item En binôme, jouez le jeu de rôle 1 en inversant les cartes (A devient B).
\item Écrivez trois questions indirectes qu'un délégué pourrait poser au proviseur sur la fête de fin d'année.
\item En groupe de trois, jouez le jeu de rôle 3 avec un observateur ; l'observateur rend compte à la classe : \eng{They greeted each other, asked five questions and concluded politely.}
\end{enumerate}
\end{exercice}

\begin{corrige}[Exercice — Jeux de rôle]
\begin{enumerate}
\item Production orale libre : vérifier la salutation, au moins trois questions, une relance et une conclusion polie.
\item Exemples : \eng{Could you tell me when the end-of-year party will take place?} / \eng{May I ask you whether parents are invited?} / \eng{I would like to know what time the ceremony starts.}
\item Compte rendu attendu au prétérit : \eng{asked}, \eng{greeted}, \eng{concluded}. L'observateur doit citer au moins deux formules entendues.
\end{enumerate}
\end{corrige}

\begin{aretenir}[L'essentiel de la section]
\begin{itemize}
\item Un jeu de rôle suit le cercle \textbf{Prepare $\rightarrow$ Perform $\rightarrow$ Observe $\rightarrow$ Feedback}.
\item Adaptez le registre : familier entre camarades, formel (questions indirectes) avec une autorité, semi-formel entre voisins.
\item Question indirecte : \eng{Could you tell me + sujet + verbe}, sans inversion.
\item Parlez \textbf{lentement, fort, en articulant}, et relancez la conversation : c'est ce qui fait la différence à l'oral du Bac.
\end{itemize}
\end{aretenir}

\newpage

\coursec{Activité d'intégration}

\courssub{Objectifs de l'activité}

À la fin de cette activité d'intégration, tu seras capable de :

\begin{itemize}
\item mener un dialogue complet et naturel en anglais sur une célébration communautaire (festival, fête du village, Youth Day, carnaval) ;
\item adapter ton registre de langue à ton interlocuteur : camarade (registre familier), voisin ou adulte (registre courant), chef traditionnel, proviseur ou journaliste (registre soutenu et poli) ;
\item mobiliser toutes les formules fonctionnelles de la leçon : saluer, demander une information (questions directes et indirectes), donner une information et décrire un programme avec des connecteurs, relancer la conversation, conclure poliment ;
\item soigner ta prononciation, ton accentuation et ton intonation (intonation descendante pour les questions en \eng{wh-}, montante pour les questions fermées) ;
\item évaluer la production orale d'un camarade à l'aide d'une grille d'observation simple.
\end{itemize}

\courssub{Situation}

\begin{textebox}[Festival Interview Day — the class becomes a celebrating village]
Your school, Government High School Mfoundi, is organising a \emph{Festival Interview Day} before the end-of-year cultural week. The classroom is turned into a celebrating village. Every student picks two cards:

An \textbf{identity card}: a classmate, a neighbour, the traditional chief of the village, the principal of the school, or a CRTV journalist.

An \textbf{information card} about a celebration: the Ngondo festival in Douala, the Youth Day parade on 11 February, a village annual festival, or the Notting Hill Carnival in London.

In rotating pairs, you hold short conversations of about three minutes. You greet your partner with the right register, ask at least three questions (including one indirect question if your partner is an authority), describe the programme of the celebration with link words, keep the conversation alive with a follow-up question, and end it politely. An observer fills in a checklist. The two best dialogues are performed in front of the class.
\end{textebox}

\begin{exemplebox}[Exemples de cartes]
\textbf{Identity card 3 — Traditional chief.} \eng{You are the traditional chief of your village. You are respected by everybody. People must greet you politely and use formal language.}

\textbf{Information card B — Ngondo festival.} \eng{The Ngondo festival takes place in Douala every year, near the Wouri river. There are traditional dances, canoe races and a ceremony on the water. Thousands of people attend it.}
\end{exemplebox}

\courssub{Consignes}

\begin{enumerate}
\item \textbf{Prépare ton identité.} Lis ta carte d'identité et ta carte d'information. Souligne trois informations clés (date, lieu, activités) et note deux mots de vocabulaire utiles (par exemple \eng{canoe race}, \eng{parade}, \eng{traditional dance}).
\item \textbf{Prépare tes questions.} Écris trois questions : deux questions directes en \eng{wh-} (\eng{When does the festival take place? Where is it held?}) et, si ton interlocuteur est une autorité, une question indirecte (\eng{Could you tell me when the festival starts?}).
\item \textbf{Premier tour de table.} Par binômes tournants, mène un dialogue de trois minutes. Respecte le plan : salutation adaptée au registre, questions, description du programme avec au moins deux connecteurs (\eng{first, then, after that, finally}), une relance (\eng{Really? And what happens next?}), une conclusion polie (\eng{Thank you very much for your time.}).
\item \textbf{Observe et évalue.} Quand tu es observateur, remplis la grille : formules utilisées (oui/non), intonation correcte, registre respecté. Note un point fort et un point à améliorer.
\item \textbf{Change de rôle.} Après chaque tour, échange les cartes et recommence avec un nouveau partenaire et un nouveau registre.
\item \textbf{Finale de classe.} Les deux binômes ayant obtenu les meilleures grilles jouent leur dialogue devant la classe. La classe applaudit et justifie son choix en anglais (\eng{Their dialogue was polite and lively.}).
\end{enumerate}

\courssub{Ressources à mobiliser}

\begin{aretenir}[Boîte à outils fonctionnelle]
\begin{itemize}
\item \textbf{Saluer selon le registre} : \eng{Hi Paul, how are you?} (camarade) ; \eng{Good morning, Sir.} (autorité) ; \eng{Good afternoon, Madam. Nice to meet you.} (adulte, journaliste).
\item \textbf{Demander} : questions directes \eng{When / Where / What time / How long...?} ; question indirecte polie \eng{Could you tell me where the festival is held?} (attention : ordre sujet + verbe, pas d'inversion).
\item \textbf{Informer et décrire} : \eng{It takes place in December. It lasts three days. First, there is a parade; then, people dance; finally, the chief gives a speech.}
\item \textbf{Relancer} : \eng{Really? That's interesting! And what about the food? Could you tell me more about the dances?}
\item \textbf{Conclure} : \eng{Thank you very much for your time. It was a pleasure talking to you. Goodbye, Sir.}
\end{itemize}
\end{aretenir}

\begin{propriete}[Rappel de grammaire : la question indirecte]
Après une formule d'introduction comme \eng{Could you tell me...}, \eng{Do you know...}, \eng{I would like to know...}, la question perd son inversion et reprend l'ordre normal de la phrase affirmative : \cle{sujet + verbe}.

\begin{center}
\begin{tabularx}{\linewidth}{|X|X|}
\hline
\rowcolor{popblueL} Question directe (inversion) & Question indirecte (ordre normal) \\
\hline
\eng{When does the festival start?} & \eng{Could you tell me when the festival starts?} \\
\hline
\eng{Where is the ceremony held?} & \eng{Do you know where the ceremony is held?} \\
\hline
\eng{What time does the parade begin?} & \eng{I would like to know what time the parade begins.} \\
\hline
\end{tabularx}
\end{center}

Remarque : à la 3\up{e} personne du singulier, le \eng{-s} du présent réapparaît dans la question indirecte (\eng{the festival starts}), puisqu'il n'y a plus d'auxiliaire \eng{does}.
\end{propriete}

\begin{attention}[Pièges fréquents]
Ne dis jamais \eng{Could you tell me where is the festival?} : dans une question indirecte, on revient à l'ordre sujet + verbe : \eng{Could you tell me where the festival is?}. Évite aussi le faux ami \eng{actually} (qui signifie « en fait », pas « actuellement » — pour « actuellement », dis \eng{at the moment} ou \eng{currently}) et n'oublie pas l'intonation montante sur les questions fermées (\eng{Do you enjoy it?}) et descendante sur les questions en \eng{wh-} (\eng{When does it start?}).
\end{attention}

\courssub{Prononciation, accentuation et intonation}

\begin{propriete}[Intonation des questions]
\begin{itemize}
\item \textbf{Questions en \eng{wh-}} (qui attendent une information) : la voix \textbf{descend} à la fin. \eng{When does the festival start?} « ouène doze ze fè-sti-val start » — la voix tombe sur \eng{start}.
\item \textbf{Questions fermées} (réponse \eng{yes/no}) : la voix \textbf{monte} à la fin. \eng{Do you enjoy it?} « dou iou èn-djoï ite » — la voix monte sur \eng{it}.
\item \textbf{Accentuation} : dans une phrase, on accentue les mots porteurs de sens (noms, verbes, adjectifs) et on affaiblit les petits mots (articles, prépositions, auxiliaires). \eng{The FES-ti-val takes PLACE in De-CEM-ber.}
\end{itemize}
\end{propriete}

\begin{exemplebox}[Mots difficiles à prononcer]
\begin{itemize}
\item \eng{festival} « fè-sti-val » (accent sur la première syllabe) ;
\item \eng{neighbour} « né-beur » (le \eng{-gh-} ne se prononce pas) ;
\item \eng{authorities} « o-so-ri-tiz » ;
\item \eng{ceremony} « sé-ri-mo-ni » ;
\item \eng{canoe} « ke-nou » ;
\item \eng{parade} « pe-rède » (accent sur la deuxième syllabe).
\end{itemize}
\end{exemplebox}

\courssub{Proposition de production et corrigé commenté}

\begin{exoresolu}[Dialogue modèle : une élève interviewe le chef traditionnel]
Énoncé : rédige et joue un dialogue complet entre Aïcha, élève de Première, et le chef traditionnel de son village, au sujet de la fête annuelle du village.\tcblower
\textbf{Production modèle :}

\eng{— Good morning, Sir. Thank you for welcoming me.}

\eng{— Good morning, my daughter. You are welcome.}

\eng{— Sir, could you tell me when the village festival takes place?}

\eng{— Of course. It takes place every year in December, after the harvest.}

\eng{— And how long does it last?}

\eng{— It lasts three days. First, we have traditional dances in the chief's compound; then, there is a big meal with all the families; finally, the young people organise a football match.}

\eng{— That sounds wonderful! Could you tell me more about the dances?}

\eng{— Yes, the dancers wear beautiful traditional clothes and the drummers play all night.}

\eng{— Really? And can visitors attend the festival?}

\eng{— Certainly, everybody is welcome.}

\eng{— Thank you very much for your time, Sir. It was a pleasure talking to you.}

\eng{— Thank you, my daughter. Goodbye.}

\textbf{Commentaire des choix de langue :}
\begin{itemize}
\item \textbf{Registre soutenu} : \eng{Good morning, Sir}, \eng{Thank you for welcoming me} — adapté à une autorité traditionnelle.
\item \textbf{Questions indirectes} : \eng{Could you tell me when the festival takes place?} — ordre sujet + verbe respecté (\eng{the festival takes}, pas d'inversion), marque de politesse envers le chef.
\item \textbf{Connecteurs} : \eng{first, then, finally} structurent la description du programme.
\item \textbf{Relances} : \eng{That sounds wonderful!}, \eng{Really?} avec intonation montante — elles maintiennent la conversation vivante.
\item \textbf{Conclusion polie} : \eng{Thank you very much for your time. It was a pleasure talking to you.}
\end{itemize}
\end{exoresolu}

\begin{exercice}[À toi de jouer]
Avec un camarade, construis un dialogue de huit répliques minimum entre un journaliste de la CRTV (registre poli) et un élève, au sujet du défilé du 11 février (\eng{Youth Day parade}). Utilise au moins une question indirecte, deux connecteurs, une relance et une conclusion polie. Jouez-le devant la classe en soignant l'intonation.
\end{exercice}

\begin{corrige}[Exercice : dialogue journaliste / élève]
\textbf{Production possible :}

\eng{— Good afternoon. I am a journalist from CRTV. Could I ask you a few questions about the Youth Day parade?}

\eng{— Good afternoon, Sir. Yes, of course.}

\eng{— Could you tell me when the parade starts?}

\eng{— It starts at eight o'clock in the morning, at the ceremonial ground.}

\eng{— And what happens during the ceremony?}

\eng{— First, the students march past the authorities; then, the school bands play music; after that, there are speeches; finally, we sing the national anthem.}

\eng{— That sounds impressive! Do you enjoy taking part in it?}

\eng{— Yes, very much. It is a great honour for our school.}

\eng{— Thank you very much for your time.}

\eng{— You are welcome, Sir. Goodbye.}

Ce dialogue respecte le registre poli, contient une question indirecte correcte (\eng{when the parade starts}), quatre connecteurs, une relance et une conclusion polie.
\end{corrige}

\courssub{Bilan de l'activité}

Cette activité t'a permis de mobiliser en situation réelle toutes les compétences de la leçon : saluer selon le registre, poser des questions directes et indirectes, décrire une célébration avec des connecteurs, relancer et conclure un dialogue. Tu as également entraîné ton oreille et ta voix : intonation montante pour les questions fermées, descendante pour les questions en \eng{wh-}, accentuation des mots importants. La grille d'observation t'a appris à évaluer une production orale, compétence utile pour préparer l'épreuve d'expression orale. Pour aller plus loin, enregistre-toi chez toi en train de refaire le dialogue avec un autre interlocuteur (voisin ou journaliste de la CRTV) et compare ton registre et ton intonation avec le modèle de la classe.

\newpage

\coursec{Méthodes et exercices résolus}

\coursec{Lesson 1 — Speaking: Exchanges information about interacting with mates, neighbours, adults and school authorities on community celebrations/festivals}

\begin{savaistu}[Les grandes fêtes communautaires au Cameroun]
Le Cameroun, « Afrique en miniature », compte une multitude de célébrations culturelles. Trois sont incontournables pour l'examen :
\begin{itemize}
\item \textbf{La Fête de la Jeunesse (Youth Day)} — 11 février : défilés scolaires (\eng{march past}), discours, activités sportives et culturelles dans tout le pays.
\item \textbf{Le Ngondo (peuple Sawa, Douala)} — novembre/décembre : rites sur le Wouri, pirogues, danses traditionnelles (\eng{essewe}), culte des ancêtres (\eng{jengu}).
\item \textbf{Le Nguon (Bamoun, Foumban)} — tous les deux ans (novembre/décembre) : grande cérémonie du roi (\eng{Fon}), masques, danses, marché artisanal, unité du peuple Bamoun.
\end{itemize}
Vocabulaire clé : \eng{chief} (chef traditionnel), \eng{esplanade} (esplanade), \eng{traditional ruler} (chef supérieur), \eng{procession} (procession), \eng{heritage} (patrimoine).
\end{savaistu}

\begin{textebox}[Situation de communication : Dialogue au quartier Mvog-Ada, Yaoundé]
\eng{— Good evening, Mr Atanga. Did you hear about the cultural festival next Saturday?}\par
\eng{— Good evening, my son. Yes, I did. The quarter head announced it at the last meeting.}\par
\eng{— Great! Do you know when it starts exactly?}\par
\eng{— It starts at 10 a.m. at the esplanade. There will be traditional dances and a big meal.}\par
\eng{— That sounds wonderful. Will the chief be there?}\par
\eng{— Of course. He always blesses the ceremony. You should come with your parents.}\par
\eng{— I will, Sir. Thank you for the information. Good night!}\par
\eng{— Good night, and see you there.}
\end{textebox}

\courssub{Échanger des informations dans la situation proposée}

\begin{methode}[Adapter son langage à l'interlocuteur]
Pour échanger des informations sur une fête communautaire, la première étape est de choisir le bon \cle{registre de langue} selon la personne à qui l'on parle.
\begin{enumerate}
\item \textbf{Identifier l'interlocuteur} : un camarade (\eng{mate} / \eng{classmate}), un voisin (\eng{neighbour}), un adulte ou une autorité scolaire (\eng{the principal}, \eng{a teacher}, \eng{the quarter head}).
\item \textbf{Choisir la formule d'ouverture adaptée} : informelle avec un camarade (\eng{Hi! Guess what?}), polie avec un adulte (\eng{Good morning, Sir. May I ask you something?}).
\item \textbf{Annoncer le sujet} avec une phrase d'accroche : \eng{Have you heard about the Ngondo festival?} ou \eng{Our quarter is organising a cultural festival next month.}
\item \textbf{Enchaîner avec une question ouverte} pour lancer l'échange : \eng{What do you know about it?}
\item \textbf{Réagir aux réponses} avec des marqueurs d'écoute : \eng{Really? That's interesting! I see.}
\end{enumerate}
\begin{center}
\begin{tabularx}{\linewidth}{|X|X|}\hline
\rowcolor{popblueL} Interlocuteur & Formules adaptées \\ \hline
Camarade (\eng{mate}) & \eng{Hi Paul! / Guess what? / You know what?} \\ \hline
Voisin (\eng{neighbour}) & \eng{Good evening, Mr Atanga. / Did you hear the news?} \\ \hline
Adulte, autorité (\eng{principal, teacher}) & \eng{Good morning, Madam. / Excuse me, Sir, could I...?} \\ \hline
\end{tabularx}
\end{center}
\end{methode}

\begin{attention}[Quarter vs Neighbourhood]
Au Cameroun anglophone, \eng{quarter} (quartier) est très utilisé (\eng{Our quarter is organising...}). En anglais standard (Royaume-Uni/USA), on préfère \eng{neighbourhood} ou \eng{district}. Les deux sont acceptés au Bac, mais \eng{neighbourhood} est plus formel.
\end{attention}


\begin{exoresolu}[Choisir la bonne formule selon l'interlocuteur]
\textbf{Énoncé.} For each situation, choose the most appropriate opening sentence (A, B or C) and justify your choice.\par
1. You want to tell your classmate about the Nguon festival in Foumban.\par
A. \eng{Good morning, Sir. I would like to inform you about the Nguon festival.}\par
B. \eng{Hi! Have you heard about the Nguon festival in Foumban?}\par
C. \eng{Ladies and gentlemen, welcome to the Nguon festival.}\par
2. You want to ask the principal about the school's participation in Youth Day.\par
A. \eng{Hey! What's up with Youth Day?}\par
B. \eng{Good morning, Sir. May I ask you about our school's participation in Youth Day?}\par
C. \eng{Youth Day is coming, right?}
\tcblower
\textbf{Solution détaillée.}\par
\textbf{Question 1 : réponse B.} L'interlocuteur est un camarade de classe : le registre doit être informel et amical. La formule A conviendrait à une autorité (vouvoiement implicite, ton cérémonieux \eng{I would like to inform you}), et la formule C est un discours public, pas une conversation. La formule B utilise \eng{Hi!} et une question directe \eng{Have you heard about...?}, parfaite pour engager la conversation entre amis. \textbf{Remarque grammaticale :} le \emph{present perfect} \eng{Have you heard} s'emploie ici pour une nouvelle récente dont le moment n'est pas précisé.\par
\textbf{Question 2 : réponse B.} Le principal est une autorité scolaire : il faut un registre formel et poli. La formule A est trop familière (\eng{Hey! What's up...?} est réservé aux amis proches) ; la formule C manque de respect et de structure interrogative correcte. La formule B salue correctement (\eng{Good morning, Sir}) et utilise la demande de permission polie \eng{May I ask you about...?}, construite avec l'inversion du modal \eng{may}.\par
\textbf{Conclusion.} En anglais comme en français, on n'aborde pas un proviseur comme un copain : le choix de la formule d'ouverture est le premier critère d'évaluation d'un dialogue oral.
\end{exoresolu}

\courssub{Poser des questions et répondre de façon adaptée}

\begin{methode}[Demander des informations sur une fête]
Pour obtenir des informations précises sur une célébration, il faut maîtriser les \cle{mots interrogatifs} (\eng{wh-words}) et l'inversion sujet-auxiliaire.
\begin{enumerate}
\item \textbf{Choisir le bon mot interrogatif} selon l'information cherchée :
\begin{itemize}
\item \eng{When} (quand) : \eng{When does the festival take place?} — Prononciation : « ouène ».
\item \eng{Where} (où) : \eng{Where is it celebrated?} — Prononciation : « ouère ».
\item \eng{Who} (qui) : \eng{Who organises the ceremony?} — Prononciation : « hou ».
\item \eng{Why} (pourquoi) : \eng{Why do people wear traditional clothes?} — Prononciation : « ouaï ».
\item \eng{How long / How often} (combien de temps / à quelle fréquence) : \eng{How long does it last?}
\end{itemize}
\item \textbf{Respecter l'ordre des mots} : mot interrogatif + auxiliaire + sujet + verbe. \eng{Where does the festival take place?} et jamais \emph{Where the festival takes place?}
\item \textbf{Utiliser l'auxiliaire do/does/did} avec le \emph{present simple} et le \emph{preterit} ; pas d'auxiliaire quand \eng{who/what} est \textbf{sujet} : \eng{Who comes to the festival?}
\item \textbf{Répondre de façon complète} : reprendre l'auxiliaire pour les réponses courtes (\eng{Yes, it does. / No, they don't.}) et donner une phrase complète pour les questions ouvertes.
\item \textbf{Adoucir la question} avec un adulte (style indirect) : \eng{Could you tell me when the festival starts?} (pas d'inversion après \eng{tell me} : ordre sujet-verbe).
\end{enumerate}
\end{methode}

\begin{propriete}[Sujet vs Complément : le piège du \eng{Who} / \eng{What}]
\begin{center}
\begin{tabular}{|l|l|l|}\hline
\rowcolor{popblueL} Type de question & Structure & Exemple \\ \hline
\eng{Who/What} = SUJET & \eng{Who/What} + verbe conjugué (sans auxiliaire) & \eng{Who organises the festival?} (\eng{The committee organises it.}) \\ \hline
\eng{Who/What} = COMPLÉMENT & \eng{Who/What} + auxiliaire + sujet + verbe & \eng{Who did you see?} (\eng{I saw the chief.}) \\ \hline
\end{tabular}
\end{center}
\end{propriete}


\begin{exoresolu}[Construire les questions d'un dialogue]
\textbf{Énoncé.} Complete the dialogue between two classmates with the correct questions. Use the words in brackets.\par
A : \eng{Hi Nadia! Our quarter is organising a cultural festival.}\par
B : \eng{Great!} 1. ................................ (when / it / take place)?\par
A : \eng{It takes place on the first Saturday of December.}\par
B : 2. ................................ (who / organise / it)?\par
A : \eng{The neighbourhood committee organises it, with the help of the chief.}\par
B : 3. ................................ (what / people / do / during the festival)?\par
A : \eng{They dance, sing traditional songs and share a big meal.}\par
B : 4. ................................ (you / go / with your family)?\par
A : \eng{Yes, I am. We never miss it!}
\tcblower
\textbf{Solution détaillée.}\par
\textbf{1.} \eng{When does it take place?} — Mot interrogatif \eng{when} + auxiliaire \eng{does} (sujet \eng{it}, 3\textsuperscript{e} personne du singulier au \emph{present simple}) + sujet + base verbale \eng{take place}. Attention : sans l'auxiliaire, la phrase est fautive ; et le verbe reprend la base verbale (pas de \emph{s} après \eng{does}).\par
\textbf{2.} \eng{Who organises it?} — Ici \eng{who} est le \textbf{sujet} de la question : il n'y a ni auxiliaire ni inversion, et le verbe prend la marque de la 3\textsuperscript{e} personne : \eng{organises}. C'est le piège classique : on n'écrit pas \emph{Who does organise it?}\par
\textbf{3.} \eng{What do people do during the festival?} — \eng{What} est ici complément (on fait \emph{quoi} ?), donc inversion avec l'auxiliaire \eng{do} (sujet pluriel \eng{people}) + base verbale \eng{do}. La phrase contient donc deux fois \eng{do} : le premier est auxiliaire, le second est le verbe « faire ».\par
\textbf{4.} \eng{Are you going with your family?} — La réponse \eng{Yes, I am} montre que la question utilise le \emph{present continuous} (projet futur) : inversion de \eng{be} : \eng{Are you going...?}\par
\textbf{Conclusion.} Un bon dialogue repose sur l'alternance question ouverte (\eng{wh-}) / question fermée (auxiliaire en tête) ; vérifier toujours si le mot interrogatif est sujet ou complément.
\end{exoresolu}

\courssub{Donner des informations et décrire une célébration}

\begin{methode}[Décrire une fête communautaire en quelques phrases]
Pour informer un interlocuteur, il faut produire des phrases claires, au bon temps, avec un lexique précis.
\begin{enumerate}
\item \textbf{Présenter la fête} au \emph{present simple} (faits, traditions) : \eng{The Ngondo festival celebrates the culture of the Sawa people.}
\item \textbf{Situer dans le temps et l'espace} avec les bonnes prépositions :
\begin{itemize}
\item \eng{in} + mois / saison / année : \eng{in December}, \eng{in 2025}.
\item \eng{on} + jour précis / date : \eng{on 11th February}, \eng{on Monday}.
\item \eng{at} + lieu précis / heure : \eng{at the esplanade}, \eng{at 10 a.m.}.
\item \eng{in} + ville / pays : \eng{in Douala}, \eng{in Cameroon}.
\end{itemize}
\item \textbf{Décrire le déroulement} avec des connecteurs chronologiques : \eng{First, the chief blesses the ceremony. Then, dancers perform in the streets. After that, families share traditional dishes. Finally, the celebration ends with music.}
\item \textbf{Exprimer ses sentiments} : \eng{I really enjoy it. / It is fascinating. / What I like most is the traditional dance.}
\item \textbf{Inviter l'interlocuteur} : \eng{You should come and see it! / Why don't you join us this year?}
\end{enumerate}
\begin{center}
\begin{tabularx}{\linewidth}{|X|X|X|}\hline
\rowcolor{popblueL} Français & English & Remarque / Prononciation \\ \hline
une fête / un festival & a celebration / a festival & \eng{party} = fête privée (anniversaire) ; « fès-ti-val » \\ \hline
avoir lieu & to take place & \textbf{Jamais} \emph{to have place} \\ \hline
se déguiser, porter un costume & to wear traditional clothes & \eng{wear} (porté), pas \emph{carry} (porter qqch) ; « ouère » \\ \hline
un défilé & a parade / a march past & \eng{parade} : « pa-rèd » ; \eng{march past} (défilé officiel) \\ \hline
partager un repas & to share a meal & préposition \eng{with} : \eng{share a meal with neighbours} \\ \hline
bénir (la cérémonie) & to bless the ceremony & \eng{bless} : « blèss » (irrégulier : blessed/blest) \\ \hline
assister à (être présent) & to attend / to go to & \textbf{Faux ami :} \eng{to assist} = aider \\ \hline
\end{tabularx}
\end{center}
\end{methode}

\begin{exoresolu}[Décrire la Fête de la Jeunesse]
\textbf{Énoncé.} A new student from Nigeria asks you about Youth Day in Cameroon. Write five to six sentences to inform him, using the following cues: date — place — activities — your feelings. Use \eng{first, then, after that, finally}.
\tcblower
\textbf{Solution détaillée.}\par
\textbf{Réflexion préalable.} On décrit des faits habituels et annuels : le \textbf{present simple} s'impose. Les prépositions de date : \eng{on} + jour précis (\eng{on 11th February}). Le vocabulaire : \eng{parade, march past, songs, speeches, sports activities}.\par
\textbf{Réponse rédigée.}\par
\eng{Youth Day is celebrated in Cameroon on 11th February every year. On that day, students and pupils gather at their schools very early in the morning. First, there is a big march past in front of the authorities. Then, the head of the local administration gives a speech about the role of young people. After that, schools organise songs, dances and football matches. Finally, families celebrate together at home. I really enjoy this day because it shows that the youth is the future of our country.}\par
Traduction : « La Fête de la Jeunesse est célébrée au Cameroun le 11 février chaque année. Ce jour-là, les élèves et étudiants se rassemblent très tôt dans leurs écoles. D'abord, il y a un grand défilé devant les autorités. Ensuite, le chef de l'administration locale prononce un discours sur le rôle des jeunes. Après cela, les écoles organisent des chants, des danses et des matchs de football. Enfin, les familles fêtent cela ensemble à la maison. J'aime vraiment ce jour car il montre que la jeunesse est l'avenir de notre pays. »\par
\textbf{Justifications grammaticales.} \eng{is celebrated} : présent simple passif (la fête \emph{subit} la célébration). \eng{gather, gives, organise, celebrate} : present simple pour des habitudes. \eng{enjoy} est suivi d'un nom (\eng{this day}) ; avec un verbe, il exigerait le gérondif : \eng{I enjoy watching the parade.}
\end{exoresolu}

\courssub{Mener un court dialogue en anglais : relancer et conclure}

\begin{methode}[Relancer et conclure une conversation]
Un dialogue naturel n'est pas une suite de questions-réponses mécaniques : il faut savoir \cle{relancer} et \cle{conclure}.
\begin{enumerate}
\item \textbf{Montrer qu'on écoute} (marqueurs de réaction) : \eng{Really? / I see. / That's wonderful! / How interesting! / Sounds great!}
\item \textbf{Relancer} pour obtenir plus de détails : \eng{And what about the food? / What else happens? / Can you tell me more about the dances?}
\item \textbf{Demander une précision} si on n'a pas compris : \eng{Sorry, what do you mean by “ngondo”? / Could you repeat that, please? / Pardon?}
\item \textbf{Exprimer son accord ou son opinion} : \eng{I agree with you. / In my opinion, festivals bring people together. / I think you are right.}
\item \textbf{Conclure poliment} selon le registre :
\begin{itemize}
\item Informel (camarade) : \eng{See you at the festival! / Catch you later! / Thanks, bye!}
\item Formel (adulte/autorité) : \eng{Thank you for the information. / It was nice talking to you. / I must go now, goodbye Sir.}
\end{itemize}
\end{enumerate}
\end{methode}

\begin{attention}[Intonation des questions et accentuation de mots]
À l'oral, l'intonation porte le sens :
\begin{itemize}
\item Questions fermées (réponse \eng{yes/no}) $\rightarrow$ intonation \textbf{montante} ($\nearrow$) : \eng{Do you like the festival?}
\item Questions en \eng{wh-} (ouvertes) $\rightarrow$ intonation \textbf{descendante} ($\searrow$) : \eng{Where does it take place?}
\end{itemize}
\textbf{Accentuation des mots clés (Word Stress)} : en anglais, un mot a une syllabe tonique forte.
\begin{center}
\begin{tabular}{|l|l|l|}\hline
\rowcolor{popblueL} Mot & Syllabe accentuée & Prononciation approximative \\ \hline
festival & 1\textsuperscript{re} & « FÈS-ti-val » \\ \hline
celebration & 3\textsuperscript{e} & « cé-lé-BRÉ-chon » \\ \hline
parade & 2\textsuperscript{e} & « pa-RADE » \\ \hline
ceremony & 1\textsuperscript{re} & « CÉ-ré-mo-ni » \\ \hline
authority & 2\textsuperscript{e} & « au-THO-ri-ty » \\ \hline
\end{tabular}
\end{center}
\end{attention}

\begin{exoresolu}[Compléter et prolonger un dialogue modèle]
\textbf{Énoncé.} Complete the dialogue with appropriate expressions from the lesson, then write the final sentence to close the conversation politely.\par
A (neighbour) : \eng{Good evening, Mrs Biya. Did you enjoy the traditional festival in our quarter?}\par
B : \eng{Oh yes, it was wonderful! The dances were magnificent.}\par
A : 1. ................................ (reaction) \eng{And what about the food?}\par
B : \eng{We shared “ndolé” and “koki” with all the families.}\par
A : 2. ................................ (ask for more details about the music)\par
B : \eng{There were drums and traditional songs all night long.}\par
A : 3. ................................ (close the conversation politely)
\tcblower
\textbf{Solution détaillée.}\par
\textbf{1.} \eng{That's wonderful!} ou \eng{How interesting!} — Marqueur de réaction enthousiaste qui montre l'écoute active avant de relancer. On évite le silence ou la répétition mécanique de la question.\par
\textbf{2.} \eng{Can you tell me more about the music?} ou \eng{What about the music?} — Relance au style indirect : après \eng{Can you tell me}, pas d'inversion (\eng{the music was}, ordre sujet-verbe). La forme \emph{Can you tell me more about what was the music?} serait doublement fautive (inversion interdite + auxiliaire manquant).\par
\textbf{3.} \eng{Thank you very much for this lovely evening, Mrs Biya. Good night!} ou \eng{It was nice talking to you. See you at the next festival!} — Formule de remerciement + salutation de départ adaptée à une voisine adulte (registre poli : \eng{Mrs Biya}, pas de \eng{Bye-bye!} trop familier).\par
\textbf{Conclusion.} Un dialogue réussi contient toujours : une ouverture adaptée, des réactions, au moins une relance, et une clôture polie. C'est cette structure complète que l'examinateur évalue à l'oral.
\end{exoresolu}

\begin{attention}[Les 5 erreurs qui coûtent des points au Bac]
\begin{enumerate}
\item \textbf{Oublier l'auxiliaire dans les questions.} \emph{Where the festival takes place?} est fautif : il faut \eng{Where does the festival take place?} Sans auxiliaire \eng{do/does/did}, la question au \emph{present simple} ou au \emph{preterit} est incorrecte.
\item \textbf{Calquer « avoir lieu ».} Ne jamais traduire par \emph{to have place} : on dit \eng{to take place}. Exemple : \eng{The ceremony takes place in December.}
\item \textbf{Faux ami « assist ».} \eng{To assist} signifie « aider » ; « assister à une fête » se dit \eng{to attend a festival} ou \eng{to go to a festival}.
\item \textbf{Mauvaises prépositions de temps et de lieu.} On dit \eng{in December}, \eng{on 11th February}, \eng{at the stadium}, \eng{in Bamenda}. Le calque \emph{in 11th February} est une faute classique.
\item \textbf{Oublier le \emph{s} de la 3\textsuperscript{e} personne ou le doubler.} \eng{The festival lasts three days.} (avec \emph{s}) mais \eng{Does the festival last long?} (sans \emph{s} après \eng{does}). Jamais \emph{Does the festival lasts...?}
\end{enumerate}
\end{attention}

\begin{experience}[Jeu de rôle guidé : « La veille de la fête »]
\textbf{Consigne.} Par groupes de trois. L'un joue un \textbf{élève}, l'autre un \textbf{voisin adulte}, le troisième le \textbf{proviseur} (ou le chef de quartier).\par
\textbf{Scénario :} L'élève doit obtenir trois informations pour un article du journal scolaire : (1) date/heure/lieu exacts, (2) programme principal, (3) invitation officielle pour l'autorité.\par
\textbf{Contraintes :}
\begin{itemize}
\item Utiliser 3 registres différents (familier, standard, formel).
\item Inclure : 1 ouverture, 2 questions \eng{wh-}, 1 question fermée, 1 relance, 1 marqueur de réaction, 1 conclusion polie.
\item Durée : 2 minutes par dialogue.
\end{itemize}
\textbf{Grille d'auto-évaluation (cocher) :}
\begin{tabular}{|l|c|}\hline
\rowcolor{popblueL} Critère & Oui / Non \\ \hline
Ouverture adaptée à l'interlocuteur & \\ \hline
Questions correctes (auxiliaire / ordre des mots) & \\ \hline
Vocabulaire précis (\eng{take place, wear, attend, parade}) & \\ \hline
Connecteurs (\eng{first, then, finally}) & \\ \hline
Relance + Réaction + Conclusion & \\ \hline
Prononciation : intonation questions / accentuation mots & \\ \hline
\end{tabular}
\end{experience}

\begin{exercice}[Entraînement écrit : Préparer son oral]
\begin{enumerate}
\item \textbf{Vocabulaire.} Traduisez en anglais : 1. Le chef bénit la cérémonie. 2. Nous portons des tenues traditionnelles. 3. Le défilé a lieu sur l'esplanade. 4. J'assiste au festival chaque année. 5. Les familles partagent le repas avec les voisins.
\item \textbf{Grammaire (Questions).} Mettez les mots dans l'ordre pour former des questions correctes :\\
1. the festival / when / take place / does / ? \\
2. who / the ceremony / organise / ? \\
3. you / are / going / tomorrow / to the parade / ? \\
4. what / people / eat / do / during the festival / ? \\
5. could / me / tell / you / where / it / is / ?
\item \textbf{Expression écrite.} Vous écrivez un SMS à un ami (informel) puis un email au proviseur (formel) pour l'inviter au festival culturel de l'école samedi prochain (10h, cour de l'école, danses, chants, discours). 50 mots chacun.
\end{enumerate}
\end{exercice}

\begin{corrige}[Exercice n°1 -- Entraînement écrit]
\textbf{1. Vocabulaire.}
1. \eng{The chief blesses the ceremony.} (ou \eng{The traditional ruler blesses...})\\
2. \eng{We wear traditional clothes.} (ou \eng{traditional outfits / attire})\\
3. \eng{The parade takes place at the esplanade.} (ou \eng{The march past takes place...})\\
4. \eng{I attend the festival every year.} (ou \eng{I go to the festival...} — \textbf{pas} \emph{I assist...})\\
5. \eng{Families share the meal with neighbours.} (ou \eng{share a meal with their neighbours})\\

\textbf{2. Grammaire (Questions).}
1. \eng{When does the festival take place?} (auxiliaire \eng{does} + base verbale)\\
2. \eng{Who organises the ceremony?} (\eng{Who} = sujet $\rightarrow$ pas d'auxiliaire, \emph{s} à \eng{organises})\\
3. \eng{Are you going to the parade tomorrow?} (Present continuous : inversion \eng{be})\\
4. \eng{What do people eat during the festival?} (\eng{What} = complément $\rightarrow$ auxiliaire \eng{do})\\
5. \eng{Could you tell me where it is?} (Style indirect : \eng{where it is}, PAS \emph{where is it})\\

\textbf{3. Expression écrite (modèles).}
\textbf{SMS (ami) :} \eng{Hi Paul! Guess what? Our school cultural festival is this Saturday at 10 a.m. in the courtyard. There'll be dances, songs and speeches. You really should come! It's gonna be great. See you there?}\par
\textbf{Email (proviseur) :} \eng{Subject: Invitation to School Cultural Festival}\par
\eng{Good morning, Sir,}\par
\eng{May I invite you to our school cultural festival taking place this Saturday at 10 a.m. in the school courtyard? The programme includes traditional dances, songs and a speech by the principal. Your presence would honour us.}\par
\eng{Thank you for your consideration.}\par
\eng{Yours faithfully,}\par
\eng{[Nom de l'élève]}
\end{corrige}

\begin{popfigure}
\begin{center}\tcbox[colback=popink,colframe=popink,coltext=white,arc=9pt,boxsep=4pt,left=6pt,right=6pt]{\bfseries\small Interacting on Festivals}\end{center}
\vspace{-2pt}
\begin{tcbraster}[raster columns=2,raster equal height=rows,raster column skip=6pt,raster row skip=6pt,enhanced,blankest]
\begin{tcolorbox}[enhanced,colback=white,colframe=popgreen,boxrule=0.9pt,arc=3pt,title={Opening (Register)},fonttitle=\bfseries\scriptsize,coltitle=white,colbacktitle=popgreen,left=4pt,right=4pt,top=2pt,bottom=2pt,boxsep=1pt,toptitle=1.5pt,bottomtitle=1.5pt,halign=left]
\begin{itemize}[nosep,leftmargin=*,itemsep=1pt,font=\scriptsize]
\item Mate: \eng{Hi! Guess what?}
\item Adult: \eng{Good morning, Sir. May I...?}
\end{itemize}
\end{tcolorbox}
\begin{tcolorbox}[enhanced,colback=white,colframe=poporange,boxrule=0.9pt,arc=3pt,title={Asking (Wh- + Aux)},fonttitle=\bfseries\scriptsize,coltitle=white,colbacktitle=poporange,left=4pt,right=4pt,top=2pt,bottom=2pt,boxsep=1pt,toptitle=1.5pt,bottomtitle=1.5pt,halign=left]
\begin{itemize}[nosep,leftmargin=*,itemsep=1pt,font=\scriptsize]
\item Subject Q: \eng{Who organises?}
\item Object Q: \eng{What do people do?}
\item Indirect: \eng{Could you tell me...?}
\end{itemize}
\end{tcolorbox}
\begin{tcolorbox}[enhanced,colback=white,colframe=poppurple,boxrule=0.9pt,arc=3pt,title={Describing (Present Simple)},fonttitle=\bfseries\scriptsize,coltitle=white,colbacktitle=poppurple,left=4pt,right=4pt,top=2pt,bottom=2pt,boxsep=1pt,toptitle=1.5pt,bottomtitle=1.5pt,halign=left]
\begin{itemize}[nosep,leftmargin=*,itemsep=1pt,font=\scriptsize]
\item Time/Place: \eng{on 11 Feb, at...}
\item Sequence: \eng{First, Then, Finally}
\item Vocab: \eng{take place, wear, attend}
\end{itemize}
\end{tcolorbox}
\begin{tcolorbox}[enhanced,colback=white,colframe=poppink,boxrule=0.9pt,arc=3pt,title={Closing (Politeness)},fonttitle=\bfseries\scriptsize,coltitle=white,colbacktitle=poppink,left=4pt,right=4pt,top=2pt,bottom=2pt,boxsep=1pt,toptitle=1.5pt,bottomtitle=1.5pt,halign=left]
\begin{itemize}[nosep,leftmargin=*,itemsep=1pt,font=\scriptsize]
\item React: \eng{Really? Wonderful!}
\item Follow-up: \eng{What about...?}
\item Leave: \eng{Thank you. Goodbye Sir.}
\end{itemize}
\end{tcolorbox}
\end{tcbraster}

\legende{Carte mentale : Structure d'un échange oral sur les fêtes communautaires (4 compétences clés).}
\end{popfigure}

\newpage

\coursec{Exercices}

\courssub{Je vérifie mes connaissances}

\begin{exercice}[QCM : la bonne formule]
Pour chaque situation, choisis la formule la plus adaptée (a, b ou c).
\begin{enumerate}
\item You meet your principal in the corridor. You say :
\begin{enumerate}
\item \eng{Hey! What's up?}
\item \eng{Good morning, Sir. May I ask you a question?}
\item \eng{Hi there! Long time no see!}
\end{enumerate}
\item To ask a neighbour about the date of the village festival, you say :
\begin{enumerate}
\item \eng{Could you tell me when the festival will take place?}
\item \eng{Festival when?}
\item \eng{You must tell me the date now.}
\end{enumerate}
\item To keep the conversation going, you can say :
\begin{enumerate}
\item \eng{That's all. Goodbye.}
\item \eng{Really? Tell me more about it!}
\item \eng{I don't care.}
\end{enumerate}
\item To end a conversation politely with an adult, you say :
\begin{enumerate}
\item \eng{Stop talking now.}
\item \eng{Thank you very much for your time. Have a nice day!}
\item \eng{Bye-bye, old man!}
\end{enumerate}
\end{enumerate}
\end{exercice}

\begin{exercice}[Vrai ou faux ? Justifie]
Dis si les affirmations suivantes sont vraies (V) ou fausses (F), puis justifie ta réponse en une phrase en français.
\begin{enumerate}
\item \eng{With a close friend, I can use informal language like “What's up?”.}
\item \eng{“Could you tell me...?” is less polite than “Tell me...”.}
\item \eng{In a formal conversation with a school authority, contractions like “I'm” and “don't” are recommended.}
\item \eng{To show interest and encourage the speaker, I can say “Oh, I see!” or “How interesting!”.}
\item \eng{In English, yes/no questions are usually pronounced with a rising intonation.}
\end{enumerate}
\end{exercice}

\begin{exercice}[Vocabulaire : les fêtes communautaires]
Complète chaque phrase avec le mot anglais qui convient, choisi dans la liste : \eng{celebration -- parade -- feast -- guest -- tradition -- rehearsal}.
\begin{enumerate}
\item The whole village attended the \ldots\ldots\ldots\ldots of the chief's birthday.
\item A colourful \ldots\ldots\ldots\ldots marched through the streets of Bamenda on Youth Day.
\item After the ceremony, the family prepared a huge \ldots\ldots\ldots\ldots with ndolé and roasted fish.
\item Every \ldots\ldots\ldots\ldots received a small gift before leaving the party.
\item Wearing traditional outfits on National Day is an old \ldots\ldots\ldots\ldots in our community.
\item The choir had a \ldots\ldots\ldots\ldots on Saturday to prepare the songs for the festival.
\end{enumerate}
\end{exercice}

\begin{exercice}[Conjugaison et formes à compléter]
Complète les questions avec la forme correcte de l'auxiliaire ou du verbe entre parenthèses.
\begin{enumerate}
\item When \ldots\ldots\ldots the festival \ldots\ldots\ldots (begin) this year?
\item \ldots\ldots\ldots you \ldots\ldots\ldots (ever / attend) the Ngondo festival in Douala?
\item What \ldots\ldots\ldots the students \ldots\ldots\ldots (prepare) for Youth Day at the moment?
\item How long \ldots\ldots\ldots the celebration \ldots\ldots\ldots (last) last year?
\item \ldots\ldots\ldots your neighbours \ldots\ldots\ldots (invite) you to the ceremony next Saturday?
\end{enumerate}
\end{exercice}

\courssub{Je m'entraîne}

\begin{exercice}[Traduction]
Traduis les phrases suivantes en anglais, en utilisant les formules fonctionnelles de la leçon.
\begin{enumerate}
\item Bonjour Monsieur le Proviseur, puis-je vous poser une question sur la fête de la jeunesse ?
\item Pourriez-vous me dire où la cérémonie aura lieu ?
\item Vraiment ? Racontez-m'en plus sur le défilé !
\item Merci beaucoup pour ces informations. Bonne journée !
\item Salut Paul ! Tu viens à la fête du village samedi ?
\end{enumerate}
\end{exercice}

\begin{exercice}[Questions directes et questions polies]
Transforme les cinq questions directes suivantes en questions polies indirectes commençant par \eng{Could you tell me...} ou \eng{Do you know...}. Attention à l'ordre des mots !
\begin{enumerate}
\item When does the festival start?
\item Where will the parade take place?
\item Who organised the celebration last year?
\item How much does a ticket for the concert cost?
\item Why did the mayor cancel the ceremony?
\end{enumerate}
\end{exercice}

\begin{exercice}[Dialogue à trous]
Complète le dialogue suivant avec les formules de la liste : \eng{Have a nice day -- May I ask you a question -- Could you tell me -- That sounds wonderful -- Thank you for your time}.

\begin{textebox}[At the principal's office]
A: Good morning, Sir. \ldots\ldots\ldots\ldots\ldots\ldots about the school celebration?\\
B: Good morning. Of course, go ahead.\\
A: \ldots\ldots\ldots\ldots\ldots\ldots when the celebration will take place?\\
B: It will be held on Friday 7th February, in the school courtyard.\\
A: \ldots\ldots\ldots\ldots\ldots\ldots ! Will parents be invited?\\
B: Yes, they will. We expect many guests.\\
A: \ldots\ldots\ldots\ldots\ldots\ldots , Sir. \ldots\ldots\ldots\ldots\ldots\ldots !\\
B: You're welcome. Goodbye!
\end{textebox}
\end{exercice}

\begin{exercice}[Appariement : le bon registre]
Pour chaque réplique ci-dessous, indique si elle convient à (a) un camarade de classe, (b) un voisin adulte, (c) une autorité scolaire. Si la réplique ne convient à aucun des trois, reformule-la correctement.
\begin{enumerate}
\item \eng{Good afternoon, Madam. I would like to invite you to our community festival.}
\item \eng{Yo man! You coming to the party tonight?}
\item \eng{Give me the programme of the ceremony, now.}
\item \eng{Excuse me, Mrs Ngono, do you know what time the feast begins?}
\item \eng{Sir, I beg you to shut up and listen to me.}
\item \eng{Hi Sandra! Did you enjoy the parade yesterday?}
\end{enumerate}
\end{exercice}

\courssub{Je me prépare au Bac}

\begin{exercice}[Compréhension de l'écrit et questions de langue]
Lis le texte suivant, puis réponds aux questions en anglais, par des phrases complètes.

\begin{textebox}[A festival that brings people together]
Every year in December, the people of Buea organise a community festival to celebrate the end of the year. For one week, the streets around the market are decorated with colourful flags, and every evening, dance groups and choirs perform on a large stage built by the town council. Traders sell grilled fish, soya and roasted plantains, and children run everywhere with balloons.

Mrs Efange, a sixty-year-old neighbour, has attended the festival since she was a child. “When I was young, only a few families took part,” she remembers. “Today, thousands of people come from all over the region. It is our best moment of the year.” The festival also has a serious side: on the last day, traditional rulers and local authorities give speeches about peace, solidarity and respect between generations. The mayor always closes the ceremony with the same words: “A community that celebrates together stays together.”

This year, the students of Government High School Buea have been invited to present a play about the history of their town. They are rehearsing every afternoon after classes, and their English teacher, Mr Tabi, is very proud of them. “The festival is a wonderful opportunity for my students to speak English in front of a real audience,” he explains.
\end{textebox}

\textbf{A. Comprehension questions}
\begin{enumerate}
\item When does the festival of Buea take place every year?
\item Mention two activities people can do during the festival.
\item How has the festival changed since Mrs Efange's childhood? Quote one sentence from the text to justify your answer.
\item What do traditional rulers talk about on the last day?
\item Why is Mr Tabi proud of his students?
\end{enumerate}

\textbf{B. Language questions}
\begin{enumerate}
\item Find in the text the opposite of: (a) \eng{the beginning}; (b) \eng{war}.
\item Turn into the passive voice: \eng{The town council builds a large stage.}
\item Turn into an indirect polite question: \eng{When does the mayor close the ceremony?} (Begin with \eng{Could you tell me...})
\end{enumerate}
\end{exercice}

\begin{exercice}[Traduction type Bac]
Traduis en anglais le passage suivant (5 phrases), en soignant les temps verbaux et le registre.

\begin{textebox}[Texte à traduire]
Samedi dernier, notre village a célébré la fête traditionnelle annuelle. Le chef a invité tous les habitants, et beaucoup de visiteurs sont venus de Yaoundé. Les femmes ont préparé un grand repas pendant que les jeunes dansaient au son des tambours. Mon voisin m'a dit que c'était la plus belle cérémonie de l'année. Si tu viens l'année prochaine, tu découvriras nos traditions.
\end{textebox}
\end{exercice}

\begin{exercice}[Expression écrite et jeu de rôle type Bac]
\textbf{Sujet de composition.} Your school is preparing the celebration of Youth Day. Write an article of about 150--180 words for the school magazine in which you describe the programme of the day, explain why this celebration is important for students, and invite parents and neighbours to attend. Give your article a title.

\textbf{Jeu de rôle oral (en binôme).} À partir de la consigne suivante, improvise avec un partenaire un dialogue d'environ deux minutes, puis jouez-le devant la classe :

\begin{textebox}[Role play]
Student A: You are a journalist for the school magazine. Interview the principal about the coming school celebration (date, place, activities, guests). Ask at least five questions, react to the answers, and conclude politely.\\
Student B: You are the principal. Answer the journalist's questions, give details about the celebration, and invite the whole community.
\end{textebox}

Critères d'évaluation : salutations adaptées au registre, questions correctement formées, relances (\eng{Really? Tell me more!}), conclusion polie, prononciation et intonation.
\end{exercice}

\newpage

\coursec{Corrigés des exercices}

\begin{corrige}[Exercice 1 — QCM : la bonne formule]
\begin{enumerate}
\item \textbf{Réponse b.} \eng{Good morning, Sir. May I ask you a question?} — Le proviseur est une autorité scolaire : il faut un registre formel (salutation complète + \eng{May I...?}). La réplique a est trop familière (\eng{Hey! What's up?}) et la c est réservée aux amis.
\item \textbf{Réponse a.} \eng{Could you tell me when the festival will take place?} — Question polie indirecte, adaptée à un voisin adulte. La réplique b est grammaticalement incorrecte (question sans auxiliaire ni inversion) et la c est impolie (\eng{You must...} exprime un ordre ou une obligation, pas une demande d'information).
\item \textbf{Réponse b.} \eng{Really? Tell me more about it!} — C'est une formule de relance qui montre de l'intérêt et encourage le locuteur à continuer. Les répliques a et c mettent fin à la conversation ou expriment l'indifférence.
\item \textbf{Réponse b.} \eng{Thank you very much for your time. Have a nice day!} — Conclusion polie adaptée à un adulte : remerciement puis souhait. Les répliques a et c sont grossières.
\end{enumerate}
\textbf{Point de méthode :} avant de choisir une formule, identifie toujours deux choses : l'interlocuteur (ami, adulte, autorité) qui détermine le \emph{registre}, et le moment de l'échange (ouverture, demande, relance, clôture) qui détermine la \emph{fonction} de la formule.
\end{corrige}

\begin{corrige}[Exercice 2 — Vrai ou faux ? Justifie]
\begin{enumerate}
\item \textbf{Vrai.} Avec un ami proche, le registre informel est naturel : \eng{What's up?}, \eng{Hi!}, les contractions (\eng{I'm}, \eng{don't}) et l'argot sont acceptés.
\item \textbf{Faux.} C'est le contraire : \eng{Could you tell me...?} est \emph{plus} poli que l'impératif \eng{Tell me...}, car le conditionnel \eng{could} adoucit la demande. L'impératif direct sonne comme un ordre.
\item \textbf{Faux.} Dans un échange formel avec une autorité, on évite les contractions : on préfère les formes pleines \eng{I am}, \eng{do not}, \eng{I would like}, qui marquent le respect et le sérieux.
\item \textbf{Vrai.} \eng{Oh, I see!}, \eng{How interesting!}, \eng{Really?} sont des marqueurs d'écoute active : ils montrent qu'on suit la conversation et encouragent le locuteur à poursuivre.
\item \textbf{Vrai.} Les questions fermées (réponse \eng{yes/no}) se prononcent généralement avec une intonation montante, tandis que les questions en \eng{wh-} (\eng{when}, \eng{where}, \eng{why}...) ont une intonation descendante.
\end{enumerate}
\end{corrige}

\begin{corrige}[Exercice 3 — Vocabulaire : les fêtes communautaires]
\begin{enumerate}
\item \eng{The whole village attended the \textbf{celebration} of the chief's birthday.} (célébration)
\item \eng{A colourful \textbf{parade} marched through the streets of Bamenda on Youth Day.} (défilé — indice : le verbe \eng{marched})
\item \eng{...the family prepared a huge \textbf{feast} with ndolé and roasted fish.} (festin — indice : les plats cités)
\item \eng{Every \textbf{guest} received a small gift before leaving the party.} (invité)
\item \eng{Wearing traditional outfits on National Day is an old \textbf{tradition} in our community.} (tradition)
\item \eng{The choir had a \textbf{rehearsal} on Saturday to prepare the songs for the festival.} (répétition — indice : \eng{to prepare})
\end{enumerate}
\textbf{Point de vigilance :} ne confonds pas \eng{guest} (l'invité, celui qui est reçu) et \eng{host} (l'hôte, celui qui reçoit) ; \eng{feast} désigne un repas festif abondant, pas la fête en général (qui se dit \eng{festival} ou \eng{celebration}).
\end{corrige}

\begin{corrige}[Exercice 4 — Conjugaison et formes à compléter]
\begin{enumerate}
\item \eng{When \textbf{does} the festival \textbf{begin} this year?} — Présent simple (programme officiel, horaire fixe) : auxiliaire \eng{does} + base verbale.
\item \eng{\textbf{Have} you \textbf{ever attended} the Ngondo festival in Douala?} — Present perfect avec \eng{ever} (expérience de vie, sans date précise) : \eng{have} + participe passé \eng{attended}.
\item \eng{What \textbf{are} the students \textbf{preparing} for Youth Day at the moment?} — Présent continu (repère \eng{at the moment}) : \eng{are} + \eng{preparing} (attention, \eng{prepare} perd son \eng{-e} final devant \eng{-ing}).
\item \eng{How long \textbf{did} the celebration \textbf{last} last year?} — Prétérit simple (repère \eng{last year}) : auxiliaire \eng{did} + base verbale \eng{last}.
\item \eng{\textbf{Will} your neighbours \textbf{invite} you to the ceremony next Saturday?} — Futur avec \eng{will} (repère \eng{next Saturday}) + base verbale. (Accepté aussi : \eng{Are your neighbours going to invite you...?})
\end{enumerate}
\textbf{Règle rappelée :} dans une question avec auxiliaire (\eng{do/does/did}, \eng{will}, \eng{have}), le verbe principal reste toujours à la base verbale (ou au participe passé pour les temps composés) : jamais \eng{does the festival begins}.
\end{corrige}

\begin{corrige}[Exercice 5 — Traduction]
\begin{enumerate}
\item \eng{Good morning, Sir. May I ask you a question about Youth Day?} — Registre formel : \eng{May I...?} (et non \eng{Can I...?}, plus familier) ; \eng{Youth Day} avec majuscules (fête officielle).
\item \eng{Could you tell me where the ceremony will take place?} — Question indirecte : ordre sujet + verbe (\eng{the ceremony will}), sans inversion.
\item \eng{Really? Tell me more about the parade!} — Relance enthousiaste ; structure \eng{tell me more about} + nom.
\item \eng{Thank you very much for the information. Have a nice day!} — Attention : \eng{information} est indénombrable en anglais (jamais \eng{informations}).
\item \eng{Hi Paul! Are you coming to the village festival on Saturday?} — Registre informel ; présent continu \eng{are you coming} pour un projet futur déjà décidé ; préposition \eng{on} devant un jour.
\end{enumerate}
\textbf{Points de vigilance :} \eng{avoir lieu} se traduit par \eng{take place} (jamais \eng{have place}) ; en conclusion, \eng{bonne journée} se dit \eng{Have a nice day!} et non \eng{Good day!}.
\end{corrige}

\begin{corrige}[Exercice 6 — Questions directes et questions polies]
\begin{enumerate}
\item \eng{Could you tell me when the festival starts?}
\item \eng{Could you tell me where the parade will take place?}
\item \eng{Do you know who organised the celebration last year?}
\item \eng{Do you know how much a ticket for the concert costs?}
\item \eng{Could you tell me why the mayor cancelled the ceremony?}
\end{enumerate}
\textbf{Justification grammaticale :} dans une question indirecte introduite par \eng{Could you tell me...} ou \eng{Do you know...}, la seconde partie suit l'ordre \emph{affirmatif} : mot interrogatif + sujet + verbe. On supprime donc l'auxiliaire \eng{do/does/did} et on accorde le verbe avec le sujet (\eng{the festival starts}, \eng{a ticket costs}). Le temps ne change pas : \eng{organised} et \eng{cancelled} restent au prétérit. Erreur fréquente des francophones : conserver l'inversion (\eng{Could you tell me when does the festival start?}) — à bannir absolument.
\end{corrige}

\begin{corrige}[Exercice 7 — Dialogue à trous]
\begin{enumerate}
\item \eng{Good morning, Sir. \textbf{May I ask you a question} about the school celebration?} — Formule d'ouverture formelle : salutation + demande de permission.
\item \eng{\textbf{Could you tell me} when the celebration will take place?} — Demande d'information polie (question indirecte, ordre affirmatif après \eng{when}).
\item \eng{\textbf{That sounds wonderful}! Will parents be invited?} — Réaction enthousiaste (relance) suivie d'une nouvelle question.
\item \eng{\textbf{Thank you for your time}, Sir. \textbf{Have a nice day}!} — Conclusion polie en deux temps : remerciement puis souhait.
\end{enumerate}
\textbf{Point de méthode :} un dialogue formel bien construit suit toujours quatre étapes : salutation + demande de permission, question polie, réaction/relance, remerciement + congé.
\end{corrige}

\begin{corrige}[Exercice 8 — Appariement : le bon registre]
\begin{enumerate}
\item \eng{Good afternoon, Madam. I would like to invite you to our community festival.} $\rightarrow$ \textbf{(c) autorité scolaire} (ou tout adulte en situation officielle) : salutation formelle + \eng{I would like}.
\item \eng{Yo man! You coming to the party tonight?} $\rightarrow$ \textbf{(a) camarade de classe} : registre très familier, question informelle sans auxiliaire (acceptable à l'oral entre amis, mais à éviter à l'écrit et avec un adulte).
\item \eng{Give me the programme of the ceremony, now.} $\rightarrow$ \textbf{Ne convient à personne} (impératif brutal). Reformulation : \eng{Could you give me the programme of the ceremony, please?} (voisin) ou \eng{Could I have the programme of the ceremony, please, Sir?} (autorité).
\item \eng{Excuse me, Mrs Ngono, do you know what time the feast begins?} $\rightarrow$ \textbf{(b) voisin adulte} : poli mais pas protocolaire, question indirecte correcte.
\item \eng{Sir, I beg you to shut up and listen to me.} $\rightarrow$ \textbf{Ne convient à personne} (\eng{shut up} est grossier, même précédé de \eng{Sir}). Reformulation : \eng{Excuse me, Sir. May I have your attention, please?}
\item \eng{Hi Sandra! Did you enjoy the parade yesterday?} $\rightarrow$ \textbf{(a) camarade de classe} : prénom, registre familier, question au prétérit correctement formée.
\end{enumerate}
\textbf{Point de vigilance :} la politesse en anglais repose sur les formes (\eng{Could you...?}, \eng{please}, \eng{Excuse me}), pas seulement sur l'intention : un impératif sec reste impoli même accompagné de \eng{Sir}.
\end{corrige}

\begin{corrige}[Exercice 9 — Compréhension de l'écrit et questions de langue]
\textbf{A. Comprehension questions} (réponses modèles, en phrases complètes)
\begin{enumerate}
\item \eng{The festival of Buea takes place every year in December, to celebrate the end of the year.}
\item \eng{During the festival, people can watch dance groups and choirs perform on stage, and they can buy grilled fish, soya and roasted plantains from traders.} (Deux activités suffisent.)
\item \eng{The festival has grown enormously: only a few families took part in the past, whereas thousands of people attend today.} Justification citée : \eng{“When I was young, only a few families took part.”} ou \eng{“Today, thousands of people come from all over the region.”}
\item \eng{On the last day, traditional rulers and local authorities give speeches about peace, solidarity and respect between generations.}
\item \eng{Mr Tabi is proud of his students because they have been invited to present a play about the history of their town, which is a wonderful opportunity for them to speak English in front of a real audience.}
\end{enumerate}
\textbf{B. Language questions}
\begin{enumerate}
\item (a) L'opposé de \eng{the beginning} : \eng{the end} (« the end of the year »). (b) L'opposé de \eng{war} : \eng{peace} (« speeches about peace »).
\item \eng{A large stage is built by the town council.} — Passif au présent : \eng{is} + participe passé \eng{built} (verbe irrégulier : \eng{build -- built -- built}) ; le complément d'agent est introduit par \eng{by}.
\item \eng{Could you tell me when the mayor closes the ceremony?} — Question indirecte : ordre affirmatif, suppression de \eng{does}, verbe accordé avec le sujet (\eng{closes}).
\end{enumerate}
\textbf{Barème indicatif (sur 20) :} A : 5 $\times$ 2 = 10 points (réponse correcte 1 pt, phrase complète en anglais correct 1 pt) ; B : 2 + 4 + 4 = 10 points.
\textbf{Points de vigilance :} répondre par une phrase complète (jamais un mot isolé à l'examen) ; pour la question 3, la citation doit être exacte et placée entre guillemets ; au passif, ne pas oublier l'auxiliaire \eng{is} (\eng{A large stage built...} est faux).
\end{corrige}

\begin{corrige}[Exercice 10 — Traduction type Bac]
\textbf{Proposition de traduction :}

\eng{Last Saturday, our village celebrated the annual traditional festival. The chief invited all the inhabitants, and many visitors came from Yaoundé. The women prepared a big feast while the young people were dancing to the sound of the drums. My neighbour told me (that) it was the most beautiful ceremony of the year. If you come next year, you will discover our traditions.}

\textbf{Justifications et points de vigilance :}
\begin{itemize}
\item \eng{Samedi dernier} = \eng{Last Saturday} (sans \eng{the} ; \eng{on} facultatif) ; le repère passé impose le prétérit : \eng{celebrated}, \eng{invited}, \eng{came} (verbe irrégulier : \eng{come -- came -- come}).
\item \eng{pendant que les jeunes dansaient} : action longue en arrière-plan $\rightarrow$ prétérit continu \eng{were dancing}, introduit par \eng{while}.
\item \eng{au son des tambours} = \eng{to the sound of the drums}.
\item Discours rapporté au passé : \eng{told me (that) it was} (recul du temps : \eng{is} $\rightarrow$ \eng{was}) ; \eng{tell} exige un complément de personne (\eng{told me}), contrairement à \eng{say} (\eng{said (that)...}).
\item Superlatif : \eng{the most beautiful} (adjectif long de plus de deux syllabes).
\item Phrase en \eng{if} (condition réelle) : \eng{If you come} (présent) + \eng{you will discover} (futur) — jamais \eng{if you will come}.
\item \eng{l'année prochaine} = \eng{next year} (sans \eng{the}).
\end{itemize}
\textbf{Barème indicatif (sur 10) :} 2 points par phrase (sens global 1 pt, correction grammaticale et lexicale 1 pt).
\end{corrige}

\begin{corrige}[Exercice 11 — Expression écrite et jeu de rôle type Bac]
\textbf{Proposition de rédaction modèle (environ 165 mots) :}

\begin{textebox}[Model article]
\textbf{Youth Day: A Celebration for All of Us!}\\[2pt]
On Tuesday 11th February, our school will celebrate Youth Day, and this year's programme promises to be unforgettable. The day will begin at 8 a.m. with a colourful parade through the streets of the neighbourhood, led by the school band. Back in the school courtyard, students will present traditional dances, songs and a short play about the history of our town. At midday, a huge feast will be offered to all the guests, and the afternoon will end with a football match between students and teachers.

Youth Day is important for us students because it reminds us that young people have a role to play in the life of our community. It is also a wonderful opportunity to show our talents and to practise English in front of a real audience.

Dear parents and neighbours, you are all warmly invited! Come and celebrate with us: a community that celebrates together stays together.
\end{textebox}

\textbf{Points de vigilance pour l'article :}
\begin{itemize}
\item Donner un titre (exigé par le sujet) et organiser le texte en trois paragraphes : programme / importance / invitation.
\item Programme futur : alterner \eng{will} + base verbale et présent simple pour les horaires officiels (\eng{The day begins at 8 a.m.}).
\item Heures : \eng{at 8 a.m.} ; dates : \eng{on Tuesday 11th February}.
\item Éviter les faux amis : \eng{assister à} (être présent) = \eng{attend} (jamais \eng{assist} seul, qui signifie « aider ») ; \eng{une opportunité} = \eng{an opportunity} (correct ici).
\end{itemize}

\textbf{Grille d'évaluation du jeu de rôle (barème indicatif sur 10) :}
\begin{itemize}
\item Salutations adaptées au registre formel (\eng{Good morning, Sir. May I ask you some questions?}) : 2 pts.
\item Au moins cinq questions correctement formées (inversion ou question indirecte, auxiliaires, prépositions) : 3 pts.
\item Relances et réactions (\eng{Really? Tell me more!}, \eng{That sounds wonderful!}) : 2 pts.
\item Conclusion polie (\eng{Thank you for your time, Sir. Have a nice day!}) : 1 pt.
\item Prononciation, accentuation et intonation (montante pour les questions \eng{yes/no}, descendante pour les questions en \eng{wh-}) : 2 pts.
\end{itemize}

\textbf{Exemple d'ouverture attendue :} \eng{Good morning, Sir. I am a journalist for the school magazine. May I ask you a few questions about the coming celebration?} — \eng{Good morning. Of course, go ahead.}
\end{corrige}

\newpage

\coursec{Fiche bilan}

\begin{popfigure}
\begin{center}\tcbox[colback=popblue,colframe=popblue,coltext=white,arc=9pt,boxsep=4pt,left=6pt,right=6pt]{\bfseries\small Speaking: community celebrations}\end{center}
\vspace{-2pt}
\begin{tcbraster}[raster columns=3,raster equal height=rows,raster column skip=6pt,raster row skip=6pt,enhanced,blankest]
\begin{tcolorbox}[enhanced,colback=white,colframe=poporange,boxrule=0.9pt,arc=3pt,title={Greetings by interlocutor},fonttitle=\bfseries\scriptsize,coltitle=white,colbacktitle=poporange,left=4pt,right=4pt,top=2pt,bottom=2pt,boxsep=1pt,toptitle=1.5pt,bottomtitle=1.5pt,halign=left]
\begin{itemize}[nosep,leftmargin=*,itemsep=1pt,font=\scriptsize]
\item mate: Hi!
\item adult: Good morning, sir
\end{itemize}
\end{tcolorbox}
\begin{tcolorbox}[enhanced,colback=white,colframe=popgreen,boxrule=0.9pt,arc=3pt,title={Asking for information},fonttitle=\bfseries\scriptsize,coltitle=white,colbacktitle=popgreen,left=4pt,right=4pt,top=2pt,bottom=2pt,boxsep=1pt,toptitle=1.5pt,bottomtitle=1.5pt,halign=left]
\begin{itemize}[nosep,leftmargin=*,itemsep=1pt,font=\scriptsize]
\item When / Where / Who?
\end{itemize}
\end{tcolorbox}
\begin{tcolorbox}[enhanced,colback=white,colframe=poppurple,boxrule=0.9pt,arc=3pt,title={Giving information},fonttitle=\bfseries\scriptsize,coltitle=white,colbacktitle=poppurple,left=4pt,right=4pt,top=2pt,bottom=2pt,boxsep=1pt,toptitle=1.5pt,bottomtitle=1.5pt,halign=left]
\begin{itemize}[nosep,leftmargin=*,itemsep=1pt,font=\scriptsize]
\item It takes place...
\end{itemize}
\end{tcolorbox}
\begin{tcolorbox}[enhanced,colback=white,colframe=poppink,boxrule=0.9pt,arc=3pt,title={Keeping the talk alive},fonttitle=\bfseries\scriptsize,coltitle=white,colbacktitle=poppink,left=4pt,right=4pt,top=2pt,bottom=2pt,boxsep=1pt,toptitle=1.5pt,bottomtitle=1.5pt,halign=left]
\begin{itemize}[nosep,leftmargin=*,itemsep=1pt,font=\scriptsize]
\item Really? What else?
\end{itemize}
\end{tcolorbox}
\begin{tcolorbox}[enhanced,colback=white,colframe=popgold,boxrule=0.9pt,arc=3pt,title={Concluding},fonttitle=\bfseries\scriptsize,coltitle=white,colbacktitle=popgold,left=4pt,right=4pt,top=2pt,bottom=2pt,boxsep=1pt,toptitle=1.5pt,bottomtitle=1.5pt,halign=left]
\begin{itemize}[nosep,leftmargin=*,itemsep=1pt,font=\scriptsize]
\item See you there!
\end{itemize}
\end{tcolorbox}
\begin{tcolorbox}[enhanced,colback=white,colframe=popblue!60,boxrule=0.9pt,arc=3pt,title={Intonation},fonttitle=\bfseries\scriptsize,coltitle=white,colbacktitle=popblue!60,left=4pt,right=4pt,top=2pt,bottom=2pt,boxsep=1pt,toptitle=1.5pt,bottomtitle=1.5pt,halign=left]
\begin{itemize}[nosep,leftmargin=*,itemsep=1pt,font=\scriptsize]
\item rising = question
\end{itemize}
\end{tcolorbox}
\begin{tcolorbox}[enhanced,colback=white,colframe=popgreen!60,boxrule=0.9pt,arc=3pt,title={Register},fonttitle=\bfseries\scriptsize,coltitle=white,colbacktitle=popgreen!60,left=4pt,right=4pt,top=2pt,bottom=2pt,boxsep=1pt,toptitle=1.5pt,bottomtitle=1.5pt,halign=left]
\begin{itemize}[nosep,leftmargin=*,itemsep=1pt,font=\scriptsize]
\item formal / informal
\end{itemize}
\end{tcolorbox}
\end{tcbraster}

\legende{Carte mentale de la leçon : échanger des informations sur les fêtes communautaires.}
\end{popfigure}

\begin{bilan}
\textbf{1. Lexique clé (à connaître par cœur)}
\begin{itemize}
  \item \eng{a celebration / a festival} : une fête / un festival ; \eng{a feast} : un festin.
  \item \eng{a traditional dance} : une danse traditionnelle ; \eng{a chief / a traditional ruler} : un chef traditionnel.
  \item \eng{to take place} : avoir lieu ; \eng{to gather} : se rassembler ; \eng{to attend} : assister à.
  \item \eng{neighbours} : les voisins ; \eng{school authorities} : les autorités scolaires ; \eng{mates} : camarades.
  \item \eng{food and drinks are shared} : on partage nourriture et boissons ; \eng{gifts are offered} : on offre des cadeaux.
\end{itemize}

\textbf{2. Formules fonctionnelles selon l'interlocuteur}
\begin{center}
\begin{tabularx}{\linewidth}{|X|X|X|}\hline
\rowcolor{popblueL} Fonction & Avec un camarade (informel) & Avec un adulte / une autorité (formel) \\ \hline
Saluer & \eng{Hi Peter! How are you?} & \eng{Good morning, sir. How do you do?} \\ \hline
Demander une information & \eng{When is the festival?} & \eng{Could you tell me when the festival takes place?} \\ \hline
Donner une information & \eng{It's on Saturday at the chief's palace.} & \eng{It will be held on Saturday at the community hall.} \\ \hline
Relancer & \eng{Really? What else?} & \eng{That is interesting. Could you tell me more?} \\ \hline
Conclure & \eng{Great! See you there!} & \eng{Thank you very much, sir. Goodbye.} \\ \hline
\end{tabularx}
\end{center}

\textbf{3. Grammaire à connaître par cœur}
\begin{itemize}
  \item \cle{Questions en Wh-} : mot interrogatif + auxiliaire + sujet + verbe : \eng{Where does the festival take place?} (jamais \emph{Where takes place...?}).
  \item \cle{Questions fermées} : auxiliaire + sujet + verbe : \eng{Do you attend the festival every year?} — intonation montante.
  \item \cle{Présent simple} pour les habitudes et les faits : \eng{People gather every December.} (3\up{e} personne : \eng{gathers}).
  \item \cle{Futur avec will / be going to} pour annoncer : \eng{The festival will start at 9 a.m.} / \eng{We are going to dance.}
  \item \cle{Politesse} : \eng{Could you...?}, \eng{Would you mind...?}, \eng{Please}, \eng{Thank you}.
\end{itemize}

\textbf{4. Prononciation et intonation}
\begin{itemize}
  \item Question fermée : voix qui \cle{monte} à la fin : \eng{Do you come?} (« dou you comme ? » avec la voix qui monte).
  \item Question en Wh- : voix qui \cle{descend} : \eng{When does it start?}
  \item Accentuation : \eng{FES-ti-val} (« fèss-ti-vol »), \eng{ce-le-BRA-tion} (« sè-lé-bré-chonne »), \eng{NEIGH-bours} (« nè-beurs »).
\end{itemize}

\textbf{5. Méthode express pour mener un dialogue}
\begin{enumerate}
  \item Saluer selon l'interlocuteur (formel ou informel).
  \item Poser une question claire (Wh- ou fermée).
  \item Écouter et réagir : \eng{Really?}, \eng{I see.}, \eng{That sounds great!}
  \item Donner à son tour une information.
  \item Conclure poliment : remercier et prendre congé.
\end{enumerate}
\end{bilan}

\begin{aretenir}[Checklist avant le Bac]
\begin{itemize}
  \item $\square$ Je sais saluer un camarade (\eng{Hi!}) et une autorité (\eng{Good morning, sir.}).
  \item $\square$ Je sais former une question en Wh- avec le bon ordre des mots.
  \item $\square$ Je sais former une question fermée avec \eng{do / does / did / will}.
  \item $\square$ Je connais le lexique des fêtes : \eng{festival, celebration, to take place, to attend, to gather}.
  \item $\square$ J'emploie le présent simple pour décrire une tradition (avec le \eng{-s} à la 3\up{e} personne).
  \item $\square$ J'annonce un événement futur avec \eng{will} ou \eng{be going to}.
  \item $\square$ Je sais relancer une conversation : \eng{Really? What else? And then?}
  \item $\square$ Je sais conclure poliment : \eng{Thank you. See you there!}
  \item $\square$ Je fais monter la voix pour une question fermée et descendre pour une question en Wh-.
  \item $\square$ J'adapte mon langage : informel avec un ami, formel avec un adulte ou une autorité.
  \item $\square$ Je peux décrire une fête camerounaise en cinq phrases en anglais.
  \item $\square$ Je peux jouer un dialogue complet de huit répliques sans notes.
\end{itemize}
\end{aretenir}

\findecours

\end{document}
